% Angles orientés et arcs capables — Mathématiques 1ère C — Cours signé OnBuch+
% Programme officiel MINESEC. Compiler avec : tectonic N15-arcs-capables.tex
\newcommand{\DOCMATIERE}{Mathématiques}
\newcommand{\DOCNIVEAU}{1ère C}
\newcommand{\DOCMODULE}{Configurations et transformations élémentaires du plan}
\newcommand{\DOCLECON}{N15}
\newcommand{\DOCTITRE}{Angles orientés et arcs capables}
% OnBuch+ — préambule « cours signés » ENRICHI (compilé avec tectonic / XeLaTeX)
% Système visuel « Pop » d'OnBuch+ : fond crème (jamais blanc), bordures encre
% épaisses, ombres « dures » décalées, titres Archivo Black en capitales.
% Version MATHÉMATIQUES : graphes (pgfplots/tikz), boîtes pédagogiques
% (définition, propriété, méthode, attention, le savais-tu, exercice résolu,
% exercice, TP, bilan), page de garde et signature OnBuch+ ancrée en bas.
%
% Macros attendues dans main.tex AVANT \input{preamble} :
%   \DOCMATIERE  \DOCNIVEAU  \DOCMODULE  \DOCLECON (numéro)  \DOCTITRE
\documentclass[11pt]{article}

\usepackage{fontspec}
\usepackage[french]{babel}
\usepackage[a4paper,margin=2cm,top=3.4cm,bottom=2.8cm,headheight=1.4cm,footskip=1.3cm]{geometry}
\usepackage{amsmath,amssymb,amsfonts}
\usepackage{mathtools} % \xmapsto, \coloneqq... souvent utilisés par les modèles
\usepackage{cancel} % simplifications barrées
% \abs{z} (module, valeur absolue) et \norm{u} (norme), utilisés par les modèles
\providecommand{\abs}{}\let\abs\relax\DeclarePairedDelimiter{\abs}{\lvert}{\rvert}
\providecommand{\norm}{}\let\norm\relax\DeclarePairedDelimiter{\norm}{\lVert}{\rVert}
\usepackage[table]{xcolor} % [table] : fournit aussi \rowcolors (alternance de lignes)
\usepackage{pagecolor}
\usepackage{tikz}
\usetikzlibrary{arrows.meta,positioning,calc,shapes.geometric,shapes.misc,shapes.arrows,decorations.pathmorphing,decorations.pathreplacing,decorations.markings,patterns,shadows,mindmap,trees,fit,backgrounds,matrix,intersections,angles,quotes}
\usepackage{pgfplots}
\usepackage{pgf-pie} % diagrammes circulaires \pie (statistiques)
\pgfplotsset{compat=1.18}
\usepgfplotslibrary{fillbetween,groupplots} % aires entre courbes (calcul intégral)
\usepackage{enumitem}
\usepackage{fancyhdr}
% [most] charge toutes les bibliothèques tcolorbox (listings, minted, raster,
% documentation...) et gonfle inutilement la mémoire TeX sur un document long
% (~300 boîtes) : on ne charge que ce qui est réellement utilisé.
\usepackage[skins,breakable]{tcolorbox}
\usepackage{graphicx}
\usepackage{colortbl}
\usepackage{booktabs}
\usepackage{tabularx}
\usepackage{diagbox} % case d'en-tête diagonale des tableaux à double entrée
% Colonnes extensibles centrée / gauche / droite (C, L, R), souvent utilisées par les modèles
\newcolumntype{C}{>{\centering\arraybackslash}X}
\newcolumntype{L}{>{\raggedright\arraybackslash}X}
\newcolumntype{R}{>{\raggedleft\arraybackslash}X}
\usepackage{array}
\usepackage{multicol}
\usepackage{siunitx}
% parse-numbers=false : accepte aussi \SI{2,5 \times 10^{-3}}{...}
\sisetup{locale=FR,output-decimal-marker={,},group-minimum-digits=4,parse-numbers=false}
\DeclareSIUnit\ohm{\ensuremath{\symup{\Omega}}} % U+2126 absent de la police
\usepackage{qrcode}
% \degree : commande fréquemment utilisée par les modèles pour un angle ou une
% température en dehors de \SI{}{} (non fournie par les paquets chargés).
\providecommand{\degree}{^{\circ}}
% \qed (fin de démonstration), utilisé par les modèles sans amsthm
\providecommand{\qed}{\hfill$\square$}
% \barre : double barre des valeurs interdites dans un tableau de variations (tkz-tab)
\providecommand{\barre}{\ensuremath{\|}}
% environnement proof (amsthm non chargé) utilisé par les modèles
\ifdefined\proof\else\newenvironment{proof}[1][Démonstration]{\par\noindent\textit{#1.}\ }{\qed\par}\fi
\usepackage{lastpage}
\usepackage{hyperref}
\hypersetup{hidelinks}

% ============================================================
% Palette « Pop » OnBuch+ — mêmes hex que la classe P (plus_ui.dart)
% ============================================================
\definecolor{poporange}{HTML}{F59321}
\definecolor{popdark}{HTML}{E07A0C}
\definecolor{popcream}{HTML}{FFF6E8}
\definecolor{popink}{HTML}{1C1714}
\definecolor{poppurple}{HTML}{7A5AE0}
\definecolor{popgreen}{HTML}{2E9E6A}
\definecolor{poppink}{HTML}{E0457B}
\definecolor{popblue}{HTML}{2F7DE1}
\definecolor{popgold}{HTML}{C98A12}
\definecolor{popmuted}{HTML}{8A7F76}
\definecolor{popline}{HTML}{EADFD2}
\definecolor{popgray}{HTML}{8A7F76}
\colorlet{popgrey}{popgray}
% Alias tolérants (le modèle nomme parfois les couleurs différemment)
\colorlet{popwhite}{white}
\colorlet{popblack}{popink}
\colorlet{popred}{poppink}
\colorlet{popyellow}{popgold}
% Teintes claires pour fonds de tableaux / zones
\colorlet{popblueL}{popblue!10!white}
\colorlet{poporangeL}{poporange!14!white}
\colorlet{popgreenL}{popgreen!12!white}
\colorlet{poppurpleL}{poppurple!12!white}
\colorlet{poppinkL}{poppink!12!white}
\colorlet{popgoldL}{popgold!14!white}

\pagecolor{popcream}

% ============================================================
% Typographie
% ============================================================
\defaultfontfeatures{Ligatures=TeX}
\setmainfont{PlusJakartaSans-Medium.ttf}[Path=../../../fonts/,BoldFont=PlusJakartaSans-Bold.ttf]
% Maths en sans-serif (Fira Math) avec chiffres et lettres droites de Plus
% Jakarta, pour que les formules se fondent dans le texte.
\usepackage[math-style=ISO,bold-style=ISO]{unicode-math}
\setmathfont{FiraMath-Regular.otf}[Path=../../../fonts/]
\setmathfont{PlusJakartaSans-Medium.ttf}[Path=../../../fonts/,range={up/num,up/{Latin,latin},\mathup}]
\usepackage{icomma} % virgule décimale sans espace en mode math
% Caractères mathématiques Unicode absents des polices de texte (Plus Jakarta,
% Archivo Black) : rendus par leur équivalent mathématique, en texte comme en
% formule — sinon ils s'affichent en carré vide (ex. « Arithmétique dans ℤ »).
\usepackage{newunicodechar}
\newunicodechar{ℝ}{\ensuremath{\mathbb{R}}}
\newunicodechar{ℤ}{\ensuremath{\mathbb{Z}}}
\newunicodechar{ℂ}{\ensuremath{\mathbb{C}}}
\newunicodechar{ℕ}{\ensuremath{\mathbb{N}}}
\newunicodechar{ℚ}{\ensuremath{\mathbb{Q}}}
\newunicodechar{𝔻}{\ensuremath{\mathbb{D}}}
\newunicodechar{α}{\ensuremath{\alpha}}
\newunicodechar{β}{\ensuremath{\beta}}
\newunicodechar{γ}{\ensuremath{\gamma}}
\newunicodechar{δ}{\ensuremath{\delta}}
\newunicodechar{ε}{\ensuremath{\varepsilon}}
\newunicodechar{θ}{\ensuremath{\theta}}
\newunicodechar{λ}{\ensuremath{\lambda}}
\newunicodechar{μ}{\ensuremath{\mu}}
\newunicodechar{σ}{\ensuremath{\sigma}}
\newunicodechar{τ}{\ensuremath{\tau}}
\newunicodechar{φ}{\ensuremath{\varphi}}
\newunicodechar{ψ}{\ensuremath{\psi}}
\newunicodechar{ω}{\ensuremath{\omega}}
\newunicodechar{Σ}{\ensuremath{\Sigma}}
% majuscules produites par \MakeUppercase dans les titres (α-aminés -> Α)
\newunicodechar{Α}{\ensuremath{\alpha}}
\newunicodechar{Β}{\ensuremath{\beta}}
\newunicodechar{Γ}{\ensuremath{\Gamma}}
\newunicodechar{Δ}{\ensuremath{\Delta}}
\newunicodechar{Ε}{\ensuremath{\varepsilon}}
\newunicodechar{Θ}{\ensuremath{\Theta}}
\newunicodechar{Λ}{\ensuremath{\Lambda}}
\newunicodechar{Μ}{\ensuremath{\mu}}
\newunicodechar{Τ}{\ensuremath{\tau}}
\newunicodechar{Φ}{\ensuremath{\Phi}}
\newunicodechar{Ψ}{\ensuremath{\Psi}}
\newunicodechar{Ω}{\ensuremath{\Omega}}
\newunicodechar{↦}{\ensuremath{\mapsto}}
\newunicodechar{∈}{\ensuremath{\in}}
\newunicodechar{∉}{\ensuremath{\notin}}
\newunicodechar{∧}{\ensuremath{\wedge}}
\newunicodechar{⇔}{\ensuremath{\Leftrightarrow}}
\newunicodechar{⟺}{\ensuremath{\Longleftrightarrow}}
\newunicodechar{⇒}{\ensuremath{\Rightarrow}}
\newunicodechar{⟹}{\ensuremath{\Longrightarrow}}
\newunicodechar{∘}{\ensuremath{\circ}}
\newunicodechar{∪}{\ensuremath{\cup}}
\newunicodechar{∩}{\ensuremath{\cap}}
\newunicodechar{≡}{\ensuremath{\equiv}}
\newunicodechar{⊂}{\ensuremath{\subset}}
\newunicodechar{⊥}{\ensuremath{\perp}}
\newunicodechar{∃}{\ensuremath{\exists}}
\newunicodechar{∀}{\ensuremath{\forall}}
\newunicodechar{∛}{\ensuremath{\sqrt[3]{\,}}}
\newunicodechar{∜}{\ensuremath{\sqrt[4]{\,}}}
\newunicodechar{⃗}{\ensuremath{{}^{\to}}}
\newunicodechar{ⁿ}{\ifmmode^{n}\else\textsuperscript{\ensuremath{n}}\fi}
\newunicodechar{ˣ}{\ifmmode^{x}\else\textsuperscript{\ensuremath{x}}\fi}
\newunicodechar{ᵇ}{\ifmmode^{b}\else\textsuperscript{\ensuremath{b}}\fi}
\newunicodechar{ʳ}{\ifmmode^{r}\else\textsuperscript{\ensuremath{r}}\fi}
\newunicodechar{ᵘ}{\ifmmode^{u}\else\textsuperscript{\ensuremath{u}}\fi}
\newunicodechar{ʸ}{\ifmmode^{y}\else\textsuperscript{\ensuremath{y}}\fi}
\newunicodechar{ᵃ}{\ifmmode^{a}\else\textsuperscript{\ensuremath{a}}\fi}
\newunicodechar{ᵉ}{\ifmmode^{e}\else\textsuperscript{\ensuremath{e}}\fi}
\newunicodechar{ᵏ}{\ifmmode^{k}\else\textsuperscript{\ensuremath{k}}\fi}
\newunicodechar{ᵖ}{\ifmmode^{p}\else\textsuperscript{\ensuremath{p}}\fi}
\newunicodechar{ᴺ}{\ifmmode^{N}\else\textsuperscript{\ensuremath{N}}\fi}
\newunicodechar{ᶜ}{\ifmmode^{c}\else\textsuperscript{\ensuremath{c}}\fi}
\newunicodechar{ʰ}{\ifmmode^{h}\else\textsuperscript{\ensuremath{h}}\fi}
\newunicodechar{ᵠ}{\ifmmode^{q}\else\textsuperscript{\ensuremath{q}}\fi}
\newunicodechar{ₙ}{\ifmmode_{n}\else\textsubscript{\ensuremath{n}}\fi}
\newunicodechar{ₐ}{\ifmmode_{a}\else\textsubscript{\ensuremath{a}}\fi}
\newunicodechar{ᵢ}{\ifmmode_{i}\else\textsubscript{\ensuremath{i}}\fi}
\newunicodechar{ᵣ}{\ifmmode_{r}\else\textsubscript{\ensuremath{r}}\fi}
\newunicodechar{ₖ}{\ifmmode_{k}\else\textsubscript{\ensuremath{k}}\fi}
\newunicodechar{ᵦ}{\ifmmode_{\beta}\else\textsubscript{\ensuremath{\beta}}\fi}
\newunicodechar{ₓ}{\ifmmode_{x}\else\textsubscript{\ensuremath{x}}\fi}
\newfontfamily\popdisplay{ArchivoBlack-Regular.ttf}[Path=../../../fonts/]
\newfontfamily\poplogo{SpaceGrotesk-Bold.ttf}[Path=../../../fonts/]
\newfontfamily\popsemi{PlusJakartaSans-SemiBold.ttf}[Path=../../../fonts/]
\newfontfamily\popxbold{PlusJakartaSans-ExtraBold.ttf}[Path=../../../fonts/]
\newfontfamily\popxboldls{PlusJakartaSans-ExtraBold.ttf}[Path=../../../fonts/,LetterSpace=3.0]

\setlist[enumerate]{leftmargin=1.7em,itemsep=5pt,topsep=4pt}
\setlist[itemize]{leftmargin=1.7em,itemsep=4pt,topsep=4pt,label={\color{poporange}\textbullet}}
\setlength{\parindent}{0pt}
\setlength{\parskip}{5pt}
\renewcommand{\arraystretch}{1.25}

% pgfplots : style Pop par défaut
\pgfplotsset{
  every axis/.append style={axis line style={popink,line width=1pt},tick style={popink},
    grid style={popline,line width=0.5pt},label style={font=\small},tick label style={font=\footnotesize}},
  popcurve/.style={poporange,line width=1.8pt,no markers,smooth},
}
% Même style disponible dans un \draw TikZ simple (hors pgfplots)
\tikzset{popcurve/.style={poporange,line width=1.8pt}}

% ============================================================
% Pill / squiggle
% ============================================================
\newcommand{\poppill}[3]{%
  \tikz[baseline=-0.55ex]\node[draw=popink,line width=1.2pt,rounded corners=2.6pt,%
    fill=#2,inner xsep=6.5pt,inner ysep=3pt]{%
    \popxboldls\fontsize{7}{7}\selectfont\color{#3}\MakeUppercase{#1}};%
}
\newcommand{\popsquiggle}[1]{%
  \tikz[baseline=-0.4ex]{
    \draw[#1,line width=2.6pt,line cap=round]
      plot [smooth] coordinates {(0,0.09) (0.34,0.01) (0.68,0.21) (1.02,0.01) (1.36,0.10)};
  }%
}
% Mot-clé surligné (marqueur orange sous le texte)
\newcommand{\cle}[1]{{\popxbold\color{popdark}#1}}

% ============================================================
% En-tête / pied
% ============================================================
\pagestyle{fancy}
\fancyhf{}
\renewcommand{\headrulewidth}{0pt}
\renewcommand{\footrulewidth}{0pt}
\lhead{\raisebox{-0.15ex}{\poplogo\fontsize{13}{13}\selectfont\color{popink}OnBuch\color{poporange}+}}
\rhead{\poppill{\DOCMATIERE\ \textbullet\ \DOCNIVEAU\ \textbullet\ Leçon \DOCLECON}{white}{popink}}
\lfoot{\popxbold\fontsize{7.6}{7.6}\selectfont\color{popmuted}\MakeUppercase{Cours signé OnBuch+}\ \textbullet\ {\popsemi\DOCTITRE}}
\rfoot{\tikz[baseline=-0.6ex]\node[rounded corners=8.5pt,fill=popink,minimum height=17pt,minimum width=17pt,inner xsep=5pt,inner sep=0]{\popxbold\fontsize{8}{8}\selectfont\color{popcream}\thepage\,/\,\pageref*{LastPage}};}

% ============================================================
% Titres
% ============================================================
\newcommand{\chaptitre}[2]{%
  \begin{tcolorbox}[enhanced,colback=poporange,colframe=popink,boxrule=2.6pt,arc=11pt,
      drop shadow=popink,left=15pt,right=15pt,top=12pt,bottom=11pt,before skip=3pt,after skip=18pt]
    \poppill{#2}{popink}{popcream}\par\vspace{9pt}
    {\raggedright\hyphenpenalty=10000\exhyphenpenalty=10000\popdisplay\fontsize{22}{24}\selectfont\color{popcream}\MakeUppercase{#1}\par}
  \end{tcolorbox}
}
% Section numérotée : pastille orange avec le numéro + titre Archivo
\newcounter{popsec}
\newcommand{\coursec}[1]{%
  \stepcounter{popsec}\setcounter{popsub}{0}%
  \par\vspace{16pt}\noindent
  \begin{minipage}{\linewidth}
  \tikz[baseline=-0.7ex]\node[draw=popink,line width=1.6pt,fill=poporange,rounded corners=4pt,minimum size=22pt,inner sep=0]{\popdisplay\fontsize{12}{12}\selectfont\color{popcream}\thepopsec};%
  \hspace{8pt}{\popdisplay\fontsize{15}{17}\selectfont\color{popink}\MakeUppercase{#1}}\par
  \vspace{4pt}\noindent{\color{poporange}\rule{36pt}{3.6pt}}
  \end{minipage}\par\nopagebreak\vspace{8pt}
}
\newcounter{popsub}
\newcommand{\courssub}[1]{%
  \stepcounter{popsub}%
  \par\vspace{9pt}\noindent{\popxbold\fontsize{11.5}{13}\selectfont\color{popink}{\color{poporange}\thepopsec.\thepopsub}\hspace{6pt}#1}\par\nopagebreak\vspace{4pt}
}

% ============================================================
% Boîtes pédagogiques — même recette : fond blanc, bordure épaisse colorée,
% ombre dure encre, badge plein. Titre optionnel [#1].
% ============================================================
\tcbset{popbase/.style={breakable,enhanced,colback=white,boxrule=2.3pt,arc=9pt,
  shadow={3pt}{-3pt}{0pt}{popink},left=14pt,right=14pt,top=11pt,bottom=11pt,before skip=11pt,after skip=15pt,
  fontupper=\popsemi\fontsize{10.6}{14.5}\selectfont\color{popink},
  fontlower=\popsemi\fontsize{10.6}{14.5}\selectfont\color{popink}}}
% Coche dessinée (le glyphe ✓ n'existe pas dans les polices du document)
\newcommand{\popcheck}[1]{\tikz[baseline=-0.5ex]\draw[#1,line width=1.6pt,line cap=round,line join=round](-0.13,0.0)--(-0.04,-0.09)--(0.13,0.1);}
\newcommand{\popboxhead}[3]{\poppill{#1}{#2}{white}\ifx\relax#3\relax\else\hspace{6pt}{\popxbold\fontsize{10.6}{12}\selectfont\color{popink}#3}\fi\par\vspace{7pt}}

\newtcolorbox{aretenir}[1][]{popbase,colframe=popgreen,before upper={\popboxhead{À retenir}{popgreen}{#1}}}
\newtcolorbox{exemplebox}[1][]{popbase,colframe=poporange,before upper={\popboxhead{Exemple}{poporange}{#1}}}
\newtcolorbox{definition}[1][]{popbase,colframe=popblue,before upper={\popboxhead{Définition}{popblue}{#1}}}
\newtcolorbox{propriete}[1][]{popbase,colframe=poppurple,before upper={\popboxhead{Propriété}{poppurple}{#1}}}
\newtcolorbox{methode}[1][]{popbase,colframe=popgold,before upper={\popboxhead{Méthode}{popgold}{#1}}}
\newtcolorbox{attention}[1][]{popbase,colframe=poppink,before upper={\popboxhead{Attention}{poppink}{#1}}}
\newtcolorbox{experience}[1][]{popbase,colframe=popblue,colback=popblueL,before upper={\popboxhead{Expérience}{popblue}{#1}}}
% « Le savais-tu ? » : carte violette pleine (contexte, culture, Cameroun)
\newtcolorbox{savaistu}[1][]{popbase,colback=poppurple,colframe=popink,
  fontupper=\popsemi\fontsize{10.4}{14.2}\selectfont\color{white},
  before upper={\poppill{Le savais-tu\,?}{white}{poppurple}\ifx\relax#1\relax\else\hspace{6pt}{\popxbold\color{white}#1}\fi\par\vspace{7pt}}}
% Exercice résolu : énoncé en haut, \tcblower puis la solution sur fond crème
\newcounter{popexo}\newcounter{popexr}
\newtcolorbox{exoresolu}[1][]{popbase,colframe=poporange,colbacklower=popcream,
  segmentation style={popink,line width=1.2pt,dashed},
  before upper={\stepcounter{popexr}\popboxhead{Exercice résolu \thepopexr}{poporange}{#1}},
  before lower={\poppill{Solution}{popink}{popcream}\par\vspace{6pt}}}
% Exercice (non résolu en place) — la correction est dans la partie Corrigés
\newtcolorbox{exercice}[1][]{popbase,colframe=popink,
  before upper={\stepcounter{popexo}\popboxhead{Exercice \thepopexo}{popink}{#1}}}
% Corrigé
\newtcolorbox{corrige}[1][]{popbase,colframe=popgreen,colback=popgreenL,
  before upper={\popboxhead{Corrigé}{popgreen}{#1}}}
% Objectifs / prérequis / bilan
\newtcolorbox{objectifs}{popbase,colframe=popink,colback=poporangeL,
  before upper={\popboxhead{Objectifs de la leçon}{popink}{}}}
\newtcolorbox{prerequis}{popbase,colframe=popink,colback=popblueL,
  before upper={\popboxhead{Prérequis}{popblue}{}}}
\newtcolorbox{bilan}{popbase,colframe=popink,colback=white,boxrule=2.8pt,
  before upper={\popboxhead{Fiche bilan}{poporange}{L'essentiel à connaître pour l'examen}}}
% Figure encadrée avec légende
\newtcolorbox{popfigure}[1][]{enhanced,breakable=false,colback=white,colframe=popline,boxrule=1.4pt,arc=8pt,
  left=10pt,right=10pt,top=10pt,bottom=8pt,before skip=10pt,after skip=12pt,halign=center,
  lower separated=false,halign lower=center,
  fontlower=\popsemi\fontsize{9}{11.5}\selectfont\color{popmuted},
  #1}
\newcommand{\legende}[1]{\par\vspace{6pt}{\popsemi\fontsize{9}{11.5}\selectfont\color{popmuted}#1}}

% ============================================================
% Signature OnBuch+
% ============================================================
\newcommand{\poplogoinline}[1]{{\poplogo\fontsize{#1}{#1}\selectfont\color{popink}OnBuch\color{poporange}+}}
\newcommand{\signatureonbuch}{%
  \begin{tcolorbox}[enhanced,colback=popink,colframe=popink,boxrule=2.2pt,arc=10pt,
      drop shadow=poporange,left=14pt,right=14pt,top=10pt,bottom=10pt,before skip=0pt,after skip=0pt]
    \tikz[baseline=-0.7ex]\node[circle,fill=popgreen,draw=popcream,line width=1.3pt,minimum size=22pt,inner sep=0]{\popcheck{white}};%
    \hspace{9pt}\begin{minipage}[c]{0.62\linewidth}
      {\popxbold\fontsize{12}{13}\selectfont\color{popcream}Signé\ \textbullet\ {\poplogo OnBuch\color{poporange}+}}\par\vspace{2pt}
      {\raggedright\popsemi\fontsize{8}{10.5}\selectfont\color{popline}\MakeUppercase{Équipe pédagogique OnBuch+}\\\MakeUppercase{Conforme au programme officiel MINESEC}\par}
    \end{minipage}\hfill
    {\poplogo\fontsize{22}{22}\selectfont\color{popcream}OnBuch\color{poporange}+}
  \end{tcolorbox}%
}

% ============================================================
% Page de garde
% #1 titre · #2 matière · #3 niveau · #4 module · #5 n° leçon · #6 durée ·
% #7 description / objectifs (texte ou itemize)
% ============================================================
\newcommand{\pagegarde}[7]{%
  \thispagestyle{empty}
  \newgeometry{a4paper,margin=1.7cm,top=1.6cm,bottom=1.4cm}%
  \begin{tikzpicture}[remember picture,overlay]
    \fill[poporange!18] ([xshift=-3.2cm,yshift=-3.6cm]current page.north east) circle (4.6cm);
    \fill[poppurple!14] ([xshift=2.4cm,yshift=7.6cm]current page.south west) circle (3.8cm);
  \end{tikzpicture}%
  {\poplogo\fontsize{30}{30}\selectfont\color{popink}OnBuch\color{poporange}+}\hfill
  \raisebox{0.6ex}{\poppill{Cours signé \textbullet\ Premium}{popink}{popcream}}\par
  \vspace{1pt}\popsquiggle{poppurple}\par
  \vspace{8pt}
  \begin{tcolorbox}[enhanced,colback=poporange,colframe=popink,boxrule=2.8pt,arc=13pt,
      drop shadow=popink,left=18pt,right=18pt,top=12pt,bottom=13pt,after skip=10pt]
    \poppill{#2 \textbullet\ #3}{popink}{popcream}\hspace{5pt}\poppill{#4}{white}{popink}\par\vspace{8pt}
    {\popdisplay\fontsize{14}{14}\selectfont\color{popink}LEÇON #5}\par\vspace{4pt}
    {\raggedright\hyphenpenalty=10000\exhyphenpenalty=10000\popdisplay\fontsize{25}{27}\selectfont\color{popcream}\MakeUppercase{#1}\par}
    \vspace{8pt}
    {\popxbold\fontsize{9.5}{9.5}\selectfont\color{popink}\MakeUppercase{Durée conseillée : #6}}
  \end{tcolorbox}
  \begin{tcolorbox}[enhanced,colback=white,colframe=popink,boxrule=2.2pt,arc=10pt,
      drop shadow=popink,left=16pt,right=16pt,top=10pt,bottom=10pt,after skip=8pt]
    \poppill{Dans cette leçon}{poppurple}{white}\par\vspace{5pt}
    \popsemi\fontsize{9.3}{12.6}\selectfont\color{popink}#7
  \end{tcolorbox}
  \noindent\begin{minipage}[t]{0.487\linewidth}
    \begin{tcolorbox}[enhanced,colback=popgreen,colframe=popink,boxrule=2.2pt,arc=10pt,
        drop shadow=popink,left=11pt,right=11pt,top=10pt,bottom=10pt,height=3.1cm,valign=center]
      \tcbox[colback=white,colframe=popink,boxrule=1.3pt,arc=5pt,boxsep=0pt,nobeforeafter,
        left=3pt,right=3pt,top=3pt,bottom=3pt,valign=center]{\qrcode[height=1.9cm]{https://play.google.com/store/apps/details?id=com.luvvix.onbuch&hl=fr}}%
      \hspace{8pt}\begin{minipage}[c]{0.5\linewidth}
        {\popdisplay\fontsize{11}{12.5}\selectfont\color{white}\MakeUppercase{Télécharge OnBuch}}\par\vspace{4pt}
        {\raggedright\popsemi\fontsize{8.4}{11}\selectfont\color{white}Gratuit \textbullet\ Google Play\\Tuteur IA, annales, concours, résultats\par}
      \end{minipage}
    \end{tcolorbox}
  \end{minipage}\hfill
  \begin{minipage}[t]{0.487\linewidth}
    \begin{tcolorbox}[enhanced,colback=poppurple,colframe=popink,boxrule=2.2pt,arc=10pt,
        drop shadow=popink,left=11pt,right=11pt,top=10pt,bottom=10pt,height=3.1cm,valign=center]
      \tikz[baseline=-0.7ex]\node[circle,fill=white,draw=popink,line width=1.3pt,minimum size=1.6cm,inner sep=0]{\popdisplay\fontsize{24}{24}\selectfont\color{poppurple}+};%
      \hspace{8pt}\begin{minipage}[c]{0.6\linewidth}
        {\popdisplay\fontsize{11}{12.5}\selectfont\color{white}\MakeUppercase{Passe à OnBuch+}}\par\vspace{4pt}
        {\raggedright\popsemi\fontsize{8.4}{11}\selectfont\color{white}Tous les cours signés, Tuteur IA illimité, fiches PDF — onglet OnBuch+ de l'app\par}
      \end{minipage}
    \end{tcolorbox}
  \end{minipage}
  \vspace{8pt}
  \popcontenu
  \vfill
  \signatureonbuch
  \restoregeometry
  \newpage
}

% Rangée « Ce cours contient » (page de garde)
\newcommand{\poptile}[3]{% #1 couleur #2 gros texte #3 légende
  \begin{tcolorbox}[enhanced,colback=#1,colframe=popink,boxrule=2pt,arc=9pt,shadow={2.5pt}{-2.5pt}{0pt}{popink},
      left=6pt,right=6pt,top=9pt,bottom=9pt,halign=center,valign=center,height=1.75cm,width=\linewidth,nobeforeafter]
    {\popdisplay\fontsize{12.5}{13}\selectfont\color{white}\MakeUppercase{#2}}\par\vspace{3pt}
    {\centering\popsemi\fontsize{7.8}{9.5}\selectfont\color{white}#3\par}
  \end{tcolorbox}}
\newcommand{\popcontenu}{%
  \par\noindent{\popxboldls\fontsize{7.5}{7.5}\selectfont\color{popmuted}CE COURS SIGNÉ CONTIENT}\par\vspace{7pt}
  \noindent\begin{minipage}[t]{0.235\linewidth}\poptile{popblue}{Cours}{complet, illustré, pas à pas}\end{minipage}\hfill
  \begin{minipage}[t]{0.235\linewidth}\poptile{poporange}{Méthodes}{et exercices résolus}\end{minipage}\hfill
  \begin{minipage}[t]{0.235\linewidth}\poptile{poppink}{Exercices}{type examen + corrigés}\end{minipage}\hfill
  \begin{minipage}[t]{0.235\linewidth}\poptile{popgreen}{Fiche bilan}{l'essentiel à retenir}\end{minipage}\par}

% Sommaire
\newtcolorbox{sommaire}{enhanced,colback=white,colframe=popink,boxrule=2.2pt,arc=10pt,
  drop shadow=popink,left=18pt,right=18pt,top=13pt,bottom=13pt,before skip=6pt,after skip=20pt,
  before upper={\poppill{Sommaire}{poppurple}{white}\par\vspace{9pt}\popsemi\fontsize{10.6}{14.5}\selectfont\color{popink}}}

% Signature de fin de cours — poussée tout en bas de la dernière page
\newcommand{\findecours}{%
  \par\vfill\nopagebreak
  \begin{center}\popsquiggle{poporange}\end{center}\vspace{4pt}
  \signatureonbuch
}
\AtBeginDocument{\def\llbracket{\mathopen{[\mkern-3.5mu[}}\def\rrbracket{\mathclose{]\mkern-3.5mu]}}} % intervalles d'entiers (glyphe ⟦ absent de la police)
% Glyphes absents de Fira Math : redéfinis à partir de glyphes disponibles
\AtBeginDocument{%
  \def\setminus{\mathbin{\backslash}}\let\smallsetminus\setminus
  \def\vdots{\mathord{\vbox{\baselineskip3.2pt\lineskiplimit0pt\kern4pt\hbox{.}\hbox{.}\hbox{.}}}}%
  \def\ddots{\mathinner{\mkern1mu\raise6.5pt\hbox{.}\mkern2mu\raise3.6pt\hbox{.}\mkern2mu\raise0.7pt\hbox{.}\mkern1mu}}%
  \def\triangle{\mathord{\scalebox{0.9}{$\symup{\Delta}$}}}%
  \def\nabla{\mathord{\raisebox{\depth}{\scalebox{1}[-1]{$\symup{\Delta}$}}}}%
  \def\checkmark{\popcheck{popdark}}%
  \def\star{\mathbin{\ast}}\def\bigstar{\mathord{\ast}}%
  \def\diamond{\mathbin{\scalebox{0.6}{$\Diamond$}}}%
  \def\bigcup{\mathop{\mathchoice{\raisebox{-0.3em}{\scalebox{1.7}{$\cup$}}}{\scalebox{1.25}{$\cup$}}{\cup}{\cup}}\displaylimits}%
  \def\bigcap{\mathop{\mathchoice{\raisebox{-0.3em}{\scalebox{1.7}{$\cap$}}}{\scalebox{1.25}{$\cap$}}{\cap}{\cap}}\displaylimits}%
  \def\lesseqqgtr{\mathrel{\lessgtr}}\def\bigoplus{\mathop{\oplus}\displaylimits}%
}
\providecommand{\ssi}{si et seulement si} % abréviation fréquente
\AtBeginDocument{\providecommand{\wideparen}{\overparen}} % arc de cercle

%% FIGFIX-BEGIN v4 — correctifs globaux des figures (docs/onbuch-generation/tools/figfix)
\usepackage{adjustbox}\usepackage{etoolbox}
% toute figure TikZ/pgfplots plus large que la ligne est réduite à la largeur disponible
\BeforeBeginEnvironment{tikzpicture}{\begin{adjustbox}{max width=\linewidth}}
\AfterEndEnvironment{tikzpicture}{\end{adjustbox}}
% légendes pgfplots lisibles par-dessus les courbes
\pgfplotsset{every axis/.append style={legend style={fill=white,fill opacity=0.92,text opacity=1,draw=black!15,inner sep=3pt}}}
% étiquettes d'arêtes (syntaxe edge["..."]) avec fond blanc
\tikzset{every edge quotes/.append style={fill=white,inner sep=1.2pt}}
%% FIGFIX-END
\begin{document}

\pagegarde{Angles orientés et arcs capables}{Mathématiques}{1ère C}{Configurations et transformations élémentaires du plan}{N15}{8 h}{Cette leçon approfondit les angles orientés de vecteurs et introduit la notion d'arc capable, ensemble des points d'où un segment est vu sous un angle donné. Elle fournit des outils puissants pour résoudre des problèmes de lieux géométriques, de constructions et de cocyclicité, très prisés au Probatoire C et E.
\vspace{4pt}
\begin{multicols}{2}\small
\begin{itemize}
\item À la fin de cette leçon, je sais définir et utiliser la mesure principale d'un angle orienté de vecteurs
\item je sais utiliser la relation de Chasles sur les angles orientés pour prouver des alignements, des orthogonalités et des cocyclicités
\item je sais énoncer et utiliser le théorème de l'angle inscrit et celui de l'angle formé par une corde et une demi-tangente
\item je sais déterminer et construire l'ensemble des points M tels que mes(AMB) = x, avec x ∈ [0 ; π]
\item je sais décrire et construire, dans le plan orienté, l'ensemble des points M tels que mes(MA, MB) = x + 2kπ
\item je sais décrire et construire l'ensemble des points M tels que mes(MA, MB) = x + kπ
\end{itemize}\end{multicols}}

\begin{sommaire}
\begin{multicols}{2}
\begin{enumerate}
\item Rappels sur les angles orientés de vecteurs
\item Angle inscrit, angle au centre et demi-tangente
\item Arc capable : ensemble des points M tels que mes(AMB) = x
\item Arcs capables orientés : mes(MA, MB) = x + 2kπ et mes(MA, MB) = x + kπ
\item Applications : lieux géométriques et cocyclicité
\item Activité d'intégration
\item Méthodes et exercices résolus
\item Exercices
\item Corrigés des exercices
\item Fiche bilan
\end{enumerate}
\end{multicols}
\end{sommaire}

\begin{objectifs}
\begin{itemize}
\item À la fin de cette leçon, je sais définir et utiliser la mesure principale d'un angle orienté de vecteurs
\item je sais utiliser la relation de Chasles sur les angles orientés pour prouver des alignements, des orthogonalités et des cocyclicités
\item je sais énoncer et utiliser le théorème de l'angle inscrit et celui de l'angle formé par une corde et une demi-tangente
\item je sais déterminer et construire l'ensemble des points M tels que mes(AMB) = x, avec x ∈ [0 ; π]
\item je sais décrire et construire, dans le plan orienté, l'ensemble des points M tels que mes(MA, MB) = x + 2kπ
\item je sais décrire et construire l'ensemble des points M tels que mes(MA, MB) = x + kπ
\item je sais utiliser les arcs capables pour résoudre des problèmes de lieux géométriques
\item je sais caractériser la cocyclicité de quatre points à l'aide des angles orientés
\end{itemize}
\end{objectifs}

\begin{prerequis}
\begin{itemize}
\item Angles orientés de vecteurs : définition, mesure principale, relation de Chasles (début d'année de Première)
\item Trigonométrie : cercle trigonométrique, valeurs remarquables des lignes trigonométriques
\item Géométrie du triangle et du cercle de Seconde : angle inscrit, angle au centre, médiatrice, tangente à un cercle
\item Colinéarité et orthogonalité de vecteurs, produit scalaire (notions de base)
\end{itemize}
\end{prerequis}

\newpage

\coursec{Rappels sur les angles orientés de vecteurs}

\textbf{Situation d'accroche.} Un technicien de la SODECAO à Bafoussam veut installer un projecteur au point $M$ d'une colline pour éclairer un champ délimité par deux poteaux $A$ et $B$, avec un angle de visée constant de $60^\circ$. Où peut-il placer son projecteur ? Plus au sud, un capitaine de pirogue à Kribi doit garder un cap tel que l'angle sous lequel il voit deux balises reste constant, afin d'éviter un haut-fond dangereux. Ces deux problèmes, en apparence très différents, se ramènent à une même question géométrique : \emph{quel est l'ensemble des points du plan d'où l'on voit un segment $[AB]$ sous un angle donné ?}

Pour y répondre, il ne suffit pas de savoir mesurer un angle « en valeur absolue » comme au collège : il faut pouvoir dire si l'on tourne dans le \cle{sens direct} (sens inverse des aiguilles d'une montre) ou dans le \cle{sens indirect}. C'est précisément l'objet des \cle{angles orientés}, outil fondamental que tu as découvert en classe de Seconde et que nous allons consolider ici avant d'attaquer les arcs capables.

\textbf{Problématique de la leçon.} Comment mesurer rigoureusement l'écart angulaire entre deux directions, et comment exploiter ces mesures pour décrire des lieux géométriques ?

\courssub{Le plan orienté et le cercle trigonométrique}

Pour parler d'angle orienté, il faut d'abord \emph{orienter} le plan, c'est-à-dire choisir une convention de rotation.

\begin{definition}[Plan orienté, cercle trigonométrique]
Orienter le plan, c'est choisir un sens de rotation appelé \cle{sens direct} (ou sens trigonométrique) : par convention, c'est le sens inverse des aiguilles d'une montre. Le sens opposé est dit \cle{sens indirect} (ou rétrograde).

Le \cle{cercle trigonométrique} est le cercle de centre $O$ et de rayon $1$, orienté dans le sens direct. On y mesure les angles en \cle{radians} : un tour complet vaut $2\pi$ radians, un demi-tour $\pi$ radians, un quart de tour direct $\dfrac{\pi}{2}$ radians.
\end{definition}

\begin{savaistu}[Pourquoi le radian ?]
Le radian est l'unité « naturelle » des mathématiciens : sur le cercle trigonométrique, un angle de $1$ radian intercepte un arc de longueur exactement $1$ (le rayon). C'est ce qui rendra les formules de dérivation des fonctions trigonométriques si simples en Terminale. Retiens les correspondances : $180^\circ = \pi$ rad, $60^\circ = \dfrac{\pi}{3}$ rad, $45^\circ = \dfrac{\pi}{4}$ rad, $30^\circ = \dfrac{\pi}{6}$ rad.
\end{savaistu}

\courssub{Angle orienté de deux vecteurs non nuls}

\textbf{Intuition.} Prends deux vecteurs non nuls $\vec{u}$ et $\vec{v}$. Place-toi le long de $\vec{u}$ et regarde dans sa direction : pour t'aligner sur la direction de $\vec{v}$, tu dois tourner d'un certain angle. Mais attention : tu peux tourner dans le sens direct, ou faire presque un tour complet dans le sens indirect ! Il y a donc une \emph{infinité} de mesures possibles, qui diffèrent toutes d'un nombre entier de tours, c'est-à-dire d'un multiple de $2\pi$.

\begin{definition}[Angle orienté de deux vecteurs]
Soient $\vec{u}$ et $\vec{v}$ deux vecteurs \textbf{non nuls} du plan orienté. L'\cle{angle orienté} du couple $(\vec{u}, \vec{v})$, noté $(\vec{u}, \vec{v})$, est l'angle dont les mesures sont les nombres $\theta$ tels que la rotation d'angle $\theta$ transforme la direction de $\vec{u}$ en celle de $\vec{v}$.

Si $\theta_0$ est une mesure de $(\vec{u}, \vec{v})$, alors \textbf{toutes} ses mesures sont les nombres $\theta_0 + 2k\pi$ avec $k \in \mathbb{Z}$. On écrit :
\[ (\vec{u}, \vec{v}) = \theta_0 \ [2\pi] \qquad \text{ou} \qquad \operatorname{mes}(\vec{u}, \vec{v}) = \theta_0 + 2k\pi,\ k \in \mathbb{Z}. \]
Parmi toutes ces mesures, il en existe une et une seule dans l'intervalle $]-\pi \, ; \pi]$ : c'est la \cle{mesure principale} de l'angle.
\end{definition}

\begin{attention}[L'angle orienté n'existe que pour des vecteurs non nuls]
Le vecteur nul $\vec{0}$ n'a pas de direction : l'écriture $(\vec{u}, \vec{0})$ n'a \textbf{aucun sens}. Avant d'écrire un angle orienté, vérifie toujours que les deux vecteurs sont non nuls. De même, dans un angle $(\overrightarrow{AB}, \overrightarrow{AC})$, les points doivent être distincts : $A \neq B$ et $A \neq C$.
\end{attention}

\begin{popfigure}
\begin{center}
\begin{tikzpicture}[scale=2.2]
  % Cercle trigonométrique
  \draw[popline, thick] (0,0) circle (1);
  \draw[->, popmuted] (-1.3,0) -- (1.3,0);
  \draw[->, popmuted] (0,-1.3) -- (0,1.3);
  \fill[popdark] (0,0) circle (0.015) node[below left=1pt] {$O$};
  % Vecteur u
  \draw[->, very thick, popblue] (0,0) -- (1,0) node[below right] {$\vec{u}$};
  % Vecteur v
  \draw[->, very thick, poporange] (0,0) -- ({cos(70)},{sin(70)}) node[above right] {$\vec{v}$};
  % Arc fléché mesure principale
  \draw[->, very thick, popgreen] (0.55,0) arc (0:70:0.55);
  \node[popgreen] at (35:0.82) {$\theta$};
  % Sens direct rappel
  \draw[->, popmuted] (150:1.12) arc (150:120:1.12);
  \node[popmuted] at (135:1.28) {\footnotesize sens direct};
\end{tikzpicture}
\end{center}
\legende{La mesure principale $\theta$ de $(\vec{u}, \vec{v})$ est l'unique mesure dans $]-\pi\,;\pi]$. Ici $\theta = \dfrac{7\pi}{18}$ ($70^\circ$), dans le sens direct.}
\end{popfigure}

\begin{exemplebox}[Réduction à la mesure principale]
Déterminons la mesure principale de deux angles dont on connaît une mesure.

\textbf{1.} $\theta = \dfrac{13\pi}{4}$. On retire des tours complets ($2\pi = \dfrac{8\pi}{4}$) :
\[ \frac{13\pi}{4} = \frac{8\pi}{4} + \frac{5\pi}{4} = 2\pi + \frac{5\pi}{4}. \]
Or $\dfrac{5\pi}{4} > \pi$, donc ce n'est pas encore la mesure principale : on retire encore un tour :
\[ \frac{5\pi}{4} - 2\pi = -\frac{3\pi}{4} \in \left]-\pi\,;\pi\right]. \]
La mesure principale est $\boxed{-\dfrac{3\pi}{4}}$.

\textbf{2.} $\theta = -\dfrac{7\pi}{6}$. On ajoute un tour :
\[ -\frac{7\pi}{6} + 2\pi = \frac{5\pi}{6} \in \left]-\pi\,;\pi\right]. \]
La mesure principale est $\boxed{\dfrac{5\pi}{6}}$.
\end{exemplebox}

\begin{methode}[Déterminer la mesure principale d'un angle]
Pour réduire une mesure quelconque $\theta$ à sa mesure principale :
\begin{enumerate}
  \item Écrire $\theta$ sous la forme $\theta = \theta_0 + 2k\pi$ en ajoutant ou retranchant des multiples de $2\pi$ (astuce : mettre au même dénominateur, $2\pi = \dfrac{2\pi n}{n}$).
  \item S'arrêter dès que $\theta_0 \in \left]-\pi \, ; \pi\right]$.
  \item Vérifier mentalement le signe : une mesure principale négative correspond à un angle dans le sens indirect (horaire).
\end{enumerate}
\end{methode}

\courssub{Relation de Chasles et propriétés opératoires}

Toute la puissance du calcul sur les angles orientés repose sur une propriété qui ressemble à la relation de Chasles sur les vecteurs : les angles orientés \emph{s'additionnent}.

\begin{propriete}[Relation de Chasles pour les angles orientés — admise]
Pour tous vecteurs non nuls $\vec{u}$, $\vec{v}$, $\vec{w}$ :
\[ (\vec{u}, \vec{v}) + (\vec{v}, \vec{w}) = (\vec{u}, \vec{w}) \ [2\pi]. \]
Cette propriété est admise au Probatoire ; elle se justifie en additionnant les rotations : passer de la direction de $\vec{u}$ à celle de $\vec{v}$, puis de celle de $\vec{v}$ à celle de $\vec{w}$, revient à passer directement de $\vec{u}$ à $\vec{w}$.
\end{propriete}

\begin{propriete}[Conséquences de la relation de Chasles]
Pour tous vecteurs non nuls $\vec{u}$, $\vec{v}$ et tous réels $a, b$ strictement positifs :
\begin{itemize}
  \item $(\vec{u}, \vec{u}) = 0 \ [2\pi]$ et $(\vec{u}, -\vec{u}) = \pi \ [2\pi]$ ;
  \item $(\vec{u}, \vec{v}) = -(\vec{v}, \vec{u}) \ [2\pi]$ \quad (l'angle orienté change de signe quand on échange les vecteurs) ;
  \item $(a\vec{u}, b\vec{v}) = (\vec{u}, \vec{v}) \ [2\pi]$ : multiplier les vecteurs par des réels \textbf{strictement positifs} ne change pas l'angle (on ne change ni la direction ni le sens) ;
  \item $(-\vec{u}, \vec{v}) = (\vec{u}, -\vec{v}) = (\vec{u}, \vec{v}) + \pi \ [2\pi]$ : changer \textbf{un seul} vecteur en son opposé ajoute $\pi$ à l'angle ;
  \item $(-\vec{u}, -\vec{v}) = (\vec{u}, \vec{v}) \ [2\pi]$.
\end{itemize}
\end{propriete}

\begin{experience}[Démonstration guidée : $(\vec{u}, \vec{v}) = -(\vec{v}, \vec{u})$]
Montrons cette propriété à l'aide de la relation de Chasles.
\begin{enumerate}
  \item Écris la relation de Chasles avec les vecteurs $\vec{u}$, $\vec{v}$, $\vec{u}$ :
  \[ (\vec{u}, \vec{v}) + (\vec{v}, \vec{u}) = (\vec{u}, \vec{u}) \ [2\pi]. \]
  \item Or $(\vec{u}, \vec{u}) = 0 \ [2\pi]$ (aucune rotation n'est nécessaire pour passer de $\vec{u}$ à lui-même).
  \item Donc $(\vec{u}, \vec{v}) + (\vec{v}, \vec{u}) = 0 \ [2\pi]$, c'est-à-dire $(\vec{u}, \vec{v}) = -(\vec{v}, \vec{u}) \ [2\pi]$. \hfill $\square$
\end{enumerate}
Essaie de prouver de même que $(-\vec{u}, \vec{v}) = (\vec{u}, \vec{v}) + \pi \ [2\pi]$ en écrivant $(-\vec{u}, \vec{v}) = (-\vec{u}, \vec{u}) + (\vec{u}, \vec{v})$.
\end{experience}

\begin{attention}[Deux erreurs classiques à bannir]
\begin{itemize}
  \item \textbf{Confondre $(\vec{u}, \vec{v})$ et $(\vec{v}, \vec{u})$.} Ce sont des angles \emph{opposés} : si $(\vec{u}, \vec{v}) = \dfrac{\pi}{3} \ [2\pi]$, alors $(\vec{v}, \vec{u}) = -\dfrac{\pi}{3} \ [2\pi]$, et non $\dfrac{\pi}{3}$ ! L'ordre des vecteurs compte.
  \item \textbf{Oublier le modulo.} Une égalité d'angles orientés de vecteurs est toujours vraie \emph{modulo $2\pi$} : écrire « $(\vec{u}, \vec{v}) = \dfrac{\pi}{4}$ » sans le $[2\pi]$ est une faute de rédaction au Probatoire. Pour les droites (voir plus bas), le modulo est $\pi$.
\end{itemize}
\end{attention}

\courssub{Caractérisations angulaires : colinéarité et orthogonalité}

Les angles orientés fournissent un langage très efficace pour traduire des propriétés géométriques.

\begin{propriete}[Colinéarité et orthogonalité]
Soient $\vec{u}$ et $\vec{v}$ deux vecteurs non nuls.
\begin{itemize}
  \item $\vec{u}$ et $\vec{v}$ sont \cle{colinéaires} si et seulement si
  \[ (\vec{u}, \vec{v}) = 0 \ [\pi] \quad \text{(c'est-à-dire } 0 \ [2\pi] \text{ pour le même sens, ou } \pi \ [2\pi] \text{ pour des sens opposés).} \]
  \item $\vec{u}$ et $\vec{v}$ sont \cle{orthogonaux} si et seulement si
  \[ (\vec{u}, \vec{v}) = \frac{\pi}{2} \ [\pi] \quad \text{(c'est-à-dire } \frac{\pi}{2} \text{ ou } -\frac{\pi}{2} \text{ modulo } 2\pi\text{).} \]
\end{itemize}
\end{propriete}

\textbf{Justification.} Dire que $\vec{u}$ et $\vec{v}$ sont colinéaires signifie qu'ils ont la même direction : la rotation qui envoie la direction de $\vec{u}$ sur celle de $\vec{v}$ est donc d'angle $0$ (même sens) ou $\pi$ (sens opposés), c'est-à-dire un multiple de $\pi$. De même, l'orthogonalité signifie que les directions font entre elles un quart de tour, dans un sens ou dans l'autre : $\dfrac{\pi}{2}$ ou $-\dfrac{\pi}{2}$, soit $\dfrac{\pi}{2}$ modulo $\pi$. \hfill $\square$

\courssub{Angle orienté de droites}

Deux droites ne portent pas de sens de parcours : la droite $(AB)$ est aussi la droite $(BA)$. Passer d'une droite à l'autre revient à ajouter $\pi$ à l'angle des vecteurs directeurs. C'est pourquoi l'angle de deux droites n'est défini que \emph{modulo $\pi$}.

\begin{definition}[Angle orienté de droites]
Soient $\Delta$ et $\Delta'$ deux droites de vecteurs directeurs respectifs $\vec{u}$ et $\vec{v}$ (non nuls). L'\cle{angle orienté des droites} $(\Delta, \Delta')$ est défini par :
\[ (\Delta, \Delta') = (\vec{u}, \vec{v}) \ [\pi]. \]
Le résultat ne dépend pas du choix des vecteurs directeurs, car changer $\vec{u}$ en $-\vec{u}$ ajoute $\pi$, ce qui est absorbé par le modulo $\pi$.
\end{definition}

\begin{aretenir}[Les deux modulos à connaître par cœur]
\begin{center}
\begin{tabular}{|l|c|c|}
\hline
\rowcolor{popblueL} \textbf{Objet} & \textbf{Notation} & \textbf{Égalité modulo} \\
\hline
Angle de vecteurs $(\vec{u}, \vec{v})$ & $(\vec{u}, \vec{v})$ & $2\pi$ \\
\hline
Angle de droites $(\Delta, \Delta')$ & $(\Delta, \Delta')$ & $\pi$ \\
\hline
\end{tabular}
\end{center}
Conséquence immédiate : deux droites sont perpendiculaires si et seulement si $(\Delta, \Delta') = \dfrac{\pi}{2} \ [\pi]$, et parallèles si et seulement si $(\Delta, \Delta') = 0 \ [\pi]$.
\end{aretenir}

\courssub{Lecture d'angles orientés dans les figures usuelles}

Savoir lire rapidement un angle orienté dans une figure de référence est un savoir-faire exigible. La règle d'or : \textbf{identifier d'abord le sens} (direct ou indirect), \textbf{puis la valeur absolue} de l'angle géométrique.

\begin{popfigure}
\begin{center}
\begin{tikzpicture}[scale=2.6]
  % Triangle équilatéral ABC direct
  \coordinate (A) at (0,0);
  \coordinate (B) at (2,0);
  \coordinate (C) at (1,1.732);
  \draw[thick, popdark] (A) -- (B) -- (C) -- cycle;
  \fill[popdark] (A) circle (0.02) node[below left] {$A$};
  \fill[popdark] (B) circle (0.02) node[below right] {$B$};
  \fill[popdark] (C) circle (0.02) node[above] {$C$};
  % Vecteurs AB et AC
  \draw[->, very thick, popblue] (A) -- ($(A)!0.55!(B)$);
  \draw[->, very thick, poporange] (A) -- ($(A)!0.55!(C)$);
  \draw[->, thick, popgreen] (0.75,0) arc (0:60:0.75);
  \node[popgreen] at (30:1.0) {$\frac{\pi}{3}$};
  % Vecteurs BA et BC
  \draw[->, very thick, popblue] (B) -- ($(B)!0.45!(A)$);
  \draw[->, very thick, poppink] (B) -- ($(B)!0.45!(C)$);
  % Arc de 180 à 120 (sens horaire = indirect)
  \draw[->, thick, poppurple] ($(B)+(180:0.6)$) arc (180:120:0.6);
  \node[poppurple] at ($(B)+(150:0.9)$) {$-\frac{\pi}{3}$};
\end{tikzpicture}
\end{center}
\legende{Triangle équilatéral $ABC$ direct : $(\overrightarrow{AB}, \overrightarrow{AC}) = \dfrac{\pi}{3} \ [2\pi]$ et $(\overrightarrow{BA}, \overrightarrow{BC}) = -\dfrac{\pi}{3} \ [2\pi]$.}
\end{popfigure}

\begin{exoresolu}[Angles orientés dans un triangle équilatéral]
Soit $ABC$ un triangle équilatéral \emph{direct} (c'est-à-dire que le sens $A \to B \to C$ est le sens direct). Déterminer les mesures principales de $(\overrightarrow{AB}, \overrightarrow{AC})$ et de $(\overrightarrow{BA}, \overrightarrow{BC})$.
\tcblower
\textbf{Solution détaillée.}

\textbf{1. Angle $(\overrightarrow{AB}, \overrightarrow{AC})$.} Les deux vecteurs partent de $A$. L'angle géométrique $\widehat{BAC}$ vaut $60^\circ = \dfrac{\pi}{3}$. Comme le triangle est direct, on passe de $\overrightarrow{AB}$ à $\overrightarrow{AC}$ en tournant dans le sens direct :
\[ (\overrightarrow{AB}, \overrightarrow{AC}) = \frac{\pi}{3} \ [2\pi]. \]
La mesure principale est $\dfrac{\pi}{3}$.

\textbf{2. Angle $(\overrightarrow{BA}, \overrightarrow{BC})$.} Les deux vecteurs partent de $B$. L'angle géométrique $\widehat{ABC}$ vaut aussi $\dfrac{\pi}{3}$. En $B$, le vecteur $\overrightarrow{BA}$ pointe vers $A$ (direction Ouest, angle $\pi$) et $\overrightarrow{BC}$ pointe vers $C$ (direction Nord-Ouest, angle $\frac{2\pi}{3}$). Pour passer de la direction de $\overrightarrow{BA}$ à celle de $\overrightarrow{BC}$, on tourne de $-\dfrac{\pi}{3}$ (sens indirect, horaire) :
\[ (\overrightarrow{BA}, \overrightarrow{BC}) = -\frac{\pi}{3} \ [2\pi]. \]
La mesure principale est $-\dfrac{\pi}{3}$.

\textbf{Vérification par le calcul (relation de Chasles).}
\begin{align*}
(\overrightarrow{BA}, \overrightarrow{BC}) &= (-\overrightarrow{AB}, \overrightarrow{BC}) \\
&= (\overrightarrow{AB}, \overrightarrow{BC}) + \pi \ [2\pi] \quad \text{(changer un vecteur en son opposé ajoute $\pi$)}.
\end{align*}
Calculons $(\overrightarrow{AB}, \overrightarrow{BC})$ par Chasles via $A$ et $C$ :
\begin{align*}
(\overrightarrow{AB}, \overrightarrow{BC}) &= (\overrightarrow{AB}, \overrightarrow{AC}) + (\overrightarrow{AC}, \overrightarrow{BC}) \ [2\pi].
\end{align*}
On a $(\overrightarrow{AB}, \overrightarrow{AC}) = \dfrac{\pi}{3}$. Le vecteur $\overrightarrow{AC}$ a pour direction $60^\circ$ ($\frac{\pi}{3}$), le vecteur $\overrightarrow{BC}$ a pour direction $120^\circ$ ($\frac{2\pi}{3}$). L'angle orienté $(\overrightarrow{AC}, \overrightarrow{BC})$ vaut donc $\frac{2\pi}{3} - \frac{\pi}{3} = \frac{\pi}{3}$ (sens direct).
Donc $(\overrightarrow{AB}, \overrightarrow{BC}) = \frac{\pi}{3} + \frac{\pi}{3} = \frac{2\pi}{3} \ [2\pi]$.
Finalement : $(\overrightarrow{BA}, \overrightarrow{BC}) = \frac{2\pi}{3} + \pi = \frac{5\pi}{3} \equiv -\frac{\pi}{3} \ [2\pi]$. \quad $\square$

\textbf{Morale :} La lecture directe sur la figure, en identifiant le sens de rotation, reste la méthode la plus sûre et la plus rapide. Le calcul algébrique est un excellent moyen de vérification.
\end{exoresolu}

\begin{exemplebox}[Carré et hexagone : les valeurs à reconnaître instantanément]
\textbf{Dans un carré $ABCD$ direct} (sommets pris dans le sens direct, par exemple $A$ bas-gauche, $B$ bas-droite, $C$ haut-droite, $D$ haut-gauche) :
\[ (\overrightarrow{AB}, \overrightarrow{AD}) = \frac{\pi}{2} \ [2\pi], \qquad (\overrightarrow{AB}, \overrightarrow{AC}) = \frac{\pi}{4} \ [2\pi] \quad (\text{diagonale}), \qquad (\overrightarrow{AB}, \overrightarrow{CD}) = \pi \ [2\pi]. \]
\textbf{Dans un hexagone régulier direct $ABCDEF$ de centre $O$} : chaque angle au centre vaut $\dfrac{\pi}{3}$ :
\[ (\overrightarrow{OA}, \overrightarrow{OB}) = \frac{\pi}{3} \ [2\pi], \qquad (\overrightarrow{OA}, \overrightarrow{OD}) = \pi \ [2\pi], \qquad (\overrightarrow{AB}, \overrightarrow{AF}) = \frac{2\pi}{3} \ [2\pi]. \]
Dans tous les cas, si la figure est \emph{indirecte} (sens horaire), toutes les mesures changent de signe.
\end{exemplebox}

\begin{attention}[Direct ou indirect : toujours vérifier l'énoncé]
Quand l'énoncé dit « $ABC$ est un triangle équilatéral » \emph{sans préciser le sens}, les deux configurations sont possibles : $(\overrightarrow{AB}, \overrightarrow{AC}) = \dfrac{\pi}{3}$ ou $-\dfrac{\pi}{3}$ selon la position de $C$. Au Probatoire, l'énoncé précise en général « triangle équilatéral direct » ou fournit une figure : appuie-toi dessus et annonce clairement le sens de rotation dans ta rédaction.
\end{attention}

\begin{exercice}[Pour t'entraîner]
\begin{enumerate}
  \item Déterminer la mesure principale de chacun des angles : $\dfrac{17\pi}{6}$ ; $-\dfrac{11\pi}{4}$ ; $\dfrac{25\pi}{3}$.
  \item $ABC$ est un triangle rectangle isocèle en $A$, direct. Déterminer $(\overrightarrow{AB}, \overrightarrow{AC})$, $(\overrightarrow{BA}, \overrightarrow{BC})$ et $(\overrightarrow{CA}, \overrightarrow{CB})$.
  \item Soient $\vec{u}$ et $\vec{v}$ non nuls tels que $(\vec{u}, \vec{v}) = \dfrac{2\pi}{5} \ [2\pi]$. Déterminer $(\vec{v}, \vec{u})$, $(-\vec{u}, \vec{v})$ et $(3\vec{u}, -2\vec{v})$.
\end{enumerate}
\end{exercice}

\begin{corrige}[Exercice : Rappels sur les angles orientés]
\textbf{1.} $\dfrac{17\pi}{6} = \dfrac{12\pi}{6} + \dfrac{5\pi}{6} = 2\pi + \dfrac{5\pi}{6}$ : mesure principale $\dfrac{5\pi}{6}$. \quad $-\dfrac{11\pi}{4} = -\dfrac{11\pi}{4} + \dfrac{8\pi}{4} = -\dfrac{3\pi}{4} \in \left]-\pi\,;\pi\right]$ : mesure principale $-\dfrac{3\pi}{4}$. \quad $\dfrac{25\pi}{3} = \dfrac{24\pi}{3} + \dfrac{\pi}{3} = 8\pi + \dfrac{\pi}{3}$ : mesure principale $\dfrac{\pi}{3}$.

\textbf{2.} On place le triangle : $A(0,0)$, $B(1,0)$, $C(0,1)$ (direct, rectangle en $A$).
\begin{itemize}
  \item $(\overrightarrow{AB}, \overrightarrow{AC}) = \dfrac{\pi}{2} \ [2\pi]$ (angle droit, sens direct).
  \item En $B$ : $\overrightarrow{BA}=(-1,0)$ (angle $\pi$), $\overrightarrow{BC}=(-1,1)$ (angle $\frac{3\pi}{4}$). Rotation de $\pi$ vers $\frac{3\pi}{4}$ : $-\frac{\pi}{4}$. Donc $(\overrightarrow{BA}, \overrightarrow{BC}) = -\dfrac{\pi}{4} \ [2\pi]$.
  \item En $C$ : $\overrightarrow{CA}=(0,-1)$ (angle $-\frac{\pi}{2}$), $\overrightarrow{CB}=(1,-1)$ (angle $-\frac{\pi}{4}$). Rotation de $-\frac{\pi}{2}$ vers $-\frac{\pi}{4}$ : $+\frac{\pi}{4}$. Donc $(\overrightarrow{CA}, \overrightarrow{CB}) = \dfrac{\pi}{4} \ [2\pi]$.
\end{itemize}

\textbf{3.} $(\vec{v}, \vec{u}) = -\dfrac{2\pi}{5} \ [2\pi]$ ; \quad $(-\vec{u}, \vec{v}) = \dfrac{2\pi}{5} + \pi = \dfrac{7\pi}{5} \ [2\pi]$, de mesure principale $-\dfrac{3\pi}{5}$ ; \quad $(3\vec{u}, -2\vec{v}) = (\vec{u}, -\vec{v}) = \dfrac{2\pi}{5} + \pi = \dfrac{7\pi}{5} \ [2\pi]$ (le facteur $3 > 0$ ne change rien, le facteur $-2$ ajoute $\pi$).
\end{corrige}

\begin{aretenir}[L'essentiel de la section]
\begin{itemize}
  \item Un angle orienté de vecteurs non nuls possède une infinité de mesures, toutes égales \textbf{modulo $2\pi$} ; la \cle{mesure principale} est l'unique mesure dans $]-\pi\,;\pi]$.
  \item La \cle{relation de Chasles} $(\vec{u}, \vec{v}) + (\vec{v}, \vec{w}) = (\vec{u}, \vec{w}) \ [2\pi]$ est l'outil de calcul central.
  \item Échanger les vecteurs change le signe ; changer un vecteur en son opposé ajoute $\pi$ ; multiplier par des réels strictement positifs ne change rien.
  \item Colinéarité $\Leftrightarrow (\vec{u}, \vec{v}) = 0 \ [\pi]$ ; orthogonalité $\Leftrightarrow (\vec{u}, \vec{v}) = \dfrac{\pi}{2} \ [\pi]$.
  \item L'angle de deux \textbf{droites} est défini \textbf{modulo $\pi$}.
\end{itemize}
Ces outils en main, nous pourrons dans la section suivante étudier l'angle inscrit et l'angle au centre, dernière étape avant les arcs capables qui résoudront le problème du technicien de Bafoussam et du capitaine de Kribi.
\end{aretenir}

\coursec{Angle inscrit, angle au centre et demi-tangente}

Dans la section précédente, nous avons rappelé la notion d'\cle{angle orienté} de vecteurs et sa \cle{mesure principale}. Nous allons maintenant mettre ces outils au service du cercle. Le cercle est une figure remarquable : les angles qui y vivent obéissent à des lois d'une grande élégance, connues depuis Euclide. Ces lois seront la clé de voûte des \cle{arcs capables} que nous étudierons dans les sections suivantes. L'idée directrice est simple : \emph{sur un cercle, tous les points d'un même arc « voient » une corde sous le même angle}.

\courssub{Angle inscrit et angle au centre}

\begin{definition}[Angle inscrit, angle au centre]
Soit $\mathcal{C}$ un cercle de centre $O$, et $A$, $B$, $M$ trois points distincts de $\mathcal{C}$.
\begin{itemize}
\item L'angle $\widehat{AMB}$ est appelé \cle{angle inscrit} dans le cercle $\mathcal{C}$. On dit qu'il \emph{intercepte} l'arc $\wideparen{AB}$ ne contenant pas $M$.
\item L'angle $\widehat{AOB}$ est appelé \cle{angle au centre} interceptant le même arc $\wideparen{AB}$.
\end{itemize}
\end{definition}

L'intuition est la suivante : imagine-toi dans un stade. La corde $[AB]$ représente les deux poteaux d'un but. Le point $O$ est au centre du rond central, et $M$ est un joueur sur la ligne du cercle. Le théorème qui suit affirme que l'angle sous lequel le joueur $M$ voit les poteaux est exactement la \emph{moitié} de l'angle sous lequel le spectateur placé en $O$ les voit.

\begin{propriete}[Théorème de l'angle au centre]
Soit $\mathcal{C}$ un cercle de centre $O$, et $A$, $B$, $M$ trois points distincts de $\mathcal{C}$. Alors, en angles orientés :
\[ \text{mes}(\overrightarrow{OA}, \overrightarrow{OB}) = 2\,\text{mes}(\overrightarrow{MA}, \overrightarrow{MB}) \quad [2\pi]. \]
En version géométrique (angles non orientés) : l'angle au centre est le double de l'angle inscrit qui intercepte le même arc :
\[ \text{mes}\,\widehat{AOB} = 2\,\text{mes}\,\widehat{AMB}. \]
\emph{Cette démonstration est exigible au Probatoire.}
\end{propriete}

\begin{experience}[Démonstration guidée du théorème de l'angle au centre]
On démontre d'abord le cas où $O$ est à l'intérieur de l'angle $\widehat{AMB}$, la droite $(MO)$ coupant le cercle en un point $N$ diamétralement opposé à $M$.
\begin{enumerate}
\item \textbf{Triangle isocèle.} $OA = OM$ (rayons), donc le triangle $OAM$ est isocèle en $O$. Posons $\alpha = \text{mes}\,\widehat{OMA} = \text{mes}\,\widehat{OAM}$.
\item \textbf{Angle extérieur.} L'angle $\widehat{AON}$ est extérieur au triangle $OAM$ en $O$ ; or la somme des angles d'un triangle vaut $\pi$, donc $\text{mes}\,\widehat{AON} = \pi - (\pi - 2\alpha) = 2\alpha$.
\item \textbf{Même travail de l'autre côté.} Dans le triangle isocèle $OBM$, en posant $\beta = \text{mes}\,\widehat{OMB}$, on obtient $\text{mes}\,\widehat{BON} = 2\beta$.
\item \textbf{Conclusion.} $\text{mes}\,\widehat{AOB} = \text{mes}\,\widehat{AON} + \text{mes}\,\widehat{NOB} = 2\alpha + 2\beta = 2(\alpha + \beta) = 2\,\text{mes}\,\widehat{AMB}$.
\end{enumerate}
Les autres positions de $O$ (extérieur à l'angle, ou sur un côté) se traitent de la même façon avec des différences au lieu de sommes. La version orientée s'en déduit, car l'égalité est vraie dans tous les cas de figure, modulo $2\pi$.
\end{experience}

\begin{popfigure}
\begin{center}
\begin{tikzpicture}[scale=1.6]
\coordinate (O) at (0,0);
\draw[popline, thick] (O) circle (1.5);
\fill[popink] (O) circle (0.03) node[below left] {$O$};
\coordinate (A) at (210:1.5);
\coordinate (B) at (-30:1.5);
\coordinate (M) at (90:1.5);
\fill[popink] (A) circle (0.035) node[below left] {$A$};
\fill[popink] (B) circle (0.035) node[below right] {$B$};
\fill[popink] (M) circle (0.035) node[above] {$M$};
\draw[popblue, thick] (O)--(A);
\draw[popblue, thick] (O)--(B);
\draw[poporange, thick] (M)--(A);
\draw[poporange, thick] (M)--(B);
\draw[popblue, thick] (210:0.45) arc (210:330:0.45);
\node[popblue] at (270:0.72) {$2\alpha$};
\draw[poporange, thick] ($(M)+(210:0.55)$) arc (210:330:0.55);
\node[poporange] at ($(M)+(270:0.82)$) {$\alpha$};
\draw[popmuted, dashed] (A)--(B);
\end{tikzpicture}
\end{center}
\legende{L'angle au centre $\widehat{AOB}$ (bleu) est le double de l'angle inscrit $\widehat{AMB}$ (orange) interceptant le même arc $\wideparen{AB}$.}
\end{popfigure}

Une conséquence immédiate et fondamentale : si deux points $M$ et $N$ sont sur le cercle, les angles inscrits $\widehat{AMB}$ et $\widehat{ANB}$ sont liés. En angles orientés, la relation de Chasles permet d'énoncer proprement le résultat sans se soucier des positions.

\begin{propriete}[Théorème de l'angle inscrit, version orientée]
Soit $A$, $B$, $M$, $N$ quatre points distincts d'un cercle $\mathcal{C}$. Alors :
\[ \text{mes}(\overrightarrow{MA}, \overrightarrow{MB}) = \text{mes}(\overrightarrow{NA}, \overrightarrow{NB}) \quad [\pi]. \]
Autrement dit, les angles orientés sous lesquels on voit la corde $[AB]$ depuis deux points du cercle sont égaux \emph{modulo $\pi$}.
\end{propriete}

\begin{experience}[Pourquoi modulo $\pi$ et pas modulo $2\pi$ ?]
Notons $O$ le centre du cercle. Par la relation de Chasles et le théorème de l'angle au centre :
\begin{align*}
\text{mes}(\overrightarrow{MA}, \overrightarrow{MB}) &= \tfrac{1}{2}\,\text{mes}(\overrightarrow{OA}, \overrightarrow{OB}) \quad [\pi], \\
\text{mes}(\overrightarrow{NA}, \overrightarrow{NB}) &= \tfrac{1}{2}\,\text{mes}(\overrightarrow{OA}, \overrightarrow{OB}) \quad [\pi].
\end{align*}
En effet, le théorème de l'angle au centre donne l'égalité modulo $2\pi$ pour le double ; en divisant par $2$, on ne conserve l'égalité que modulo $\pi$. D'où le résultat. Géométriquement : si $M$ et $N$ sont sur le \emph{même} arc, les angles géométriques sont égaux ; s'ils sont sur des arcs \emph{différents}, les angles géométriques sont \emph{supplémentaires} — et deux angles supplémentaires sont précisément égaux modulo $\pi$ en version orientée.
\end{experience}

\begin{propriete}[Caractérisation de la cocyclicité]
Quatre points distincts $A$, $B$, $M$, $N$ sont \cle{cocycliques} (ou alignés) si et seulement si :
\[ \text{mes}(\overrightarrow{MA}, \overrightarrow{MB}) = \text{mes}(\overrightarrow{NA}, \overrightarrow{NB}) \quad [\pi]. \]
C'est l'outil numéro un pour prouver que quatre points sont sur un même cercle au Probatoire.
\end{propriete}

\begin{attention}[Le piège des arcs différents]
En version \emph{non orientée}, le théorème de l'angle inscrit ne dit PAS que $\widehat{AMB} = \widehat{ANB}$ pour tous les points du cercle ! Si $M$ et $N$ sont sur des arcs déterminés par $A$ et $B$ \emph{différents}, alors :
\[ \text{mes}\,\widehat{AMB} + \text{mes}\,\widehat{ANB} = 180^\circ. \]
Exemple : si $\text{mes}\,\widehat{AMB} = 70^\circ$ et que $N$ est sur l'autre arc, alors $\text{mes}\,\widehat{ANB} = 110^\circ$, et non $70^\circ$. Avant d'écrire une égalité d'angles géométriques, vérifie toujours que les points sont sur le même arc — ou travaille directement avec les angles orientés modulo $\pi$, qui dispensent de cette vérification.
\end{attention}

\courssub{Angle formé par une corde et une demi-tangente}

Considérons maintenant une situation limite : que devient l'angle inscrit $\widehat{AMB}$ lorsque le point $M$ glisse le long du cercle et vient se confondre avec $A$ ? La droite $(MA)$ devient alors la \cle{tangente} au cercle en $A$. L'angle inscrit « survit » sous la forme de l'angle entre la tangente et la corde $[AB]$.

\begin{definition}[Angle d'une corde et d'une demi-tangente]
Soit $\mathcal{C}$ un cercle, $A$ et $B$ deux points de $\mathcal{C}$, et $[AT)$ une demi-tangente à $\mathcal{C}$ en $A$. L'angle $\widehat{TAB}$ (ou l'angle orienté $(\overrightarrow{AT}, \overrightarrow{AB})$) est appelé \cle{angle formé par la corde $[AB]$ et la demi-tangente $[AT)$}.
\end{definition}

\begin{propriete}[Théorème de la tangente et de la corde]
Soit $\mathcal{C}$ un cercle, $A$, $B$, $M$ trois points de $\mathcal{C}$, et $[AT)$ la demi-tangente en $A$ située du côté de l'arc $\wideparen{AB}$ contenant $M$. Alors :
\[ \text{mes}(\overrightarrow{AT}, \overrightarrow{AB}) = \text{mes}(\overrightarrow{MA}, \overrightarrow{MB}) \quad [\pi]. \]
En version géométrique : l'angle entre la demi-tangente et la corde est égal à l'angle inscrit interceptant cette corde du côté opposé.
\emph{Cette démonstration est exigible au Probatoire.}
\end{propriete}

\begin{experience}[Démonstration guidée du théorème de la tangente]
Soit $O$ le centre du cercle. On raisonne en angles géométriques, $[AT)$ étant du côté de l'arc contenant $M$.
\begin{enumerate}
\item \textbf{Tangente et rayon.} La tangente en $A$ est perpendiculaire au rayon $[OA]$, donc $\text{mes}\,\widehat{OAT} = 90^\circ$.
\item \textbf{Triangle isocèle $OAB$.} $OA = OB$, donc $\text{mes}\,\widehat{OAB} = \text{mes}\,\widehat{OBA}$. La somme des angles du triangle donne :
\[ \text{mes}\,\widehat{OAB} = \frac{180^\circ - \text{mes}\,\widehat{AOB}}{2} = 90^\circ - \frac{1}{2}\,\text{mes}\,\widehat{AOB}. \]
\item \textbf{Angle entre tangente et corde.} 
\[ \text{mes}\,\widehat{TAB} = \text{mes}\,\widehat{OAT} - \text{mes}\,\widehat{OAB} = 90^\circ - \left(90^\circ - \frac{1}{2}\,\text{mes}\,\widehat{AOB}\right) = \frac{1}{2}\,\text{mes}\,\widehat{AOB}. \]
\item \textbf{Conclusion via l'angle au centre.} Par le théorème de l'angle au centre, $\frac{1}{2}\,\text{mes}\,\widehat{AOB} = \text{mes}\,\widehat{AMB}$. Donc $\text{mes}\,\widehat{TAB} = \text{mes}\,\widehat{AMB}$. La version orientée modulo $\pi$ s'obtient en reprenant le calcul avec les angles orientés.
\end{enumerate}
\end{experience}

\begin{popfigure}
\begin{center}
\begin{tikzpicture}[scale=1.6]
\coordinate (O) at (0,0);
\draw[popline, thick] (O) circle (1.5);
\fill[popink] (O) circle (0.03) node[below right] {$O$};
\coordinate (A) at (200:1.5);
\coordinate (B) at (-20:1.5);
\coordinate (M) at (80:1.5);
\fill[popink] (A) circle (0.035) node[below left] {$A$};
\fill[popink] (B) circle (0.035) node[below right] {$B$};
\fill[popink] (M) circle (0.035) node[above] {$M$};
\draw[popdark, thick] (A)--(B);
\draw[poporange, thick] (M)--(A);
\draw[poporange, thick] (M)--(B);
\draw[popgreen, thick, -{Stealth}] ($(A)+(110:2.2)$) -- ($(A)+(-70:1.2)$);
\coordinate (T) at ($(A)+(110:1.6)$);
\fill[popink] (T) circle (0.02) node[above left] {$T$};
\draw[popgreen, thick] ($(A)+(110:0.5)$) arc (110:-20:0.5);
\node[popgreen] at ($(A)+(45:0.75)$) {$\alpha$};
\draw[poporange, thick] ($(M)+(200:0.5)$) arc (200:320:0.5);
\node[poporange] at ($(M)+(260:0.75)$) {$\alpha$};
\draw[popmuted, dashed] (O)--(A);
\end{tikzpicture}
\end{center}
\legende{L'angle entre la demi-tangente $[AT)$ et la corde $[AB]$ (vert) est égal à l'angle inscrit $\widehat{AMB}$ (orange).}
\end{popfigure}

\begin{methode}[Construire un cercle connaissant une corde et l'angle sous lequel elle est vue]
On veut construire le cercle passant par $A$ et $B$ tel que, pour tout point $M$ d'un arc, $\text{mes}\,\widehat{AMB} = \alpha$ donné.
\begin{enumerate}
\item Tracer la demi-droite $[AT)$ telle que $\text{mes}\,\widehat{TAB} = \alpha$ (report d'angle au rapporteur).
\item Le centre $O$ est sur la perpendiculaire en $A$ à la droite $(AT)$ (propriété tangente--rayon).
\item Le centre $O$ est aussi sur la médiatrice de $[AB]$ (car $OA = OB$).
\item L'intersection de ces deux droites donne $O$ ; tracer le cercle de centre $O$ passant par $A$ (et $B$).
\end{enumerate}
Cette construction sera essentielle pour les arcs capables de la section suivante.
\end{methode}

\begin{exoresolu}[Angle au centre et angle tangentiel]
Soit $\mathcal{C}$ un cercle de centre $O$, et $A$, $B$, $M$ trois points de $\mathcal{C}$ tels que $\text{mes}\,\widehat{AMB} = 35^\circ$.
\begin{enumerate}
\item Calculer $\text{mes}\,\widehat{AOB}$.
\item Calculer l'angle entre la tangente en $A$ et la corde $[AB]$.
\item $N$ est un point de $\mathcal{C}$ situé sur l'arc $\wideparen{AB}$ ne contenant pas $M$. Calculer $\text{mes}\,\widehat{ANB}$.
\end{enumerate}
\tcblower
\textbf{Solution détaillée.}
\begin{enumerate}
\item Par le théorème de l'angle au centre : $\text{mes}\,\widehat{AOB} = 2 \times \text{mes}\,\widehat{AMB} = 2 \times 35^\circ = 70^\circ$.
\item Par le théorème de la tangente et de la corde, l'angle entre la tangente en $A$ et la corde $[AB]$ est égal à l'angle inscrit interceptant $[AB]$ : il vaut $35^\circ$ (du côté de l'arc contenant $M$) ou $180^\circ - 35^\circ = 145^\circ$ de l'autre côté.
\item $M$ et $N$ sont sur des arcs différents : les angles inscrits sont supplémentaires. Donc $\text{mes}\,\widehat{ANB} = 180^\circ - 35^\circ = 145^\circ$. 
\end{enumerate}
\textbf{Précision sur les angles orientés.} En version orientée, si l'on choisit $\text{mes}(\overrightarrow{MA}, \overrightarrow{MB}) = 35^\circ$ (mesure principale), alors $\text{mes}(\overrightarrow{NA}, \overrightarrow{NB}) = 35^\circ + 180^\circ = 215^\circ \equiv -145^\circ \ [2\pi]$. On a bien $215^\circ \equiv 35^\circ \ [\pi]$ car $215^\circ - 35^\circ = 180^\circ$. La mesure géométrique $145^\circ$ n'est \emph{pas} congrue à $35^\circ$ modulo $\pi$, mais l'angle \emph{orienté} associé ($215^\circ$ ou $-145^\circ$) l'est. C'est exactement le piège signalé plus haut : ne jamais confondre mesure géométrique (dans $[0,\pi]$) et mesure orientée (modulo $2\pi$) !
\end{exoresolu}

\begin{exercice}[Exercice 1]
Sur un cercle de centre $O$, on donne $\text{mes}\,\widehat{AOB} = 96^\circ$. Détermine toutes les valeurs possibles de $\text{mes}\,\widehat{AMB}$ lorsque $M$ décrit le cercle (privé de $A$ et $B$), puis l'angle entre la tangente en $B$ et la corde $[BA]$.
\end{exercice}

\begin{corrige}[Exercice 1]
Si $M$ est sur l'arc majeur (celui qui ne « voit » pas l'angle au centre de $96^\circ$), $\text{mes}\,\widehat{AMB} = \frac{96^\circ}{2} = 48^\circ$. Si $M$ est sur l'arc mineur, l'angle au centre intercepté vaut $360^\circ - 96^\circ = 264^\circ$, d'où $\text{mes}\,\widehat{AMB} = 132^\circ$ (on vérifie : $48^\circ + 132^\circ = 180^\circ$). L'angle entre la tangente en $B$ et la corde $[BA]$ vaut $48^\circ$ du côté de l'arc majeur, $132^\circ$ de l'autre.
\end{corrige}

\begin{savaistu}[Un théorème vieux de 2300 ans]
Le théorème de l'angle inscrit figure déjà dans les \emph{Éléments} d'\textbf{Euclide} (III\textsuperscript{e} siècle av. J.-C.), au livre III, proposition 20 : « Dans un cercle, l'angle au centre est double de l'angle à la circonférence. » Ce résultat a servi pendant des siècles à la navigation : un navigateur mesurant l'angle sous lequel il voit deux amer (phares, collines) peut se placer sur un arc de cercle sur sa carte — c'est exactement la technique de l'\emph{arc capable}, fondement du point au sextant. Aujourd'hui encore, ces propriétés sont utilisées en géodésie et en cartographie, par exemple pour les relevés terrain des ingénieurs du Cameroun.
\end{savaistu}

\begin{aretenir}[L'essentiel de la section]
\begin{itemize}
\item \textbf{Angle au centre} : $\text{mes}(\overrightarrow{OA}, \overrightarrow{OB}) = 2\,\text{mes}(\overrightarrow{MA}, \overrightarrow{MB}) \ [2\pi]$.
\item \textbf{Angles inscrits} : pour $A, B, M, N$ cocycliques, $\text{mes}(\overrightarrow{MA}, \overrightarrow{MB}) = \text{mes}(\overrightarrow{NA}, \overrightarrow{NB}) \ [\pi]$ — et c'est une \emph{caractérisation} de la cocyclicité.
\item \textbf{Tangente et corde} : $\text{mes}(\overrightarrow{AT}, \overrightarrow{AB}) = \text{mes}(\overrightarrow{MA}, \overrightarrow{MB}) \ [\pi]$.
\item En version géométrique, attention au \emph{supplémentaire} quand les points sont sur des arcs différents.
\end{itemize}
\end{aretenir}

\coursec{Arc capable : ensemble des points M tels que mes(AMB) = x}

Dans la section précédente, nous avons établi un résultat fondamental : tous les angles inscrits qui interceptent le même arc d'un même cercle ont la même mesure, égale à la moitié de l'angle au centre. Autrement dit, si deux points $M$ et $M'$ sont sur un même arc de cercle d'extrémités $A$ et $B$, alors $\text{mes}(\widehat{AMB}) = \text{mes}(\widehat{AM'B})$. Nous allons maintenant \emph{renverser la perspective} : au lieu de partir d'un cercle et d'observer que les angles sont égaux, nous partons d'un angle imposé $x$ et nous cherchons \textbf{tous} les points $M$ du plan d'où le segment $[AB]$ est « vu » sous cet angle. C'est le problème de l'\cle{arc capable}, un outil extraordinairement puissant pour les lieux géométriques au Probatoire.

\courssub{Le problème : voir un segment sous un angle donné}

\textbf{Une situation concrète.} Deux poteaux $A$ et $B$ bordent le terrain de football du lycée. Un photographe veut se placer de telle sorte que, sur son cliché, les deux poteaux apparaissent sous un angle de $60^{\circ}$ exactement. Où peut-il se poster ? Un seul point ? Une courbe ? C'est exactement la question que résout la théorie des arcs capables.

\begin{definition}[Angle sous lequel on voit un segment]
Soient $A$ et $B$ deux points distincts du plan et $M$ un point n'appartenant pas à la droite $(AB)$. On dit que $M$ \cle{voit le segment $[AB]$ sous l'angle} $x$ lorsque
\[
\text{mes}(\widehat{AMB}) = x.
\]
Il s'agit ici d'un \textbf{angle géométrique} (non orienté) : $x \in [0\,;\pi]$.
\end{definition}

Le problème à résoudre est donc le suivant : \emph{étant donnés deux points $A$ et $B$ et un réel $x \in [0\,;\pi]$, déterminer l'ensemble $\mathcal{E}$ des points $M$ du plan tels que $\text{mes}(\widehat{AMB}) = x$.}

Nous allons traiter d'abord les cas limites $x = 0$, $x = \pi$ et le cas remarquable $x = \dfrac{\pi}{2}$, avant le cas général.

\courssub{Les cas particuliers : $x = 0$, $x = \pi$, $x = \dfrac{\pi}{2}$}

\textbf{Cas $x = 0$.} Dire que $\text{mes}(\widehat{AMB}) = 0$ signifie que les demi-droites $[MA)$ et $[MB)$ sont confondues : les points $M$, $A$, $B$ sont alignés et $M$ n'est pas « entre » $A$ et $B$ (sinon l'angle serait plat). Réciproquement, tout point de la droite $(AB)$ situé hors du segment $[AB]$ voit $[AB]$ sous un angle nul.

\begin{propriete}[Lieu pour $x = 0$]
L'ensemble des points $M$ tels que $\text{mes}(\widehat{AMB}) = 0$ est la \textbf{droite $(AB)$ privée du segment $[AB]$} (c'est-à-dire les deux demi-droites ouvertes d'extrémités $A$ et $B$ ne contenant pas l'autre point).
\end{propriete}

\textbf{Cas $x = \pi$.} Dire que $\text{mes}(\widehat{AMB}) = \pi$ signifie que $[MA)$ et $[MB)$ sont deux demi-droites opposées : $M$ est aligné avec $A$ et $B$ et situé \emph{entre} eux.

\begin{propriete}[Lieu pour $x = \pi$]
L'ensemble des points $M$ tels que $\text{mes}(\widehat{AMB}) = \pi$ est le \textbf{segment $[AB]$ privé de ses extrémités $A$ et $B$}.
\end{propriete}

\textbf{Cas $x = \dfrac{\pi}{2}$.} Dire que $\text{mes}(\widehat{AMB}) = \dfrac{\pi}{2}$ signifie que le triangle $AMB$ est rectangle en $M$. Le théorème de Thalès (et sa réciproque) nous donne immédiatement la réponse.

\begin{propriete}[Lieu pour $x = \dfrac{\pi}{2}$ : théorème de Thalès]
L'ensemble des points $M$ tels que $\text{mes}(\widehat{AMB}) = \dfrac{\pi}{2}$ est le \textbf{cercle de diamètre $[AB]$, privé des points $A$ et $B$}.
\end{propriete}

Ce cas est un excellent indice pour la suite : le lieu est un cercle (presque entier). Pour un angle $x$ quelconque, nous devons nous attendre à des \emph{morceaux} de cercles.

\courssub{Cas général : $0 < x < \pi$}

\textbf{Intuition.} Fixons $x = 60^{\circ}$ et imaginons un point $M$ au-dessus de $[AB]$ tel que $\text{mes}(\widehat{AMB}) = 60^{\circ}$. Le point $M$ n'est pas sur le cercle de diamètre $[AB]$ (sinon l'angle vaudrait $90^{\circ}$) ; il est sur un cercle \emph{plus grand}, passant par $A$ et $B$. Par le théorème de l'angle inscrit, \emph{tous} les points du même arc de ce cercle (celui qui ne contient pas le centre $O$ relativement à la corde $[AB]$, c'est-à-dire l'arc majeur) voient $[AB]$ sous le même angle $60^{\circ}$. Et par symétrie par rapport à la droite $(AB)$, il existe un second arc, en dessous, qui convient aussi.

\begin{definition}[Arc capable]
Soient $A$ et $B$ deux points distincts et $x \in \left]0\,;\pi\right[$. On appelle \cle{arc capable} du segment $[AB]$ pour l'angle $x$ tout arc de cercle d'extrémités $A$ et $B$ dont tous les points $M$ (distincts de $A$ et $B$) vérifient $\text{mes}(\widehat{AMB}) = x$.
\end{definition}

\begin{propriete}[Théorème de l'arc capable]
Soient $A$ et $B$ deux points distincts et $x \in \left]0\,;\pi\right[$. L'ensemble des points $M$ du plan tels que
\[
\text{mes}(\widehat{AMB}) = x
\]
est la \textbf{réunion de deux arcs de cercle d'extrémités $A$ et $B$, symétriques par rapport à la droite $(AB)$, privés des points $A$ et $B$}. Ces deux arcs sont appelés \cle{arcs capables} du segment $[AB]$ pour l'angle $x$. Leurs centres $O$ et $O'$ sont situés sur la médiatrice de $[AB]$, et le rayon commun est
\[
R = \frac{AB}{2\sin x}.
\]
Ce résultat est admis au Probatoire, mais sa justification repose entièrement sur le théorème de l'angle inscrit et le théorème de la demi-tangente, que nous détaillons ci-dessous.
\end{propriete}

\begin{experience}[Pourquoi le lieu est-il un arc de cercle ?]
Raisonnons en deux temps, dans le demi-plan situé au-dessus de $(AB)$.
\begin{enumerate}
\item \textbf{Tout point d'un tel arc convient.} Soit $\mathcal{C}$ un cercle passant par $A$ et $B$, et soient $M$ et $M'$ deux points du même arc d'extrémités $A$ et $B$. Les angles inscrits $\widehat{AMB}$ et $\widehat{AM'B}$ interceptent le même arc ; par le théorème de l'angle inscrit, $\text{mes}(\widehat{AMB}) = \text{mes}(\widehat{AM'B})$. Donc si un point de l'arc voit $[AB]$ sous l'angle $x$, tous les points de cet arc le voient sous l'angle $x$.
\item \textbf{Il n'y en a pas d'autres.} Soit $M$ un point tel que $\text{mes}(\widehat{AMB}) = x$, situé au-dessus de $(AB)$. Le cercle circonscrit au triangle $AMB$ passe par $A$, $B$, $M$. Son centre $O$ est sur la médiatrice de $[AB]$, et l'angle au centre vérifie $\text{mes}(\widehat{AOB}) = 2x$ (angle au centre interceptant le même arc que l'angle inscrit en $M$). Or, $A$, $B$ et la valeur $2x$ déterminent entièrement $O$ sur la médiatrice (dans le triangle $OIA$ rectangle en $I$, milieu de $[AB]$, on a $\tan x = \dfrac{IA}{OI}$, soit $OI = \dfrac{AB}{2\tan x}$). Il n'y a donc qu'\textbf{un seul} tel cercle au-dessus de $(AB)$ : $M$ appartient nécessairement à l'arc décrit au point 1.
\end{enumerate}
Le même raisonnement par symétrie donne l'arc situé en dessous de $(AB)$. Enfin, la formule du rayon s'obtient dans le triangle $OIA$ : $\sin x = \dfrac{IA}{OA} = \dfrac{AB/2}{R}$, d'où $R = \dfrac{AB}{2\sin x}$.
\end{experience}

\begin{popfigure}
\begin{center}
\begin{tikzpicture}[scale=0.85]
% Coordonnées
\coordinate (A) at (-2,0);
\coordinate (B) at (2,0);
\coordinate (I) at (0,0);
\coordinate (O) at (0,1.155); % OI = R cos 60 = (4/sqrt3)*0.5 = 2/sqrt3 ~ 1.155
\coordinate (Op) at (0,-1.155);
% Rayon
\def\R{2.309} % 4/sqrt(3) ~ 2.309

% Médiatrice
\draw[dashed,popmuted] (0,-3.5) -- (0,3.5);
% Segment AB
\draw[very thick,popdark] (A) -- (B);

% Arc capable supérieur (Majeur, centre O, angles 330 à 570 depuis O)
% On trace de B (angle 330) vers A (angle 210+360=570) en sens trigonométrique (arc majeur haut)
\draw[very thick,popblue] (O) ++(330:\R) arc (330:570:\R);
% Arc capable inférieur (Majeur, centre Op, angles 150 à 390 depuis Op)
% B est à angle 30, A à 150. Arc majeur bas : de A (150) vers B (30+360=390) en sens trigonométrique
\draw[very thick,poporange] (Op) ++(150:\R) arc (150:390:\R);

% Centres
\fill[popblue] (O) circle (1.5pt) node[above right]{$O$};
\fill[poporange] (Op) circle (1.5pt) node[below right]{$O'$};
% Rayons OA, OB (pour visualisation)
\draw[popblue!60] (O) -- (A);
\draw[popblue!60] (O) -- (B);

% Points A, B, I
\fill[popdark] (A) circle (1.8pt) node[below left]{$A$};
\fill[popdark] (B) circle (1.8pt) node[below right]{$B$};
\fill[popmuted] (I) circle (1.2pt) node[below]{$I$};

% Point M sur l'arc supérieur (ex: angle 90° depuis O -> sommet du cercle)
\coordinate (M) at (0, 3.464); % O_y + R = 1.155 + 2.309
\fill[popgreen] (M) circle (1.8pt) node[above]{$M$};
\draw[popgreen] (M) -- (A);
\draw[popgreen] (M) -- (B);
% Marque d'angle 60° en M
% Vecteurs MA = (-2, -3.464) angle ~ -120° (240°), MB = (2, -3.464) angle ~ -60° (300°)
% Angle entre eux = 60°. On marque entre 240 et 300 (ou -120 et -60)
\draw[popgreen] ($(M)+(240:0.6)$) arc (240:300:0.6);
\node[popgreen] at ($(M)+(270:0.9)$) {$60^{\circ}$};
\end{tikzpicture}
\end{center}
\legende{Les deux arcs capables du segment $[AB]$ pour $x = 60^{\circ}$ : symétriques par rapport à $(AB)$, de centres $O$ et $O'$ sur la médiatrice de $[AB]$. Tout point $M$ de l'arc bleu (arc majeur du cercle de centre $O$) voit $[AB]$ sous $60^{\circ}$.}
\end{popfigure}

\courssub{Construction de l'arc capable}

Comment tracer effectivement l'arc capable pour un angle $x$ donné ? L'idée-clé est le \textbf{théorème de la demi-tangente} : si $[At)$ est tangente en $A$ au cercle cherché, alors l'angle entre la tangente et la corde $[AB]$ est égal à l'angle inscrit qui intercepte cette corde, c'est-à-dire $x$. Il suffit donc de construire la tangente... pour en déduire le centre !

\begin{methode}[{Construire l'arc capable de $[AB]$ pour l'angle $x$}]
\begin{enumerate}
\item \textbf{Demi-tangente.} Tracer en $A$ la demi-droite $[At)$ telle que l'angle entre $[AB]$ et $[At)$ mesure $x$, du côté opposé à celui où l'on veut placer l'arc.
\item \textbf{Perpendiculaire.} Élever en $A$ la perpendiculaire à $[At)$ : par le théorème de la demi-tangente, $[At)$ sera tangente au cercle, donc cette perpendiculaire \textbf{passe par le centre} $O$.
\item \textbf{Médiatrice.} Tracer la médiatrice de $[AB]$ : le centre $O$ lui appartient aussi (car $OA = OB$). L'intersection des deux droites est le centre $O$.
\item \textbf{Arc.} Tracer au compas, de centre $O$ et de rayon $OA$, l'arc d'extrémités $A$ et $B$ situé du même côté que $[At)$, \emph{sans} les points $A$ et $B$. Construire par symétrie le second arc si le problème demande le lieu complet.
\end{enumerate}
\end{methode}

\begin{popfigure}
\begin{center}
\begin{tikzpicture}[scale=0.9]
\coordinate (A) at (-2,0);
\coordinate (B) at (2,0);
\coordinate (I) at (0,0);
\coordinate (O) at (0,1.155);
% étape 1 : demi-droite [At) à 60° de [AB)
\draw[very thick,poporange,->] (A) -- ++(60:3.6) node[above left]{$t$};
\draw[poporange] ($(A)+(0:0.9)$) arc[start angle=0,end angle=60,radius=0.9];
\node[poporange] at ($(A)+(30:1.35)$) {$x$};
% étape 2 : perpendiculaire en A à [At)
\draw[popblue,thick] ($(A)+(150:1.2)$) -- ($(A)+(-30:4.6)$);
\node[popblue] at ($(A)+(-30:4.9)$) {$\perp [At)$};
% étape 3 : médiatrice
\draw[dashed,popmuted] (0,-1.2) -- (0,3.2);
\node[popmuted] at (0.15,3.05) {médiatrice};
% segment AB
\draw[very thick,popdark] (A) -- (B);
\fill[popdark] (A) circle (1.8pt) node[below left]{$A$};
\fill[popdark] (B) circle (1.8pt) node[below right]{$B$};
\fill[popmuted] (I) circle (1.2pt) node[below]{$I$};
% centre O
\fill[popblue] (O) circle (1.8pt) node[above right]{$O$};
% étape 4 : arc
\draw[very thick,popgreen] (A) arc[start angle=150,end angle=30,radius=2.309];
% codage angle droit
\draw[popblue] ($(A)+(60:0.35)$) -- ($(A)+(60:0.35)+(-30:0.35)$) -- ($(A)+(-30:0.35)$);
\end{tikzpicture}
\end{center}
\legende{Construction de l'arc capable : (1) demi-droite $[At)$ avec $\text{mes}(\widehat{tAB}) = x$ ; (2) perpendiculaire en $A$ à $[At)$ ; (3) médiatrice de $[AB]$ ; leur intersection est le centre $O$ ; (4) l'arc de centre $O$ passant par $A$ et $B$.}
\end{popfigure}

\textbf{Justification par le théorème de la demi-tangente.} Soit $\mathcal{C}$ le cercle de centre $O$ construit ci-dessus. Comme $OA \perp [At)$, la demi-droite $[At)$ est tangente en $A$ à $\mathcal{C}$. Pour tout point $M$ de l'arc tracé (ne contenant pas le point de tangence « côté $t$ »), le théorème de la demi-tangente affirme que l'angle inscrit $\widehat{AMB}$ et l'angle $\widehat{tAB}$ interceptent la même corde $[AB]$, donc :
\[
\text{mes}(\widehat{AMB}) = \text{mes}(\widehat{tAB}) = x.
\]
La construction fournit bien un arc capable. Réciproquement, tout arc capable admet en $A$ une tangente faisant l'angle $x$ avec $[AB]$, donc son centre est nécessairement l'intersection construite : la construction donne \emph{le} lieu, tout entier.

\begin{exemplebox}[{Lecture inverse : tout point de l'arc voit $[AB]$ sous l'angle $x$}]
Sur la figure ci-dessus avec $x = 60^{\circ}$, choisissons un point $M$ quelconque de l'arc vert, par exemple le point le plus haut de l'arc (sur la médiatrice). Le triangle $AMB$ est alors isocèle en $M$, et $\text{mes}(\widehat{AMB}) = 60^{\circ}$ : c'est même un triangle \textbf{équilatéral}. Vérifions avec la formule du rayon : $R = \dfrac{AB}{2\sin 60^{\circ}} = \dfrac{AB}{\sqrt{3}}$, et la distance $OM = R$... Tout point de l'arc, pas seulement celui-là, voit $[AB]$ sous exactement $60^{\circ}$ : c'est toute la force du théorème de l'angle inscrit.
\end{exemplebox}

\begin{exoresolu}[Arc capable pour $AB = 6$ cm et $x = 45^{\circ}$]
On donne deux points $A$ et $B$ tels que $AB = 6$ cm. Déterminer et construire l'ensemble des points $M$ tels que $\text{mes}(\widehat{AMB}) = 45^{\circ}$, et calculer le rayon des arcs capables.
\tcblower
\textbf{Nature du lieu.} Comme $0 < 45^{\circ} < 180^{\circ}$, le théorème de l'arc capable s'applique : le lieu est la réunion de deux arcs de cercle d'extrémités $A$ et $B$, symétriques par rapport à $(AB)$, privés de $A$ et $B$.

\textbf{Calcul du rayon.} D'après la formule établie dans le triangle formé par le centre $O$, le milieu $I$ de $[AB]$ et le point $A$ :
\[
R = \frac{AB}{2\sin x} = \frac{6}{2\sin 45^{\circ}} = \frac{6}{2 \times \frac{\sqrt{2}}{2}} = \frac{6}{\sqrt{2}} = 3\sqrt{2} \text{ cm} \approx 4,24 \text{ cm}.
\]

\textbf{Construction.} 
\begin{enumerate}
\item On trace en $A$ la demi-droite $[At)$ telle que $\text{mes}(\widehat{tAB}) = 45^{\circ}$ (diagonale d'un carré construit sur une perpendiculaire, ou rapporteur).
\item On élève en $A$ la perpendiculaire à $[At)$.
\item On trace la médiatrice de $[AB]$ ; elle coupe cette perpendiculaire en $O$.
\item On trace l'arc de centre $O$, de rayon $OA = 3\sqrt{2}$ cm, d'extrémités $A$ et $B$, puis son symétrique par rapport à $(AB)$.
\end{enumerate}
\textbf{Vérification.} L'angle au centre vaut $\text{mes}(\widehat{AOB}) = 2 \times 45^{\circ} = 90^{\circ}$ : le triangle $OAB$ est rectangle isocèle en $O$, d'hypoténuse $AB = R\sqrt{2} = 3\sqrt{2}\times\sqrt{2} = 6$ cm. C'est cohérent.
\end{exoresolu}

\begin{attention}[Les trois erreurs classiques sur les arcs capables]
\begin{enumerate}
\item \textbf{Oublier d'exclure $A$ et $B$.} En $M = A$ ou $M = B$, l'angle $\widehat{AMB}$ n'a \emph{aucun sens} (on ne peut pas définir les demi-droites $[MA)$ et $[MB)$). Le lieu est l'arc \emph{privé de ses extrémités}. Au Probatoire, écris toujours : « l'arc d'extrémités $A$ et $B$, privé des points $A$ et $B$ ».
\item \textbf{Oublier le second arc.} Le lieu complet pour un angle géométrique $x$ comporte \textbf{deux} arcs, symétriques par rapport à $(AB)$. N'en donner qu'un, c'est perdre la moitié des points... et la moitié des points au barème ! (Seule l'étude des arcs capables \emph{orientés}, section suivante, distinguera les deux arcs.)
\item \textbf{Confondre avec le cercle entier.} Pour $x = 90^{\circ}$ seulement, le lieu est (presque) un cercle entier : les deux arcs se rejoignent en un cercle de diamètre $[AB]$. Pour $x \neq 90^{\circ}$, ce sont deux arcs de cercles \emph{différents} (deux cercles distincts, symétriques).
\end{enumerate}
\end{attention}

\begin{aretenir}[L'essentiel sur l'arc capable]
\begin{itemize}
\item Pour $0 < x < \pi$, l'ensemble des points $M$ tels que $\text{mes}(\widehat{AMB}) = x$ est la réunion de \textbf{deux arcs de cercle} d'extrémités $A$ et $B$, symétriques par rapport à $(AB)$, \textbf{privés de $A$ et $B$}.
\item Cas limites : $x = 0$ donne la droite $(AB)$ privée de $[AB]$ ; $x = \pi$ donne le segment $[AB]$ privé de $A$ et $B$ ; $x = \dfrac{\pi}{2}$ donne le cercle de diamètre $[AB]$ privé de $A$ et $B$ (Thalès).
\item Rayon de l'arc capable : $R = \dfrac{AB}{2\sin x}$ ; angle au centre : $\text{mes}(\widehat{AOB}) = 2x$.
\item Construction en 4 étapes : angle $x$ en $A$ $\rightarrow$ perpendiculaire en $A$ $\rightarrow$ médiatrice de $[AB]$ $\rightarrow$ arc de centre $O$.
\end{itemize}
\end{aretenir}

\begin{savaistu}[L'arc capable dans la vie courante]
Les arcs capables interviennent partout où l'on « vise » un objet sous un angle imposé. En navigation maritime, un capitaine qui connaît l'angle sous lequel il voit deux balises côtières sait qu'il se trouve sur un arc capable : en croisant deux telles observations, il détermine sa position sans GPS — c'est la méthode dite des « arcs capables » encore enseignée dans les écoles de marine. Au Cameroun, les techniciens qui positionnent les antennes relais MTN ou Orange utilisent le même principe pour garantir qu'une antenne « couvre » deux localités sous un angle donné depuis un site optimal.
\end{savaistu}

\begin{exercice}[À toi de jouer]
On donne $AB = 8$ cm. 
\begin{enumerate}
\item Construire l'ensemble des points $M$ tels que $\text{mes}(\widehat{AMB}) = 30^{\circ}$.
\item Calculer le rayon des arcs capables obtenus.
\item Que vaut l'angle au centre $\widehat{AOB}$ ? En déduire la nature du triangle $OAB$... attention, $OAB$ est isocèle en $O$ : calcule $\text{mes}(\widehat{OAB})$.
\end{enumerate}
\end{exercice}

\begin{corrige}[Exercice 1]
\begin{enumerate}
\item On trace $[At)$ avec $\text{mes}(\widehat{tAB}) = 30^{\circ}$, la perpendiculaire en $A$ à $[At)$, puis la médiatrice de $[AB]$ ; leur intersection $O$ est le centre. On trace l'arc de centre $O$ passant par $A$ et $B$, puis son symétrique par rapport à $(AB)$, en excluant $A$ et $B$.
\item $R = \dfrac{AB}{2\sin 30^{\circ}} = \dfrac{8}{2 \times \frac{1}{2}} = 8$ cm. Remarque amusante : ici $R = AB$.
\item $\text{mes}(\widehat{AOB}) = 2x = 60^{\circ}$. Le triangle $OAB$ est isocèle en $O$ avec un angle au sommet de $60^{\circ}$ ; les angles à la base valent $\dfrac{180^{\circ} - 60^{\circ}}{2} = 60^{\circ}$ : le triangle $OAB$ est \textbf{équilatéral}, ce qui confirme $OA = OB = AB = 8$ cm.
\end{enumerate}
\end{corrige}

Nous savons désormais déterminer le lieu des points qui voient un segment sous un \emph{angle géométrique} donné. Mais ce lieu, fait de deux arcs symétriques, ne distingue pas les « côtés ». En munissant le plan d'une orientation et en travaillant avec des \emph{angles orientés}, nous pourrons séparer ces deux arcs et obtenir des lieux plus fins : c'est l'objet de la section suivante sur les arcs capables orientés.

\coursec{Arcs capables orientés : mes(MA, MB) = x + 2kπ et mes(MA, MB) = x + kπ}

Dans la section précédente, nous avons travaillé avec l'\emph{angle géométrique} $\widehat{AMB}$, un nombre compris entre $0$ et $\pi$, sans aucune notion de sens. Nous avons vu que l'ensemble des points $M$ tels que $\operatorname{mes}\widehat{AMB} = x$ est formé de \textbf{deux arcs de cercle} symétriques par rapport à la droite $(AB)$. Cette réponse en deux morceaux est un peu frustrante : pourquoi deux arcs alors que la condition ne mentionne qu'une seule mesure ? La réponse tient en un mot : l'\emph{orientation}. En munissant le plan d'un sens de rotation — le sens trigonométrique — et en utilisant les angles \emph{orientés}, nous allons voir que chaque arc reçoit sa propre mesure, et que les lieux deviennent plus précis, plus fins. C'est toute la puissance des angles orientés, que tu as découverts en Seconde et que nous allons maintenant exploiter à fond.

\courssub{Le plan orienté : rappel essentiel}

\begin{definition}[Plan orienté, sens direct]
Orienter le plan, c'est choisir une convention de sens de rotation. Par convention universelle, on appelle \cle{sens direct} (ou \cle{sens trigonométrique}) le sens inverse des aiguilles d'une montre ; le sens opposé est le \cle{sens indirect} (ou sens horaire). Un plan muni de cette convention est appelé \cle{plan orienté}.
\end{definition}

Dans un plan orienté, un angle de vecteurs $(\vec{u}, \vec{v})$ possède une infinité de mesures, toutes de la forme $\theta + 2k\pi$ avec $k \in \mathbb{Z}$ ; la \cle{mesure principale} est celle qui appartient à $]-\pi\,;\,\pi]$. Un angle de droites $(d, d')$, lui, est défini \emph{modulo $\pi$} seulement, car une droite n'a pas de sens de parcours privilégié : tourner d'un angle plat ramène la droite sur elle-même.

\begin{aretenir}[Les trois cadres à ne jamais confondre]
\begin{itemize}
\item \textbf{Angle géométrique} $\widehat{AMB}$ : mesure dans $[0\,;\,\pi]$, sans signe, sans orientation.
\item \textbf{Angle orienté de vecteurs} $(\overrightarrow{MA}, \overrightarrow{MB})$ : mesure définie \emph{modulo $2\pi$}, avec un signe qui dépend du sens de rotation.
\item \textbf{Angle de droites} $(MA, MB)$ : mesure définie \emph{modulo $\pi$}.
\end{itemize}
La nature du lieu géométrique dépendra entièrement du cadre choisi par l'énoncé.
\end{aretenir}

\courssub{Ensemble des points M tels que mes(MA, MB) = x + 2kπ}

\textbf{Intuition.} Reprenons le segment $[AB]$ et les deux arcs symétriques de la section précédente. Plaçons-nous dans le plan orienté. Si $M$ parcourt l'arc situé « au-dessus » de $(AB)$, le passage de $\overrightarrow{MA}$ à $\overrightarrow{MB}$ se fait, disons, dans le sens direct : la mesure orientée est $+x$. Si $M$ parcourt l'arc « en dessous », le même passage se fait dans le sens indirect : la mesure orientée est $-x$ (ou $2\pi-x$). Autrement dit, l'orientation \emph{sépare} les deux arcs : chacun porte sa propre mesure signée. La condition $\operatorname{mes}(\overrightarrow{MA}, \overrightarrow{MB}) = x \;[2\pi]$ ne retient donc qu'\textbf{un seul} des deux arcs.

\begin{propriete}[Arc capable orienté — admis]
Soient $A$ et $B$ deux points distincts du plan orienté et $x$ un réel. L'ensemble des points $M$ du plan tels que
\[ \operatorname{mes}\left(\overrightarrow{MA}, \overrightarrow{MB}\right) = x + 2k\pi, \quad k \in \mathbb{Z}, \]
est un \textbf{arc de cercle d'extrémités $A$ et $B$, privé des points $A$ et $B$}, appelé \cle{arc capable orienté}. L'autre arc du même cercle correspond à la mesure $x - \pi \;[2\pi]$ (c'est-à-dire à l'opposé de la mesure géométrique complémentaire).
\end{propriete}

\begin{experience}[Pourquoi un seul arc ? Justification guidée]
Soit $\Gamma$ le cercle passant par $A$, $B$ et un point $M_0$ tel que $\operatorname{mes}(\overrightarrow{M_0A}, \overrightarrow{M_0B}) = x \;[2\pi]$.
\begin{enumerate}
\item Pour tout point $M$ du \emph{même arc} d'extrémités $A$ et $B$ que $M_0$, le théorème de l'angle inscrit orienté donne $\operatorname{mes}(\overrightarrow{MA}, \overrightarrow{MB}) = \operatorname{mes}(\overrightarrow{M_0A}, \overrightarrow{M_0B}) = x \;[2\pi]$ : tout l'arc convient.
\item Soit maintenant $M'$ sur l'\emph{autre} arc du \emph{même cercle} $\Gamma$. Les angles inscrits orientés qui interceptent des arcs complémentaires vérifient $\operatorname{mes}(\overrightarrow{M'A}, \overrightarrow{M'B}) = x - \pi \;[2\pi]$. En effet, les mesures géométriques sont supplémentaires (elles valent $x$ et $\pi - x$ si $0 < x < \pi$) et les sens de rotation sont opposés.
\item Comme $x - \pi \neq x \;[2\pi]$ (leur différence est $\pi$, non multiple de $2\pi$), aucun point du second arc ne vérifie la condition. Le lieu est donc \emph{exactement} un arc, privé de $A$ et $B$ (où les vecteurs $\overrightarrow{MA}$, $\overrightarrow{MB}$ ne sont pas tous deux définis).
\end{enumerate}
\end{experience}

\begin{popfigure}
\begin{center}
\begin{tikzpicture}[scale=1.6,>=Stealth]
\coordinate (A) at (-1.3,0);
\coordinate (B) at (1.3,0);
\coordinate (O) at (0,0.75);
\coordinate (M) at (0,1.75);
% Arc supérieur (le lieu) en popblue, fléché sens direct
\draw[popblue,very thick] (A) arc[start angle=180,end angle=0,radius=1.3];
\draw[popblue,->,very thick] (0.92,0.92) arc[start angle=45,end angle=90,radius=1.3];
% Arc inférieur (exclu) en gris pointillé
\draw[popmuted,dashed,thick] (A) arc[start angle=-180,end angle=0,radius=1.3];
\draw[popline] (A)--(B);
\draw[popink] (M)--(A) (M)--(B);
% Flèche d'angle en M
\draw[poporange,->,thick] ($(M)+(230:0.45)$) arc[start angle=230,end angle=310,radius=0.45];
\node[poporange] at ($(M)+(-90:0.68)$) {\footnotesize $+\frac{\pi}{3}$};
\foreach \p/\pos in {A/below left,B/below right,M/above}{\fill[popdark] (\p) circle(1.4pt); \node[\pos,popdark] at (\p) {$\p$};}
\node[popblue] at (0,2.05) {\footnotesize arc : $\operatorname{mes}(\overrightarrow{MA},\overrightarrow{MB})=\frac{\pi}{3}\;[2\pi]$};
\node[popmuted] at (0,-1.75) {\footnotesize arc exclu : mesure $-\frac{2\pi}{3}\;[2\pi]$};
\end{tikzpicture}
\end{center}
\legende{Arc capable orienté : seul l'arc supérieur (fléché dans le sens direct) répond à la condition $\operatorname{mes}(\overrightarrow{MA},\overrightarrow{MB}) = \frac{\pi}{3}\;[2\pi]$.}
\end{popfigure}

\textbf{Cas particuliers fondamentaux.} Que devient l'arc capable lorsque la mesure est un multiple de $\pi$ ? L'arc « s'aplatit » sur la droite $(AB)$ :

\begin{propriete}[Alignement en termes d'angles orientés]
Soient $A$ et $B$ deux points distincts.
\begin{itemize}
\item $\operatorname{mes}(\overrightarrow{MA}, \overrightarrow{MB}) = 0 \;[2\pi]$ si et seulement si $M$ appartient à la droite $(AB)$ \textbf{privée du segment $[AB]$} (les vecteurs $\overrightarrow{MA}$ et $\overrightarrow{MB}$ sont alors colinéaires de même sens).
\item $\operatorname{mes}(\overrightarrow{MA}, \overrightarrow{MB}) = \pi \;[2\pi]$ si et seulement si $M$ appartient au \textbf{segment ouvert $]AB[$} (les vecteurs sont colinéaires de sens contraires).
\end{itemize}
\end{propriete}

Ces deux résultats sont précieux : ils permettent de démontrer qu'un point appartient à un segment ou qu'il est extérieur à celui-ci, uniquement par le calcul d'angles orientés, sans figure auxiliaire.

\courssub{Ensemble des points M tels que mes(MA, MB) = x + kπ}

\textbf{Intuition.} Passons maintenant aux \emph{droites} $(MA)$ et $(MB)$. Une droite n'a pas de sens : la droite $(MA)$ ne « sait pas » si on la parcourt de $M$ vers $A$ ou de $A$ vers $M$. Conséquence immédiate : les mesures $x$ et $x - \pi$, qui distinguaient les deux arcs dans le cadre des vecteurs, deviennent \emph{indiscernables} modulo $\pi$. Sur un \emph{même} cercle, les deux arcs se réunissent donc. Mais il existe \emph{deux} cercles symétriques par rapport à $(AB)$ sur lesquels l'angle géométrique vaut $x$ ou $\pi-x$ ; comme $x \equiv \pi-x \pmod{\pi}$, \textbf{les deux cercles entiers} conviennent (sauf cas particulier $x=\frac{\pi}{2}$ où ils se confondent).

\begin{propriete}[Lieu pour un angle de droites — admis]
Soient $A$ et $B$ deux points distincts du plan et $x$ un réel. L'ensemble des points $M$ tels que
\[ \operatorname{mes}\left((MA), (MB)\right) = x + k\pi, \quad k \in \mathbb{Z}, \]
est :
\begin{itemize}
\item si $x \not\equiv 0 \;[\pi]$ : un \textbf{cercle passant par $A$ et $B$}, privé des points $A$ et $B$ (c'est la réunion des deux arcs capables de mesures $x$ et $x - \pi$ dans le cadre des vecteurs, qui sont deux arcs d'un \emph{même} cercle) ;
\item si $x \equiv 0 \;[\pi]$ : la \textbf{droite $(AB)$} privée des points $A$ et $B$.
\end{itemize}
Le centre $\Omega$ du cercle est sur la médiatrice de $[AB]$ ; il est déterminé par la tangente en $A$, qui fait avec $(AB)$ l'angle de mesure $x$.
\end{propriete}

\begin{propriete}[Cas particulier : l'angle droit]
L'ensemble des points $M$ tels que $\operatorname{mes}((MA), (MB)) = \dfrac{\pi}{2} \;[\pi]$ est le \textbf{cercle de diamètre $[AB]$ privé de $A$ et $B$}. C'est la reformulation orientée du théorème du triangle rectangle inscrit dans un demi-cercle.
\end{propriete}

\begin{exemplebox}[Un exemple chiffré pour comparer les deux cadres]
Soient $A$ et $B$ deux points distincts. Déterminons :
\begin{enumerate}
\item l'ensemble $\mathcal{E}_1$ des points $M$ tels que $\operatorname{mes}((MA),(MB)) = \dfrac{\pi}{2} \;[\pi]$ ;
\item l'ensemble $\mathcal{E}_2$ des points $M$ tels que $\operatorname{mes}(\overrightarrow{MA},\overrightarrow{MB}) = \dfrac{\pi}{2} \;[2\pi]$.
\end{enumerate}
\textbf{Réponse.} $\mathcal{E}_1$ est le \emph{cercle entier} de diamètre $[AB]$, privé de $A$ et $B$ : en effet, pour tout point $M$ de ce cercle, le triangle $AMB$ est rectangle en $M$, et l'angle des droites $(MA)$ et $(MB)$ vaut $\frac{\pi}{2}$ modulo $\pi$, quel que soit l'arc où se trouve $M$. En revanche, $\mathcal{E}_2$ n'est qu'\emph{un seul arc} de ce même cercle : celui où le passage de $\overrightarrow{MA}$ à $\overrightarrow{MB}$ s'effectue dans le sens direct de $+\frac{\pi}{2}$. Sur l'autre arc, la mesure orientée vaut $-\frac{\pi}{2} \;[2\pi]$, ce qui ne convient pas. On a donc $\mathcal{E}_2 \subsetneq \mathcal{E}_1$ : le cadre $[2\pi]$ est deux fois plus sélectif.
\end{exemplebox}

\begin{methode}[Choisir le bon cadre avant de conclure]
Face à un problème de lieu géométrique faisant intervenir un angle :
\begin{enumerate}
\item \textbf{Lire attentivement l'écriture} : $\widehat{AMB}$ (angle géométrique), $(\overrightarrow{MA},\overrightarrow{MB})$ (vecteurs, modulo $2\pi$) ou $(MA),(MB)$ (droites, modulo $\pi$).
\item \textbf{Identifier la congruence} : $[2\pi]$ impose le signe de la mesure, donc sélectionne un seul arc sur un cercle ; $[\pi]$ ignore le signe, donc réunit les deux arcs d'un \emph{même} cercle (les arcs de mesures $x$ et $x-\pi$) — ce qui donne un cercle complet.
\item \textbf{Conclure avec précision} : deux arcs symétriques sur deux cercles (angle géométrique), un seul arc sur un cercle (vecteurs, $[2\pi]$), un cercle entier (droites, $[\pi]$, $x \not\equiv 0$).
\item \textbf{Ne jamais oublier d'exclure $A$ et $B$} : en ces points, l'un des vecteurs est nul ou l'une des droites n'est pas définie, l'angle n'existe pas.
\item Vérifier les cas dégénérés : $x = 0$ ou $\pi$ $[2\pi]$ donnent des sous-ensembles de la droite $(AB)$.
\end{enumerate}
\end{methode}

\begin{attention}[Les deux erreurs classiques au Probatoire]
\begin{itemize}
\item \textbf{Erreur n°1} : conclure « le lieu est un cercle » alors que la condition est donnée \emph{modulo $2\pi$}. Avec des vecteurs et une congruence $[2\pi]$, le lieu est \textbf{un seul arc}, pas le cercle entier. Le cercle complet est le lieu pour un angle de \emph{droites} (modulo $\pi$).
\item \textbf{Erreur n°2} : oublier d'écrire « privé de $A$ et $B$ ». En $M = A$ ou $M = B$, l'angle n'est pas défini ; un lieu qui contiendrait $A$ ou $B$ est une réponse fausse, et les examinateurs y sont attentifs.
\end{itemize}
\end{attention}

\begin{popfigure}
\begin{center}
\begin{tikzpicture}[scale=1.15,>=Stealth]
% --- Figure 1 : angle géométrique : deux arcs sur deux cercles ---
\begin{scope}[shift={(-5.2,0)}]
\coordinate (A1) at (-1,0); \coordinate (B1) at (1,0);
% Cercle 1 (centre haut)
\draw[popgreen,very thick] (A1) arc[start angle=180,end angle=0,radius=1];
% Cercle 2 (centre bas)
\draw[popgreen,very thick] (A1) arc[start angle=-180,end angle=0,radius=1];
\draw[popline] (A1)--(B1);
\fill[popdark] (A1) circle(1.5pt) node[below left]{\footnotesize $A$};
\fill[popdark] (B1) circle(1.5pt) node[below right]{\footnotesize $B$};
\node[popgreen,align=center] at (0,-1.7) {\footnotesize géométrique $[0,\pi]$ :\\ \footnotesize deux arcs sur deux cercles};
\end{scope}
% --- Figure 2 : vecteurs [2π] : un arc fléché ---
\begin{scope}[shift={(0,0)}]
\coordinate (A2) at (-1,0); \coordinate (B2) at (1,0);
\draw[popblue,very thick,->] (A2) arc[start angle=180,end angle=90,radius=1];
\draw[popblue,very thick] (0,1) arc[start angle=90,end angle=0,radius=1];
\draw[popmuted,dashed] (A2) arc[start angle=-180,end angle=0,radius=1];
\draw[popline] (A2)--(B2);
\fill[popdark] (A2) circle(1.5pt) node[below left]{\footnotesize $A$};
\fill[popdark] (B2) circle(1.5pt) node[below right]{\footnotesize $B$};
\node[popblue,align=center] at (0,-1.7) {\footnotesize vecteurs $[2\pi]$ :\\ \footnotesize un seul arc};
\end{scope}
% --- Figure 3 : droites [π] : un cercle entier ---
\begin{scope}[shift={(5.2,0)}]
\coordinate (A3) at (-1,0); \coordinate (B3) at (1,0);
\draw[poporange,very thick] (0,0.577) circle[radius=1.155];
\draw[popline] (A3)--(B3);
\fill[popdark] (A3) circle(1.5pt) node[below left]{\footnotesize $A$};
\fill[popdark] (B3) circle(1.5pt) node[below right]{\footnotesize $B$};
\node[poporange,align=center] at (0,-1.7) {\footnotesize droites $[\pi]$ :\\ \footnotesize un cercle entier};
\end{scope}
\end{tikzpicture}
\end{center}
\legende{Le même segment $[AB]$, trois conditions angulaires, trois lieux différents (dans tous les cas, $A$ et $B$ sont exclus). Pour les droites modulo $\pi$, les arcs de mesures $x$ et $x-\pi$ se réunissent en un cercle entier ; pour $x=\frac{\pi}{2}$, c'est le cercle de diamètre $[AB]$.}
\end{popfigure}

\begin{aretenir}[Tableau récapitulatif des trois types de lieux]
\begin{center}
\begin{tabularx}{\linewidth}{|l|X|X|}
\hline
\rowcolor{popblueL} \textbf{Condition} & \textbf{Nature du lieu} & \textbf{Cas particulier} \\
\hline
$\operatorname{mes}\widehat{AMB} = x$, $x \in [0\,;\,\pi]$ & Deux arcs de cercle symétriques par rapport à $(AB)$, privés de $A$ et $B$ (sur deux cercles distincts) & $x = \frac{\pi}{2}$ : cercle de diamètre $[AB]$ privé de $A$ et $B$ (les deux cercles se confondent) \\
\hline
$\operatorname{mes}(\overrightarrow{MA},\overrightarrow{MB}) = x \;[2\pi]$ & Un seul arc de cercle d'extrémités $A$ et $B$, privé de $A$ et $B$ & $x = 0$ : $(AB)\setminus[AB]$ ; $x = \pi$ : segment $]AB[$ \\
\hline
$\operatorname{mes}((MA),(MB)) = x \;[\pi]$ & Un cercle passant par $A$ et $B$, privé de $A$ et $B$ & $x = \frac{\pi}{2}$ : cercle de diamètre $[AB]$ ; $x \equiv 0$ : droite $(AB)$ privée de $A$ et $B$ \\
\hline
\end{tabularx}
\end{center}
\end{aretenir}

\begin{exoresolu}[Distinguer les trois cadres sur une même donnée]
Soient $A$ et $B$ deux points distincts du plan orienté. Déterminer et représenter mentalement :
\begin{enumerate}
\item l'ensemble $\mathcal{E}_1$ des points $M$ tels que $\operatorname{mes}\widehat{AMB} = \dfrac{\pi}{3}$ ;
\item l'ensemble $\mathcal{E}_2$ des points $M$ tels que $\operatorname{mes}(\overrightarrow{MA},\overrightarrow{MB}) = \dfrac{\pi}{3} \;[2\pi]$ ;
\item l'ensemble $\mathcal{E}_3$ des points $M$ tels que $\operatorname{mes}((MA),(MB)) = \dfrac{\pi}{3} \;[\pi]$.
\end{enumerate}
\tcblower
\textbf{Solution détaillée.}
\begin{enumerate}
\item $\mathcal{E}_1$ : l'angle géométrique ne distingue pas les positions de $M$ par rapport à $(AB)$. Le lieu est formé de \textbf{deux arcs de cercle} symétriques par rapport à la droite $(AB)$, privés de $A$ et $B$. Ces arcs appartiennent à deux cercles distincts, symétriques par rapport à $(AB)$, passant par $A$ et $B$.
\item $\mathcal{E}_2$ : la congruence est modulo $2\pi$ avec des vecteurs : le signe compte. Seul l'arc où le passage de $\overrightarrow{MA}$ à $\overrightarrow{MB}$ se fait dans le sens direct convient. $\mathcal{E}_2$ est donc \textbf{un seul arc} de cercle d'extrémités $A$ et $B$, privé de $A$ et $B$ — c'est l'une des deux moitiés de l'un des cercles de $\mathcal{E}_1$. Sur l'arc exclu de ce même cercle, la mesure vaut $\frac{\pi}{3} - \pi = -\frac{2\pi}{3} \;[2\pi]$.
\item $\mathcal{E}_3$ : modulo $\pi$, les mesures $\frac{\pi}{3}$ et $\frac{\pi}{3} - \pi = -\frac{2\pi}{3}$ sont confondues (elles diffèrent de $\pi$). Or, sur le cercle qui porte $\mathcal{E}_2$, l'arc exclu est précisément celui où la mesure vaut $-\frac{2\pi}{3}$. $\mathcal{E}_3$ est donc \textbf{le cercle entier} qui porte $\mathcal{E}_2$, privé de $A$ et $B$. (Sur le cercle symétrique, les mesures valent $-\frac{\pi}{3}$ et $\frac{2\pi}{3}$ : il n'est pas dans $\mathcal{E}_3$.)
\end{enumerate}
On observe les inclusions : $\mathcal{E}_2 \subset \mathcal{E}_1$ et $\mathcal{E}_2 \subset \mathcal{E}_3$. $\mathcal{E}_1$ est la réunion de deux arcs (un sur chaque cercle), $\mathcal{E}_3$ est un cercle complet. Cet exercice montre qu'une lecture précise de la congruence est la clé de toute la question.
\end{exoresolu}

\begin{savaistu}[Vers la Terminale : les similitudes directes]
Les arcs capables orientés que tu viens d'étudier sont bien plus qu'un exercice de style : ils constituent l'un des outils majeurs de la géométrie de Terminale C et E. Une \emph{similitude directe} est une transformation qui conserve les angles orientés ; caractériser les points $M$ dont l'image $M'$ vérifie une condition d'angle orienté revient précisément à manipuler des arcs capables orientés. Le fait que la condition $\operatorname{mes}(\overrightarrow{MA},\overrightarrow{MB}) = x \;[2\pi]$ sélectionne \emph{un seul} arc — et non deux — est exactement ce qui rend les similitudes \emph{directes} si agréables à utiliser : elles respectent le sens trigonométrique, donc la fine structure des lieux. Maîtriser cette section en Première, c'est te construire une avance décisive pour l'année prochaine et pour le Baccalauréat.
\end{savaistu}

\coursec{Applications : lieux géométriques et cocyclicité}

Dans la section précédente, tu as appris à décrire et à construire les ensembles de points $M$ vérifiant $\text{mes}(\overrightarrow{MA},\overrightarrow{MB}) = x + 2k\pi$ ou $\text{mes}(\overrightarrow{MA},\overrightarrow{MB}) = x + k\pi$ : ce sont des \cle{arcs capables} (ouverts ou privés de leurs extrémités selon les cas). Il est temps maintenant de retourner l'outil : au lieu de partir d'une condition angulaire pour trouver un lieu, nous partons d'un \emph{problème} (lieu géométrique, cocyclicité, visée) et nous le \emph{ramenons} à une condition angulaire. C'est exactement la démarche exigée au Probatoire, et c'est l'objet de cette dernière section.

\courssub{Stratégie générale de recherche d'un lieu géométrique}

\textbf{Intuition.} Chercher un lieu, c'est répondre à la question : « où peuvent se trouver les points $M$ qui vérifient telle propriété ? ». Quand la propriété fait intervenir un angle sous lequel on voit un segment fixe $[AB]$, le théorème de l'arc capable fait tout le travail : le lieu est un arc de cercle (ou une réunion d'arcs) passant par $A$ et $B$.

\begin{methode}[Déterminer un lieu géométrique par l'arc capable]
\begin{enumerate}
\item \textbf{Identifier les points fixes.} Repérer dans l'énoncé les points qui ne bougent pas (souvent notés $A$ et $B$, ou $B$ et $C$) et le point variable $M$.
\item \textbf{Traduire la condition en égalité angulaire constante.} Montrer, à l'aide des propriétés de la figure (angles inscrits, cocyclicité, transformations), que la condition équivaut à $\text{mes}(\overrightarrow{MA},\overrightarrow{MB}) = x + k\pi$ (ou $+2k\pi$), où $x$ est une constante.
\item \textbf{Conclure par le théorème de l'arc capable.} Le lieu est l'arc capable (ou la réunion des deux arcs symétriques par rapport à $(AB)$) construit sur $[AB]$ pour l'angle $x$, privé des points $A$ et $B$ et de toute position interdite par l'énoncé.
\item \textbf{Vérifier la réciproque.} S'assurer que tout point de l'ensemble trouvé vérifie bien la condition initiale, et exclure les points parasites (positions où un angle n'est pas défini, où une longueur s'annule, etc.).
\end{enumerate}
\end{methode}

\begin{attention}[Les deux oublis classiques dans un problème de lieu]
Deux erreurs coûtent cher au Probatoire :
\begin{itemize}
\item \textbf{Oublier la réciproque.} Tu as montré que tout point solution est sur l'arc capable : très bien, mais il faut aussi dire que, réciproquement, tout point de cet arc (hors positions interdites) est solution. Sans cela, le lieu n'est pas \emph{exactement} déterminé.
\item \textbf{Oublier les positions interdites.} Les points $A$ et $B$ eux-mêmes sont toujours exclus (l'angle $(\overrightarrow{MA},\overrightarrow{MB})$ n'y est pas défini). Il peut y en avoir d'autres : par exemple, si l'énoncé impose que $M$ soit le sommet d'un triangle non aplati, les points de la droite $(AB)$ sont exclus.
\end{itemize}
Rédige toujours la conclusion ainsi : « Le lieu de $M$ est l'arc capable $\Gamma$ privé des points $A$ et $B$ ».
\end{attention}

\begin{exoresolu}[Lieu du sommet d'un triangle d'angle fixe]
\textbf{Énoncé.} $B$ et $C$ sont deux points fixes tels que $BC = 5$ cm. Déterminer et construire le lieu des points $A$ tels que $\text{mes}(\widehat{BAC}) = 60^{\circ}$ (triangle $ABC$ non aplati).
\tcblower
\textbf{Solution détaillée.}

\emph{Étape 1 — Points fixes et point variable.} Les points fixes sont $B$ et $C$ ; le point variable est $A$.

\emph{Étape 2 — Traduction angulaire.} Dire que $\text{mes}(\widehat{BAC}) = 60^{\circ}$, c'est dire que l'angle géométrique $\text{mes}(\widehat{BAC}) = \dfrac{\pi}{3}$, donc
\[ \text{mes}(\overrightarrow{AB},\overrightarrow{AC}) = \pm\frac{\pi}{3} + k\pi. \]
Autrement dit, $A$ voit le segment $[BC]$ sous un angle de mesure $\dfrac{\pi}{3}$ (au sens géométrique).

\emph{Étape 3 — Conclusion par l'arc capable.} Le lieu de $A$ est la réunion des deux arcs capables construits sur le segment $[BC]$ pour l'angle $\dfrac{\pi}{3}$, symétriques l'un de l'autre par rapport à la droite $(BC)$, privés des points $B$ et $C$.

\emph{Construction.} Le cercle contenant l'arc capable a pour rayon $R$ tel que $BC = 2R\sin\dfrac{\pi}{3}$, soit $R = \dfrac{5}{2 \times \frac{\sqrt{3}}{2}} = \dfrac{5\sqrt{3}}{3} \approx \num{2,89}$ cm. On construit la demi-tangente en $B$ faisant un angle de $60^{\circ}$ avec $(BC)$, puis la perpendiculaire en $B$ à cette demi-tangente ; elle coupe la médiatrice de $[BC]$ au centre $O$ du cercle. On trace l'arc, puis son symétrique par rapport à $(BC)$.

\emph{Étape 4 — Réciproque.} Tout point $A$ de ces arcs (distinct de $B$ et $C$) voit $[BC]$ sous l'angle $\dfrac{\pi}{3}$ au sens géométrique, donc convient. Les points $B$ et $C$ sont exclus (triangle aplati).

\textbf{Conclusion :} le lieu de $A$ est la réunion des deux arcs capables d'angle $60^{\circ}$ construits sur $[BC]$, privée de $B$ et $C$.
\end{exoresolu}

\courssub{Cocyclicité : quatre points sur un même cercle}

\textbf{Intuition.} Quatre points $A$, $B$, $M$, $N$ sont \cle{cocycliques} s'il existe un cercle passant par les quatre. Comment le prouver sans tracer le cercle ? L'idée est lumineuse : si $M$ et $N$ sont sur un même cercle passant par $A$ et $B$, alors ils voient la corde $[AB]$ sous des angles égaux (théorème de l'angle inscrit). Réciproquement, si $M$ voit $[AB]$ sous le même angle (modulo $\pi$) que $N$, alors $M$ est sur le cercle $(ABN)$.

\begin{propriete}[Caractérisation de la cocyclicité]
Soient $A$, $B$, $N$ trois points non alignés et $M$ un point distinct de $A$ et $B$. Alors :
\[ M \in \text{cercle }(ABN) \iff \text{mes}(\overrightarrow{MA},\overrightarrow{MB}) = \text{mes}(\overrightarrow{NA},\overrightarrow{NB}) \ [\pi]. \]
Autrement dit, quatre points $A$, $B$, $M$, $N$ sont cocycliques si et seulement si $M$ et $N$ voient le segment $[AB]$ sous le même angle modulo $\pi$.
\end{propriete}

\begin{experience}[Démonstration guidée de la caractérisation]
On raisonne en deux temps.
\begin{enumerate}
\item \textbf{Sens direct.} Supposons $M$ sur le cercle $(ABN)$. Par le théorème de l'angle inscrit orienté, les angles $(\overrightarrow{MA},\overrightarrow{MB})$ et $(\overrightarrow{NA},\overrightarrow{NB})$ interceptent la même corde $[AB]$ du même cercle, donc
\[ \text{mes}(\overrightarrow{MA},\overrightarrow{MB}) = \text{mes}(\overrightarrow{NA},\overrightarrow{NB}) \ [\pi]. \]
\item \textbf{Réciproque.} Supposons l'égalité modulo $\pi$. Le lieu des points $P$ tels que $\text{mes}(\overrightarrow{PA},\overrightarrow{PB}) = \text{mes}(\overrightarrow{NA},\overrightarrow{NB}) + k\pi$ est, d'après la section précédente, le cercle passant par $A$, $B$ et $N$ (privé de $A$ et $B$). Comme $M$ vérifie la condition, $M$ appartient à ce cercle. \hfill$\square$
\end{enumerate}
\end{experience}

\begin{methode}[Prouver que quatre points sont cocycliques]
\begin{enumerate}
\item Choisir deux des quatre points comme « corde de référence », disons $A$ et $B$.
\item Calculer ou justifier $\text{mes}(\overrightarrow{MA},\overrightarrow{MB})$ et $\text{mes}(\overrightarrow{NA},\overrightarrow{NB})$ (souvent à l'aide de la relation de Chasles sur les angles orientés).
\item Conclure : si les deux mesures sont égales modulo $\pi$, alors $A$, $B$, $M$, $N$ sont cocycliques.
\end{enumerate}
Astuce : on peut aussi montrer que $(\overrightarrow{MA},\overrightarrow{MB}) = (\overrightarrow{NA},\overrightarrow{NB}) \ [\pi]$ en passant par un angle intermédiaire grâce à Chasles.
\end{methode}

\begin{popfigure}
\begin{tikzpicture}[scale=2.2]
\draw[popline] (0,0) circle (1);
\coordinate (A) at (200:1);
\coordinate (B) at (340:1);
\coordinate (M) at (80:1);
\coordinate (N) at (140:1);
\draw[popblue,thick] (A)--(M)--(B);
\draw[poporange,thick] (A)--(N)--(B);
\draw[popmuted] (A)--(B);
\foreach \p/\pos in {A/below left, B/below right, M/above, N/above left}
  \fill[popdark] (\p) circle (0.018) node[\pos, popdark] {$\p$};
\draw[popblue,->] ($(M)+(210:0.28)$) arc (210:285:0.28);
\node[popblue] at ($(M)+(-90:0.42)$) {\small $\alpha$};
\draw[poporange,->] ($(N)+(230:0.28)$) arc (230:305:0.28);
\node[poporange] at ($(N)+(-90:0.42)$) {\small $\alpha$};
\end{tikzpicture}
\legende{Quatre points cocycliques : $M$ et $N$ voient la corde $[AB]$ sous le même angle $\alpha$ (modulo $\pi$), donc $A$, $B$, $M$, $N$ sont sur un même cercle.}
\end{popfigure}

\begin{exoresolu}[Montrer une cocyclicité]
\textbf{Énoncé.} Soit $ABC$ un triangle et $H$ son orthocentre. On note $A'$ le pied de la hauteur issue de $A$ et $B'$ le pied de la hauteur issue de $B$. Montrer que les points $A$, $B$, $A'$, $B'$ sont cocycliques.
\tcblower
\textbf{Solution.} Par définition des hauteurs, $(AA') \perp (BC)$ en $A'$ et $(BB') \perp (AC)$ en $B'$. Donc :
\[ \text{mes}(\overrightarrow{A'A},\overrightarrow{A'B}) = \text{mes}(\overrightarrow{B'A},\overrightarrow{B'B}) = \frac{\pi}{2} \ [\pi], \]
car $\overrightarrow{A'A}$ est porté par $(AA')$ et $\overrightarrow{A'B}$ par $(BC)$ (angle droit), et de même $\overrightarrow{B'A}$ est porté par $(AC)$ et $\overrightarrow{B'B}$ par $(BB')$ (angle droit).

Ainsi $A'$ et $B'$ voient le segment $[AB]$ sous le même angle droit modulo $\pi$ : d'après la caractérisation de la cocyclicité, $A$, $B$, $A'$, $B'$ sont cocycliques. (Le cercle en question est le cercle de diamètre $[AB]$, puisqu'un angle inscrit droit intercepte un diamètre.)
\end{exoresolu}

\courssub{Intersection de lieux : double condition angulaire}

\textbf{Intuition.} Un point du plan est déterminé par deux conditions. Si chacune se traduit par un arc capable, le point cherché est à l'\emph{intersection des deux arcs}. C'est le principe du « tir croisé » utilisé en navigation ou en géolocalisation : viser deux repères connus sous deux angles connus suffit à se positionner.

\begin{methode}[Point vérifiant deux conditions angulaires]
\begin{enumerate}
\item Traduire la première condition : $M$ appartient à l'arc capable $\Gamma_1$ construit sur $[AB]$ pour l'angle $x$.
\item Traduire la seconde condition : $M$ appartient à l'arc capable $\Gamma_2$ construit sur $[CD]$ pour l'angle $y$ (les segments peuvent partager un point).
\item Les points solutions sont les intersections de $\Gamma_1$ et $\Gamma_2$, hors positions interdites. Discuter le nombre de solutions selon la figure (0, 1 ou 2 points en général).
\end{enumerate}
\end{methode}

\begin{exemplebox}[Double visée]
Soient $A$, $B$, $C$ trois points fixes avec $C$ distinct de $A$ et $B$. On cherche les points $M$ tels que
\[ \text{mes}(\overrightarrow{MA},\overrightarrow{MB}) = \frac{\pi}{3} + k\pi \quad \text{et} \quad \text{mes}(\overrightarrow{MB},\overrightarrow{MC}) = \frac{\pi}{4} + k\pi. \]
Le point $M$ est à l'intersection de l'arc capable construit sur $[AB]$ pour l'angle $\dfrac{\pi}{3}$ et de l'arc capable construit sur $[BC]$ pour l'angle $\dfrac{\pi}{4}$. Comme les deux cercles passent par $B$, ils se recoupent en un second point $M$ (sauf cas tangentiel) : c'est en général l'unique solution.
\end{exemplebox}

\begin{popfigure}
\begin{tikzpicture}[scale=1.5]
% Coordinates
\coordinate (A) at (0,0);
\coordinate (B) at (4,0);
\coordinate (C) at (4,3);
% Circle 1: Center O1 (2, 1.1547), R1=2.3094 (Capable arc 60 deg on AB)
\coordinate (O1) at (2, 1.1547);
% Circle 2: Center O2 (2.5, 1.5), R2=2.1213 (Capable arc 45 deg on BC)
\coordinate (O2) at (2.5, 1.5);
% Intersection M (approx 1.63, 3.43)
\coordinate (M) at (1.63, 3.43);

% Arc Capable 1 (Gamma_1) on AB, angle 60 deg (Major arc)
% Start at A (angle 210 from O1), end at B (angle 330 from O1) going CCW (210 -> 450)
\draw[popblue, thick] (A) arc (210:450:2.3094);

% Arc Capable 2 (Gamma_2) on BC, angle 45 deg (Major arc)
% Start at C (angle 45 from O2), end at B (angle 315 from O2) going CCW (45 -> 315)
\draw[poporange, thick] (C) arc (45:315:2.1213);

% Points
\fill[popdark] (A) circle (0.06) node[below left] {$A$};
\fill[popdark] (B) circle (0.06) node[below right] {$B$};
\fill[popdark] (C) circle (0.06) node[right] {$C$};
\fill[popdark] (M) circle (0.06) node[above left] {$M$};
\fill[popdark] (O1) circle (0.04) node[below] {$O_1$};
\fill[popdark] (O2) circle (0.04) node[below right] {$O_2$};

% Segments
\draw[popmuted, dashed] (M)--(A) (M)--(B) (M)--(C);
\draw[popmuted, thin] (A)--(B) (B)--(C);

% Labels
\node[popblue] at (2, 3.8) {$\Gamma_1$};
\node[poporange] at (0.8, 1.8) {$\Gamma_2$};
\end{tikzpicture}
\legende{Le point $M$, intersection des arcs capables $\Gamma_1$ (sur $[AB]$, angle $\frac{\pi}{3}$) et $\Gamma_2$ (sur $[BC]$, angle $\frac{\pi}{4}$), est l'unique point d'où l'on voit $[AB]$ sous $60^{\circ}$ et $[BC]$ sous $45^{\circ}$ du même côté.}
\end{popfigure}

\courssub{Application concrète : voir deux segments sous des angles donnés}

Les problèmes de visée sont l'application la plus parlante des arcs capables : un photographe au stade Ahmadou Ahidjo veut cadrer les deux poteaux de but sous un angle précis ; un technicien de CAMTEL positionne un relais pour couvrir deux antennes sous un angle donné ; un navigateur sur le Wouri se repère en visant deux balises. Dans tous les cas, le lieu des positions possibles est un arc capable, et deux visées croisées déterminent la position.

\begin{exoresolu}[Position par double visée]
\textbf{Énoncé.} Deux balises $A$ et $B$ distantes de $4$ km balisent l'entrée d'un port, et un phare $C$ est situé à $3$ km de $B$. Un bateau $M$ voit $[AB]$ sous un angle de $90^{\circ}$ et $[BC]$ sous un angle de $90^{\circ}$. Décrire l'ensemble des positions possibles de $M$.
\tcblower
\textbf{Solution.} Voir $[AB]$ sous un angle droit signifie $\text{mes}(\overrightarrow{MA},\overrightarrow{MB}) = \dfrac{\pi}{2} + k\pi$ : $M$ est sur le cercle de diamètre $[AB]$ (privé de $A$ et $B$), de centre le milieu de $[AB]$ et de rayon $2$ km. De même, $M$ est sur le cercle de diamètre $[BC]$ (privé de $B$ et $C$), de rayon $1,5$ km. Les positions possibles sont donc les points d'intersection de ces deux cercles : ils passent tous deux par $B$ et se recoupent en un second point $M$ (les centres et $B$ ne sont pas alignés dans la configuration générale). Le point $B$ est exclu (le bateau n'est pas sur la balise). \textbf{Conclusion :} le bateau occupe le second point d'intersection des cercles de diamètres $[AB]$ et $[BC]$.
\end{exoresolu}

\begin{savaistu}[L'arc capable sur le terrain de football]
Les gardiens de but et les tireurs de coups francs exploitent \emph{intuitivement} l'arc capable ! L'angle sous lequel un attaquant voit les deux poteaux est constant le long d'un arc de cercle passant par les deux poteaux, et il est d'autant plus grand que le cercle est petit. Quand un attaquant « ouvre son angle de tir », il se déplace en réalité vers un arc capable plus favorable. Le gardien, lui, avance pour réduire cet angle : il se place sur la bissectrice pour couvrir symétriquement les deux poteaux. Les meilleurs buteurs camerounais, de Roger Milla à Samuel Eto'o, ont fait de cette géométrie un réflexe — sans jamais écrire $\text{mes}(\overrightarrow{MA},\overrightarrow{MB})$ !
\end{savaistu}

\courssub{Lien avec les transformations et rédaction type}

\textbf{Un mot pour la Terminale.} L'arc capable possède une propriété remarquable vis-à-vis des transformations : l'image d'un cercle par une similitude (rotation, homothétie, similitude directe) est un cercle, et les similitudes directes \emph{conservent les angles orientés}. Ainsi, l'image d'un arc capable par une similitude directe est encore un arc capable de même angle : on dit que le lieu est \cle{invariant} en structure. Cette remarque, que tu retrouveras en Terminale avec les similitudes et les nombres complexes, permet de résoudre élégamment des problèmes de lieux quand le point variable est l'image d'un autre point décrivant un cercle.

\begin{methode}[Rédaction type d'un problème de lieu (modèle Probatoire)]
\begin{enumerate}
\item \textbf{Analyse.} « Supposons le point $M$ solution. » Traduire l'hypothèse en égalité angulaire : $\text{mes}(\overrightarrow{MA},\overrightarrow{MB}) = x + k\pi$. En déduire que $M$ appartient à l'arc capable $\Gamma$ construit sur $[AB]$ pour l'angle $x$.
\item \textbf{Construction.} Décrire la construction de $\Gamma$ (demi-tangente, perpendiculaire, médiatrice, centre, arc) et le tracer proprement.
\item \textbf{Justification (réciproque).} « Réciproquement, soit $M$ un point de $\Gamma$ distinct de $A$ et $B$... » Vérifier que $M$ satisfait la condition initiale et exclure les positions interdites.
\item \textbf{Conclusion.} « Le lieu cherché est l'arc capable $\Gamma$ privé des points $A$ et $B$ » (préciser les deux arcs symétriques si l'angle est géométrique).
\end{enumerate}
\end{methode}

\begin{attention}[Pièges de rédaction à éviter absolument]
\begin{itemize}
\item Ne jamais écrire « le lieu est le cercle » : c'est un \emph{arc} de cercle (souvent deux arcs symétriques), privé de points.
\item Ne pas confondre les congruences : $+2k\pi$ donne un seul arc orienté ; $+k\pi$ donne le cercle entier (privé de $A$ et $B$) ; un angle géométrique fixé donne deux arcs symétriques.
\item Dans les problèmes de cocyclicité, l'égalité à utiliser est toujours \emph{modulo $\pi$}, jamais modulo $2\pi$ (sinon tu imposes en plus la position du même côté de $(AB)$).
\end{itemize}
\end{attention}

\begin{exercice}[Pour t'entraîner]
Soit $[AB]$ un segment fixe de longueur $6$ cm.
\begin{enumerate}
\item Déterminer et construire le lieu des points $M$ tels que $\text{mes}(\overrightarrow{MA},\overrightarrow{MB}) = \dfrac{\pi}{4} + k\pi$.
\item Soit $C$ un point tel que $ABC$ soit équilatéral. Montrer que le cercle précédent et le cercle circonscrit à $ABC$ se coupent en deux points dont on précisera la nature angulaire.
\end{enumerate}
\end{exercice}

\begin{corrige}[Exercice 1]
\begin{enumerate}
\item La condition est du type $\text{mes}(\overrightarrow{MA},\overrightarrow{MB}) = x + k\pi$ avec $x = \dfrac{\pi}{4}$ : le lieu est le cercle passant par $A$ et $B$ dont l'arc capable est associé à l'angle $\dfrac{\pi}{4}$, privé de $A$ et $B$. Le rayon est $R = \dfrac{AB}{2\sin\frac{\pi}{4}} = \dfrac{6}{\sqrt{2}} = 3\sqrt{2}$ cm. Construction : demi-tangente en $A$ faisant $45^{\circ}$ avec $(AB)$, perpendiculaire en $A$, intersection avec la médiatrice de $[AB]$ : centre $O$.
\item Le cercle circonscrit à $ABC$ est l'arc capable de $[AB]$ pour l'angle $\dfrac{\pi}{3}$ (angle au sommet du triangle équilatéral). Les deux cercles passent par $A$ et $B$ : ce sont leurs deux points d'intersection. En ces points, les angles de visée de $[AB]$ valent respectivement $\dfrac{\pi}{4}$ et $\dfrac{\pi}{3}$ modulo $\pi$ selon le cercle considéré.
\end{enumerate}
\end{corrige}

\begin{aretenir}[L'essentiel de la section]
\begin{itemize}
\item Un lieu géométrique à condition angulaire se traite en quatre temps : points fixes, traduction $\text{mes}(\overrightarrow{MA},\overrightarrow{MB}) = \text{constante}$, arc capable, réciproque avec exclusion des positions interdites.
\item Cocyclicité : $M \in \text{cercle}(ABN) \iff \text{mes}(\overrightarrow{MA},\overrightarrow{MB}) = \text{mes}(\overrightarrow{NA},\overrightarrow{NB}) \ [\pi]$.
\item Un point vérifiant deux conditions angulaires est l'intersection de deux arcs capables : c'est le principe de la double visée.
\item Les similitudes directes conservent les angles orientés : l'image d'un arc capable est un arc capable (complément utile pour la Terminale).
\end{itemize}
\end{aretenir}

\newpage

\coursec{Activité d'intégration}

\courssub{Rappels sur les angles orientés de vecteurs}

\begin{definition}[Angle orienté de deux vecteurs]
Soit $\vec{u}$ et $\vec{v}$ deux vecteurs non nuls du plan. L'\cle{angle orienté} $(\vec{u},\vec{v})$ est la mesure de la rotation qui transforme $\vec{u}$ en $\vec{v}$. On le note $\mathrm{mes}(\vec{u},\vec{v})$. Il est défini \emph{modulo $2\pi$} : $\mathrm{mes}(\vec{u},\vec{v}) \in \mathbb{R}/2\pi\mathbb{Z}$.
\end{definition}

\begin{propriete}[Mesure principale]
On appelle \cle{mesure principale} l'unique représentant de $\mathrm{mes}(\vec{u},\vec{v})$ dans l'intervalle $]-\pi\,;\pi]$ (ou $[0\,;2\pi[$ selon les conventions). Au Probatoire, on utilise généralement $]-\pi\,;\pi]$.
\end{propriete}

\begin{propriete}[Relation de Chasles]
Pour trois vecteurs non nuls $\vec{u},\vec{v},\vec{w}$ :
\[ \mathrm{mes}(\vec{u},\vec{w}) = \mathrm{mes}(\vec{u},\vec{v}) + \mathrm{mes}(\vec{v},\vec{w}) \quad [2\pi]. \]
\end{propriete}

\begin{propriete}[Opposé et inverse]
\[ \mathrm{mes}(\vec{v},\vec{u}) = -\mathrm{mes}(\vec{u},\vec{v}) \quad [2\pi]. \]
\end{propriete}

\begin{aretenir}[Angle orienté et angle géométrique (non orienté)]
L'angle géométrique $\widehat{AMB}$ (ou $\mathrm{mes}(\widehat{AMB})$) est toujours pris dans $[0\,;\pi]$. Il est lié à l'angle orienté par :
\[ \mathrm{mes}(\widehat{AMB}) = \min\bigl(|\mathrm{mes}(\overrightarrow{MA},\overrightarrow{MB})|,\ 2\pi - |\mathrm{mes}(\overrightarrow{MA},\overrightarrow{MB})|\bigr). \]
Autrement dit, l'angle géométrique « oublie » le sens de rotation.
\end{aretenir}

\courssub{Objectifs de l'activité}

À l'issue de cette activité, tu seras capable de :
\begin{itemize}
\item traduire une contrainte concrète d'implantation en une condition angulaire du type $\mathrm{mes}(\widehat{AMB}) = x$ ;
\item construire un \cle{arc capable} à partir d'un angle inscrit dont un côté est une demi-tangente ;
\item calculer le rayon du cercle support d'un arc capable grâce à la relation entre \cle{angle inscrit} et \cle{angle au centre} ;
\item déterminer une position précise par \cle{intersection de lieux} géométriques (arc capable et médiatrice) ;
\item distinguer, dans le plan orienté, les lieux définis par $\mathrm{mes}(\overrightarrow{MA},\overrightarrow{MB}) = x + 2k\pi$ et par $\mathrm{mes}(\overrightarrow{MA},\overrightarrow{MB}) = x + k\pi$, et comprendre ce que l'orientation ajoute ou retire au lieu.
\end{itemize}

\courssub{Situation}

Le stade omnisports de Bafoussam, dans la région de l'Ouest, prépare la retransmission télévisée d'un match de la MTN Elite One. Deux poteaux d'éclairage $A$ et $B$, plantés de part et d'autre de la ligne médiane, sont distants de $AB = \SI{40}{m}$. La régie technique souhaite installer une caméra mobile $M$ sur le pourtour du terrain, de sorte que les deux poteaux — qui serviront de repères de cadrage — soient filmés sous un angle \emph{exact} de $60^{\circ}$ : autrement dit, on exige
\[ \mathrm{mes}(\widehat{AMB}) = 60^{\circ} = \frac{\pi}{3}\ \text{rad}. \]
Le chef cadreur se pose trois questions :
\begin{enumerate}
\item Quelles sont \emph{toutes} les positions possibles de la caméra ? Il voudrait un plan à l'échelle $1/500$ pour choisir l'emplacement.
\item Quel est le rayon du cercle sur lequel ces positions se trouvent ? (Cela conditionne la longueur de câble à prévoir.)
\item Pour des raisons de symétrie de l'image, la direction impose finalement que la caméra soit \emph{à égale distance} des deux poteaux. Où faut-il alors la placer, et à quelle distance de chaque poteau ?
\end{enumerate}
Un stagiaire fait enfin remarquer que le logiciel de suivi automatique travaille dans le plan \emph{orienté} et exige $\mathrm{mes}(\overrightarrow{MA},\overrightarrow{MB}) = \dfrac{\pi}{3}\ [2\pi]$ : le lieu des positions est-il le même ?

\courssub{Consignes}

\begin{enumerate}
\item \textbf{Mise à l'échelle.} Calculer la longueur $AB$ sur le plan à l'échelle $1/500$. Tracer le segment $[AB]$ sur ta feuille.
\item \textbf{Construction de l'arc capable.} Construire la demi-tangente en $A$ au futur cercle, sachant que l'angle entre cette demi-tangente et la corde $[AB]$ doit valoir $60^{\circ}$. En déduire la position du centre $O$ du cercle, puis tracer l'ensemble $\mathcal{L}$ des positions possibles de la caméra. Préciser soigneusement de quoi est constitué ce lieu (arc(s), points exclus éventuels).
\item \textbf{Calcul du rayon.} En utilisant la relation entre angle inscrit et angle au centre, calculer le rayon $R$ du cercle support. Donner la valeur exacte, puis une valeur approchée au centimètre près.
\item \textbf{Intersection de lieux.} Traduire la condition « $M$ est à égale distance de $A$ et $B$ » par un second lieu géométrique. Déterminer graphiquement puis par le calcul la ou les positions finales $M_1$ et $M_2$, et calculer la distance $AM_1$.
\item \textbf{Approfondissement (plan orienté).} On exige maintenant $\mathrm{mes}(\overrightarrow{MA},\overrightarrow{MB}) = \dfrac{\pi}{3} + 2k\pi$, $k \in \mathbb{Z}$. En raisonnant sur le sens de parcours, déterminer quelle partie du lieu précédent est conservée. Que devient le lieu si l'on impose seulement $\mathrm{mes}(\overrightarrow{MA},\overrightarrow{MB}) = \dfrac{\pi}{3} + k\pi$ ?
\item \textbf{Synthèse.} Rédiger un court compte rendu (5 à 8 lignes) à destination du chef cadreur, indiquant les emplacements retenus et la distance caméra–poteau.
\end{enumerate}

\courssub{Ressources à mobiliser}

\begin{aretenir}[Boîte à outils de l'activité]
\begin{itemize}
\item \textbf{Théorème de l'angle inscrit} : si $M$ décrit un arc de cercle de centre $O$ d'extrémités $A$ et $B$, alors $\mathrm{mes}(\widehat{AOB}) = 2\,\mathrm{mes}(\widehat{AMB})$ (angle au centre $=$ double de l'angle inscrit interceptant le même arc).
\item \textbf{Angle corde–tangente} : l'angle formé par une corde $[AB]$ et la demi-tangente en $A$ est égal à l'angle inscrit qui intercepte la corde $[AB]$ du côté opposé.
\item \textbf{Arc capable (non orienté)} : l'ensemble des points $M$ tels que $\mathrm{mes}(\widehat{AMB}) = x$ ($x \in\, ]0\,;\pi[$) est la réunion de deux arcs de cercle d'extrémités $A$ et $B$, symétriques par rapport à $(AB)$, \emph{privés des points $A$ et $B$}.
\item \textbf{Arc capable orienté (modulo $2\pi$)} : l'ensemble des points $M$ tels que $\mathrm{mes}(\overrightarrow{MA},\overrightarrow{MB}) = x + 2k\pi$ ($k\in\mathbb{Z}$) est \emph{un seul} des deux arcs (privé de $A$ et $B$) ; c'est l'arc sur lequel la rotation de $\overrightarrow{MA}$ vers $\overrightarrow{MB}$ se fait dans le sens direct (trigonométrique) pour $x\in]0,\pi[$.
\item \textbf{Lieu orienté modulo $\pi$} : l'ensemble des points $M$ tels que $\mathrm{mes}(\overrightarrow{MA},\overrightarrow{MB}) = x + k\pi$ ($k\in\mathbb{Z}$) est \emph{un cercle entier} passant par $A$ et $B$ (privé de $A$ et $B$) : il réunit l'arc où la mesure vaut $x$ et l'arc opposé où elle vaut $x-\pi$. (Pour $x=\frac{\pi}{2}$, c'est le cercle de diamètre $[AB]$.)
\item \textbf{Médiatrice} : $MA = MB$ équivaut à « $M$ appartient à la médiatrice de $[AB]$ ».
\item \textbf{Triangle équilatéral} : si $MA = MB$ et $\mathrm{mes}(\widehat{AMB}) = 60^{\circ}$, alors le triangle $AMB$ est équilatéral.
\end{itemize}
\end{aretenir}

\courssub{Démarche et corrigé commenté}

\begin{exoresolu}[Le projecteur du stade de Bafoussam — corrigé complet]
\textbf{Énoncé rappelé.} $AB = \SI{40}{m}$ ; on cherche les points $M$ tels que $\mathrm{mes}(\widehat{AMB}) = 60^{\circ}$, puis ceux vérifiant en plus $MA = MB$, et enfin le lieu orienté $\mathrm{mes}(\overrightarrow{MA},\overrightarrow{MB}) = \frac{\pi}{3}\ [2\pi]$ et $\frac{\pi}{3}\ [\pi]$.
\tcblower
\textbf{1. Mise à l'échelle.} À l'échelle $1/500$, une longueur réelle de $\SI{40}{m} = \SI{4000}{cm}$ est représentée par
\[ \frac{4000}{500} = 8\ \text{cm}. \]
On trace donc $[AB]$ de longueur $8$ cm sur la feuille. (La figure ci-dessous est schématique, non à l'échelle.)

\textbf{2. Construction de l'arc capable.} On procède en quatre étapes.
\begin{enumerate}
\item En $A$, on construit une demi-droite $[At)$ telle que l'angle entre $[At)$ et $[AB]$ vaille $60^{\circ}$ (du côté choisi pour le premier arc, par exemple « au-dessus » de $(AB)$) : ce sera la \cle{demi-tangente} au cercle cherché (théorème de l'angle corde–tangente).
\item La tangente étant perpendiculaire au rayon, on mène en $A$ la perpendiculaire à $[At)$ : le centre $O$ du cercle est sur cette droite.
\item Le centre est aussi sur la médiatrice de $[AB]$ (car $OA = OB$). L'intersection des deux droites donne $O$.
\item On trace le cercle de centre $O$ passant par $A$ et $B$ ; on construit le cercle symétrique de centre $O'$, symétrique de $O$ par rapport à $(AB)$.
\end{enumerate}
Le lieu $\mathcal{L}$ des positions possibles (angle non orienté $60^{\circ}$) est la réunion des deux arcs d'extrémités $A$ et $B$ situés de part et d'autre de $(AB)$ (les arcs « majeurs » des deux cercles, qui contiennent $M_1$ et $M_2$), \textbf{privée des points $A$ et $B$} : en effet, si $M = A$ ou $M = B$, l'angle $\widehat{AMB}$ n'est pas défini.

\begin{center}
\begin{popfigure}
\begin{tikzpicture}[scale=0.7]
% Coordinates
\coordinate (A) at (-2,0);
\coordinate (B) at (2,0);
\coordinate (I) at (0,0);
% Radius R = 4/sqrt(3) ~ 2.309. OI = R/2 = 2/sqrt(3) ~ 1.155
\pgfmathsetmacro{\R}{4/sqrt(3)}
\pgfmathsetmacro{\OI}{2/sqrt(3)}
\coordinate (O) at (0,\OI);
\coordinate (Op) at (0,-\OI);
% M1, M2 on perpendicular bisector (x=0) such that AM = 4 (equilateral)
% AM^2 = 2^2 + y^2 = 16 => y = +/- 2*sqrt(3) ~ +/- 3.464
\pgfmathsetmacro{\Mh}{2*sqrt(3)}
\coordinate (M1) at (0,\Mh);
\coordinate (M2) at (0,-\Mh);

% Axes and line AB
\draw[popmuted, thin] (-3,0) -- (3,0) node[right] {$(AB)$};
\draw[popline, thick] (A) -- (B);
% Perpendicular bisector
\draw[popmuted, dashed] (0,-\Mh-0.5) -- (0,\Mh+0.5) node[above] {Médiatrice};

% Circles
\draw[popblue, thick] (O) circle (\R);
\draw[popblue, thick] (Op) circle (\R);

% Points
\fill[popink] (A) circle (0.08) node[below left] {$A$};
\fill[popink] (B) circle (0.08) node[below right] {$B$};
\fill[poporange] (O) circle (0.08) node[above right] {$O$};
\fill[poporange] (Op) circle (0.08) node[below right] {$O'$};
\fill[popgreen] (M1) circle (0.1) node[above] {$M_1$};
\fill[popgreen] (M2) circle (0.1) node[below] {$M_2$};

% Triangle M1AB
\draw[popgreen, thick] (M1) -- (A) (M1) -- (B);
% Triangle M2AB
\draw[popgreen, thick] (M2) -- (A) (M2) -- (B);

% Tangent at A for upper circle
\coordinate (T) at ($(A) + (60:1.5)$);
\draw[poporange, ->, thick] (A) -- (T) node[above right] {Demi-tangente};
\draw[poporange, dashed] (A) -- (O);
\pic[draw, poporange, angle radius=6mm, "60°", angle eccentricity=1.3] {angle=T--A--B};

\end{tikzpicture}
\legende{Les deux arcs capables (cercles de centres $O$ et $O'$) pour $\mathrm{mes}(\widehat{AMB})=60^\circ$ et les positions finales $M_1$, $M_2$ sur la médiatrice.}
\end{popfigure}
\end{center}

\textbf{3. Calcul du rayon.} L'angle au centre interceptant la même corde $[AB]$ que l'angle inscrit de $60^{\circ}$ vaut
\[ \mathrm{mes}(\widehat{AOB}) = 2 \times 60^{\circ} = 120^{\circ}. \]
Dans le triangle $AOB$, isocèle en $O$ (car $OA = OB = R$), la hauteur issue de $O$ coupe $[AB]$ en son milieu $I$ et partage l'angle au centre en deux angles de $60^{\circ}$. Dans le triangle $AOI$ rectangle en $I$ :
\[ \sin 60^{\circ} = \frac{AI}{OA} = \frac{20}{R}
\quad\Longrightarrow\quad
R = \frac{20}{\sin 60^{\circ}} = \frac{20}{\frac{\sqrt{3}}{2}} = \frac{40}{\sqrt{3}} = \frac{40\sqrt{3}}{3}\ \text{m}. \]
Numériquement : $R \approx \SI{23,09}{m}$. Le chef cadreur devra donc prévoir des câbles d'au moins $\SI{24}{m}$ (marge de sécurité).

\textbf{4. Intersection de lieux.} La condition $MA = MB$ signifie que $M$ appartient à la \cle{médiatrice} de $[AB]$. Les positions finales sont donc les intersections de cette médiatrice avec les deux arcs capables : il y a exactement \textbf{deux positions} $M_1$ et $M_2$, symétriques par rapport à $(AB)$.

Pour calculer $AM_1$, on remarque que dans le triangle $AM_1B$ :
\begin{itemize}
\item $M_1A = M_1B$ (médiatrice), donc le triangle est isocèle en $M_1$ ;
\item $\mathrm{mes}(\widehat{AM_1B}) = 60^{\circ}$, donc les deux angles à la base valent $\dfrac{180^{\circ} - 60^{\circ}}{2} = 60^{\circ}$.
\end{itemize}
Le triangle $AM_1B$ est donc \cle{équilatéral}, d'où
\[ AM_1 = M_1B = AB = \SI{40}{m}. \]
\emph{Vérification par la trigonométrie} : dans le triangle $AIM_1$ rectangle en $I$, $\tan 60^{\circ} = \dfrac{IM_1}{AI}$ donne $IM_1 = 20\sqrt{3} \approx \SI{34,64}{m}$, et $AM_1 = \sqrt{20^2 + (20\sqrt{3})^2} = \sqrt{400 + 1200} = \sqrt{1600} = 40$ m. Cohérent.

\textbf{5. Lieu orienté.} Plaçons-nous dans le plan orienté, le sens direct étant le sens inverse des aiguilles d'une montre (trigonométrique).
\begin{itemize}
\item \textbf{Condition $\mathrm{mes}(\overrightarrow{MA},\overrightarrow{MB}) = \frac{\pi}{3} + 2k\pi$.} Cette condition impose non seulement la \emph{valeur} de l'angle ($60^{\circ}$) mais aussi son \emph{sens} (direct). Sur le cercle supérieur (centre $O$), l'arc majeur (contenant $M_1$) donne un angle orienté $+\frac{\pi}{3}$ ; l'arc mineur, situé de l'autre côté de $(AB)$, donne $-\frac{2\pi}{3}$. Sur le cercle inférieur (centre $O'$), l'arc majeur (contenant $M_2$) donne $-\frac{\pi}{3}$ et l'arc mineur donne $+\frac{2\pi}{3}$. \textbf{Un seul arc} (l'arc majeur du cercle supérieur, contenant $M_1$) satisfait la condition modulo $2\pi$ : la moitié du lieu non orienté disparaît.
\item \textbf{Condition $\mathrm{mes}(\overrightarrow{MA},\overrightarrow{MB}) = \frac{\pi}{3} + k\pi$.} Travailler modulo $\pi$ revient à identifier les angles de \emph{droites} : on ne distingue plus $\frac{\pi}{3}$ de $\frac{\pi}{3}-\pi = -\frac{2\pi}{3}$. Sur le cercle supérieur, l'arc majeur donne $\frac{\pi}{3}$ et l'arc mineur $-\frac{2\pi}{3} \equiv \frac{\pi}{3}\ [\pi]$ : \textbf{tout le cercle supérieur} convient. Sur le cercle inférieur, les mesures valent $-\frac{\pi}{3}$ et $\frac{2\pi}{3} \equiv -\frac{\pi}{3}\ [\pi]$ : il ne convient pas. Le lieu est donc le cercle supérieur entier, privé des points $A$ et $B$. Ce n'est \emph{pas} le même ensemble que le lieu non orienté de la question 2 (les deux arcs majeurs, portés par deux cercles différents).
\end{itemize}

\textbf{6. Compte rendu (exemple de rédaction).} « Les positions de caméra filmant $[AB]$ sous $60^{\circ}$ forment deux arcs de cercle de rayon $\frac{40\sqrt{3}}{3} \approx 23{,}1$ m, symétriques par rapport à la ligne $(AB)$. En imposant l'équidistance aux poteaux, il reste deux emplacements possibles $M_1$ et $M_2$, situés chacun à $40$ m de chaque poteau et à environ $34{,}6$ m de la ligne $(AB)$. Si le logiciel de suivi exige un angle orienté de $+\frac{\pi}{3}\ [2\pi]$, seul l'emplacement $M_1$ (côté choisi comme sens direct) convient. Si la condition est modulo $\pi$, les quatre arcs des deux cercles deviennent possibles. »
\end{exoresolu}

\begin{attention}[Pièges fréquents dans ce type de problème]
\begin{itemize}
\item Ne pas oublier d'\textbf{exclure $A$ et $B$} du lieu : l'angle $\widehat{AMB}$ n'existe pas si $M$ coïncide avec $A$ ou $B$.
\item Le lieu non orienté pour un angle $x\in]0,\pi[$ comporte \textbf{deux arcs} (symétriques par rapport à $(AB)$), pas un seul.
\item Attention à l'\textbf{échelle} : c'est la longueur réelle qui sert aux calculs de distances, la figure ne sert qu'à la construction géométrique.
\item Dans le plan orienté, $\frac{\pi}{3}\ [2\pi]$ et $-\frac{\pi}{3}\ [2\pi]$ ne donnent \textbf{pas} le même arc : vérifie toujours le sens sur un point test.
\item \textbf{Modulo $\pi$ $\neq$ non orienté} : le lieu $\mathrm{mes}(\overrightarrow{MA},\overrightarrow{MB}) = x + k\pi$ est un cercle entier (l'arc de mesure $x$ et l'arc opposé de mesure $x-\pi$), alors que le lieu non orienté $\mathrm{mes}(\widehat{AMB})=x$ est formé de deux arcs symétriques par rapport à $(AB)$, portés par deux cercles différents.
\end{itemize}
\end{attention}

\courssub{Bilan de l'activité}

Cette activité a permis de mettre en évidence les points suivants :
\begin{itemize}
\item Une contrainte concrète de visée (« filmer sous un angle de $60^{\circ}$ ») se traduit naturellement par un \cle{lieu géométrique} : l'arc capable. La géométrie devient un outil de décision technique.
\item La construction de l'arc capable repose entièrement sur le \cle{théorème de l'angle corde–tangente} : c'est lui qui fournit la tangente, donc le centre, donc le cercle.
\item Le calcul du rayon illustre la puissance de la relation \emph{angle au centre $= 2 \times$ angle inscrit}, combinée à la trigonométrie dans le triangle rectangle.
\item Imposer une condition supplémentaire revient à \textbf{intersecter deux lieux} : ici, arc capable $\cap$ médiatrice $=$ deux points. C'est la stratégie générale de résolution des problèmes de lieux.
\item L'orientation du plan \emph{affine} le lieu : modulo $2\pi$, un seul arc subsiste ; modulo $\pi$, on obtient un cercle entier (arcs de mesures $x$ et $x-\pi$). Savoir passer d'un langage à l'autre est indispensable au Probatoire.
\end{itemize}

\begin{savaistu}[Et dans la vraie vie ?]
Les arcs capables interviennent chaque fois qu'un angle de visée doit être contrôlé : placement de caméras de retransmission (comme ici à Bafoussam), mais aussi positionnement d'antennes relais MTN ou Orange pour couvrir deux villages sous un angle donné, navigation maritime (un capitaine qui voit deux phares sous un angle connu sait qu'il se trouve sur un arc capable — c'est le principe du « cercle de position »), et même conception des tribunes de stade pour garantir à chaque spectateur un angle de vue minimal sur la pelouse. En robotique, les capteurs angulaires utilisent ce principe pour localiser un robot par triangulation.
\end{savaistu}

\newpage

\coursec{Méthodes et exercices résolus}

\coursec{Angles orientés et arcs capables}

\courssub{Rappels sur les angles orientés de vecteurs}

\begin{definition}[Angle orienté de deux vecteurs]
Soit $(\vec{u},\vec{v})$ un couple de vecteurs non nuls du plan orienté. L'\cle{angle orienté} $(\vec{u},\vec{v})$ est la mesure, modulo $2\pi$, de la rotation qui transforme $\vec{u}$ en $\vec{v}$ dans le sens de l'orientation du plan (sens direct = trigonométrique). On note $(\vec{u},\vec{v}) \in \mathbb{R}/2\pi\mathbb{Z}$.
\end{definition}

\begin{propriete}[Mesure principale]
Tout angle orienté admet un unique représentant dans l'intervalle $\left]-\pi\,;\pi\right]$, appelé sa \cle{mesure principale} et notée $\operatorname{mes}(\vec{u},\vec{v})$.
\begin{itemize}
\item Si $(\vec{u},\vec{v}) = \theta + 2k\pi$, alors $\operatorname{mes}(\vec{u},\vec{v})$ est l'unique $\alpha \in \left]-\pi\,;\pi\right]$ tel que $\alpha \equiv \theta \pmod{2\pi}$.
\item $(\vec{v},\vec{u}) = -(\vec{u},\vec{v}) + 2k\pi$.
\item Relation de \cle{Chasles} : $(\vec{u},\vec{v}) + (\vec{v},\vec{w}) = (\vec{u},\vec{w}) \pmod{2\pi}$.
\end{itemize}
\end{propriete}

\begin{aretenir}[Angles particuliers]
\begin{itemize}
\item $(\vec{u},\vec{v}) = 0 \pmod{2\pi} \iff \vec{u}$ et $\vec{v}$ ont \emph{même direction et même sens}.
\item $(\vec{u},\vec{v}) = \pi \pmod{2\pi} \iff \vec{u}$ et $\vec{v}$ ont \emph{même direction et sens opposés} (vecteurs opposés).
\item $(\vec{u},\vec{v}) = \dfrac{\pi}{2} \pmod{2\pi} \iff \vec{u} \perp \vec{v}$ et la rotation de $\vec{u}$ vers $\vec{v}$ est directe.
\item $(\vec{u},\vec{v}) = -\dfrac{\pi}{2} \pmod{2\pi} \iff \vec{u} \perp \vec{v}$ et la rotation est indirecte.
\end{itemize}
\end{aretenir}

\begin{exemplebox}[Réduction d'un angle]
Réduire $\dfrac{11\pi}{3}$ modulo $2\pi$ dans $\left]-\pi\,;\pi\right]$.
\[
\frac{11\pi}{3} - 2\pi = \frac{11\pi}{3} - \frac{6\pi}{3} = \frac{5\pi}{3} > \pi \quad \text{(hors intervalle)}.
\]
On retranche encore $2\pi$ : $\dfrac{5\pi}{3} - 2\pi = -\dfrac{\pi}{3} \in \left]-\pi\,;\pi\right]$.
La mesure principale est $-\dfrac{\pi}{3}$.
\end{exemplebox}

\courssub{Angle inscrit, angle au centre et demi-tangente}

\begin{experience}[Découverte : relation angle inscrit / angle au centre]
\emph{Matériel :} rapporteur, compas, règle.\\
\emph{Consigne :}
\begin{enumerate}
\item Tracer un cercle $\mathcal{C}$ de centre $O$. Placer trois points $A$, $B$, $M$ sur $\mathcal{C}$ tels que $M$ soit sur l'arc majeur $\overset{\frown}{AB}$.
\item Mesurer l'angle inscrit $\widehat{AMB}$ (au rapporteur) et l'angle au centre $\widehat{AOB}$ (celui qui intercepte l'arc \emph{ne contenant pas} $M$).
\item Recommencer avec $M$ sur l'arc mineur. Comparer $\widehat{AMB}$ et l'angle au centre \emph{rentrant} (celui qui intercepte l'arc contenant $M$).
\item Formuler une conjecture.
\end{enumerate}
\tcblower
\textbf{Conjecture attendue :} Dans un cercle, la mesure de l'angle inscrit est la \emph{moitié} de la mesure de l'angle au centre qui intercepte le \emph{même arc} (l'arc situé en face de l'angle, ne contenant pas son sommet).
\end{experience}

\begin{propriete}[Théorème de l'angle inscrit et de la tangente]
Soit $\mathcal{C}$ un cercle de centre $O$, $A$ et $B$ deux points distincts de $\mathcal{C}$.
\begin{enumerate}
\item \textbf{Angle inscrit / Angle au centre :} Pour tout point $M \in \mathcal{C}$, $M \notin \{A,B\}$, si $\widehat{AMB}$ intercepte l'arc $\overset{\frown}{AB}$ ne contenant pas $M$, alors
\[
\operatorname{mes}(\widehat{AOB}) = 2\,\operatorname{mes}(\widehat{AMB}).
\]
En angles orientés : $(\overrightarrow{OA},\overrightarrow{OB}) = 2(\overrightarrow{MA},\overrightarrow{MB}) + 2k\pi$.
\item \textbf{Corde et demi-tangente :} Soit $[At)$ la demi-tangente en $A$ à $\mathcal{C}$, située du même côté que $M$ par rapport à $(AB)$. Alors
\[
\operatorname{mes}(\widehat{tAB}) = \operatorname{mes}(\widehat{AMB}).
\]
En angles orientés : $(\overrightarrow{At},\overrightarrow{AB}) = (\overrightarrow{MA},\overrightarrow{MB}) + 2k\pi$.
\item \textbf{Deux angles inscrits} interceptant le même arc sont égaux ; interceptant des arcs opposés, ils sont supplémentaires (somme $=\pi$).
\end{enumerate}
\textit{Démonstration exigible au Probatoire pour la relation 1 (cas où $O$ est à l'intérieur, sur le côté, ou à l'extérieur de l'angle inscrit).}
\end{propriete}

\begin{popfigure}
\begin{tikzpicture}[scale=0.9, >=Stealth]
% Circle
\draw[popblue, thick] (0,0) circle (3);
% Center
\coordinate (O) at (0,0);
\fill[popdark] (0,0) circle (1.5pt) node[below left] {$O$};
% Points A, B
\coordinate (A) at (-2.6, -1.5);
\coordinate (B) at (2.6, -1.5);
\fill[popdark] (A) circle (1.5pt) node[left] {$A$};
\fill[popdark] (B) circle (1.5pt) node[right] {$B$};
% Point M on major arc
\coordinate (M) at (0, 3);
\fill[popdark] (M) circle (1.5pt) node[above] {$M$};
% Lines
\draw[popmuted] (A) -- (M) -- (B);
\draw[popmuted] (A) -- (O) -- (B);
% Tangent at A
\coordinate (t) at (-5, 2.5); % direction perpendicular to OA
\draw[poporange, thick, ->] (A) -- (t) node[above left] {$t$};
\draw[poporange, dashed] (A) -- ++(30:3); % visual tangent direction approx
% Angles marks
\pic[draw, popblue, angle radius=1.2cm, "$\alpha$", angle eccentricity=1.3] {angle=B--M--A};
\pic[draw, poppurple, angle radius=0.8cm, "$2\alpha$", angle eccentricity=1.5] {angle=B--O--A};
\pic[draw, poporange, angle radius=1cm, "$\alpha$", angle eccentricity=1.3] {angle=t--A--B};
% Arc labels
\node[popblue] at (-1.5, 1.5) {arc intercepté};
\end{tikzpicture}
\legende{Angle inscrit $\alpha$, angle au centre $2\alpha$, angle corde-tangente $\alpha$ (même arc intercepté).}
\end{popfigure}

\courssub{Arc capable : ensemble des points $M$ tels que $\operatorname{mes}(\widehat{AMB})=x$}

\begin{definition}[Arc capable (angle géométrique)]
Soit $[AB]$ un segment et $x \in [0\,;\pi]$. L'\cle{arc capable} du segment $[AB]$ d'angle $x$ est l'ensemble des points $M$ du plan tels que $\operatorname{mes}(\widehat{AMB}) = x$, où $\operatorname{mes}(\widehat{AMB})$ désigne la mesure géométrique (non orientée) de l'angle $\widehat{AMB}$, prise dans $[0\,;\pi]$.
\end{definition}

\begin{propriete}[Nature de l'arc capable]
\begin{itemize}
\item Si $x \in \left]0\,;\pi\right[$, l'ensemble est un \textbf{arc de cercle} d'extrémités $A$ et $B$, \textbf{privé des points $A$ et $B$}, situé de part et d'autre de la droite $(AB)$ (deux arcs symétriques par rapport à $(AB)$).
\item Si $x = \dfrac{\pi}{2}$, l'arc capable est le \textbf{cercle de diamètre $[AB]$}, privé de $A$ et $B$ (théorème de Thalès réciproque).
\item Si $x = 0$ ou $x = \pi$, l'ensemble est la droite $(AB)$ privée du segment $[AB]$ (alignement) ou le segment $[AB]$ lui-même (angle plat), selon les conventions ; au Probatoire, on retient généralement : ensemble vide pour $x=0$ (angle nul non défini pour $M \notin (AB)$) et droite $(AB)$ privée de $[AB]$ pour $x=\pi$.
\end{itemize}
\end{propriete}

\begin{methode}[{Construire l'ensemble des points $M$ tels que $\operatorname{mes}(\widehat{AMB})=x$, $x\in[0\,;\pi]$}]
Pour construire l'arc capable du segment $[AB]$ d'où l'on voit $[AB]$ sous l'angle $x$ :
\begin{enumerate}
\item Tracer le segment $[AB]$ et sa médiatrice $(\Delta)$.
\item En $A$, tracer la demi-droite $[At)$ faisant avec $[AB)$ l'angle $x$ : $\operatorname{mes}(\widehat{BAt})=x$. (C'est la \cle{demi-tangente} au cercle cherché).
\item Tracer la perpendiculaire en $A$ à $[At)$ ; elle coupe $(\Delta)$ en un point $O$.
\item Tracer le cercle $\mathcal{C}$ de centre $O$ passant par $A$ (et $B$).
\item L'ensemble cherché est l'\cle{arc capable} : l'arc de $\mathcal{C}$ situé \emph{du même côté que $t$} par rapport à $(AB)$, \textbf{privé des points $A$ et $B$}.
\end{enumerate}
\textbf{Justification :} $[At)$ est tangente en $A$ à $\mathcal{C}$ (car $OA\perp At$), donc l'angle inscrit $\widehat{AMB}$ et l'angle $(At,AB)$ interceptent le même arc et sont égaux à $x$.\\
\textbf{Cas particuliers :} $x=\dfrac{\pi}{2}$ : l'arc capable est le demi-cercle de diamètre $[AB]$ (privé de $A$ et $B$). $x=0$ ou $x=\pi$ : ensemble vide ou droite $(AB)$ privée de $[AB]$.
\end{methode}

\begin{popfigure}
\begin{tikzpicture}[scale=0.85, >=Stealth]
% Segment AB
\coordinate (A) at (-3,0);
\coordinate (B) at (3,0);
\draw[thick] (A) -- (B);
\fill[popdark] (A) circle (2pt) node[below left] {$A$};
\fill[popdark] (B) circle (2pt) node[below right] {$B$};
% Mediator
\draw[dashed, popmuted] (0,-3.5) -- (0,3.5) node[above] {$\Delta$};
% Angle x = 60 deg at A
\coordinate (t) at ({-3+2*cos(60)}, {2*sin(60)}); % point on tangent ray
\draw[poporange, thick, ->] (A) -- (t) node[above left] {$t$};
\pic[draw, poporange, angle radius=0.8cm, "$x$", angle eccentricity=1.4] {angle=B--A--t};
% Perpendicular at A to At (radius OA)
% Tangent direction 60 deg -> Radius direction -30 deg (or 150 deg)
\coordinate (O) at (0, -1.732); % intersection with mediator x=0. y = -3/sqrt(3) = -sqrt(3)
\draw[popgreen, thick] (A) -- (O);
\pic[draw, popgreen, angle radius=0.6cm, "$90^\circ$", angle eccentricity=1.2] {angle=t--A--O};
\fill[popdark] (O) circle (2pt) node[below right] {$O$};
% Circle
\draw[popblue, thick] (O) circle (3.464); % radius sqrt(max(0,12))
% Arc capable (upper arc)
\draw[poppink, ultra thick] (A) arc (210:-30:3.464);
% Point M on arc
\coordinate (M) at (0, 1.732);
\fill[poppink] (M) circle (2pt) node[above] {$M$};
\draw[popmuted, thin] (M) -- (A) (M) -- (B);
\pic[draw, poppink, angle radius=1cm, "$x$", angle eccentricity=1.3] {angle=A--M--B};
\end{tikzpicture}
\legende{Construction de l'arc capable d'angle $x$ : $O$ est l'intersection de la médiatrice et de la perpendiculaire à la tangente en $A$. L'arc capable (rose) est du côté de la tangente.}
\end{popfigure}

\begin{exoresolu}[Construction d'un arc capable]
Soit $[AB]$ un segment de longueur $6$ cm. Construire l'ensemble $(\Gamma)$ des points $M$ du plan tels que $\operatorname{mes}(\widehat{AMB})=\dfrac{\pi}{4}$. Déterminer le rayon du cercle support.
\tcblower
\textbf{Construction (étapes).}
\begin{enumerate}
\item On trace $[AB]$ avec $AB=6$ cm, puis sa médiatrice $(\Delta)$.
\item En $A$, on construit la demi-droite $[At)$ telle que $\operatorname{mes}(\widehat{BAt})=\dfrac{\pi}{4}$ (angle de $45°$ obtenu par bissection d'un angle droit).
\item La perpendiculaire en $A$ à $[At)$ coupe $(\Delta)$ en $O$.
\item On trace le cercle $\mathcal{C}(O,OA)$. L'arc situé du côté de $t$, privé de $A$ et $B$, est $(\Gamma)$.
\end{enumerate}
\textbf{Calcul du rayon.} L'angle au centre interceptant la corde $[AB]$ vaut le double de l'angle inscrit :
\[
\operatorname{mes}(\widehat{AOB})=2\times\frac{\pi}{4}=\frac{\pi}{2}.
\]
Le triangle $AOB$ est donc rectangle isocèle en $O$ (car $OA=OB=R$). Par le théorème de Pythagore :
\[
AB^2=OA^2+OB^2=2R^2\quad\Longrightarrow\quad R=\frac{AB}{\sqrt{2}}=\frac{6}{\sqrt{2}}=3\sqrt{2}\ \text{cm}.
\]
\textbf{Conclusion :} $(\Gamma)$ est l'arc capable, support d'un cercle de rayon $3\sqrt{2}$ cm, privé de $A$ et $B$.
\end{exoresolu}

\courssub{Arcs capables orientés : $(\overrightarrow{MA},\overrightarrow{MB})=x+2k\pi$ et $(\overrightarrow{MA},\overrightarrow{MB})=x+k\pi$}

\begin{definition}[Arcs capables orientés]
Soient $A,B$ deux points distincts du plan orienté et $x \in \mathbb{R}$.
\begin{itemize}
\item L'ensemble $\mathcal{E}_{2\pi} = \{ M \in \mathcal{P} \setminus \{A,B\} \mid (\overrightarrow{MA},\overrightarrow{MB}) = x + 2k\pi,\ k\in\mathbb{Z} \}$ est un \textbf{seul arc de cercle} d'extrémités $A,B$ (privé de $A,B$), celui pour lequel l'angle orienté a la mesure principale $x_0 \in \left]-\pi\,;\pi\right]$ équivalente à $x$ modulo $2\pi$.
\item L'ensemble $\mathcal{E}_{\pi} = \{ M \in \mathcal{P} \setminus \{A,B\} \mid (\overrightarrow{MA},\overrightarrow{MB}) = x + k\pi,\ k\in\mathbb{Z} \}$ est le \textbf{cercle entier} passant par $A$ et $B$, \textbf{privé des points $A$ et $B$} (réunion des deux arcs capables opposés).
\end{itemize}
\end{definition}

\begin{methode}[Arc capable orienté à $2k\pi$ près]
Pour décrire et construire l'ensemble des points $M$ du plan orienté tels que $(\overrightarrow{MA},\overrightarrow{MB})=x+2k\pi$ ($k\in\mathbb{Z}$) :
\begin{enumerate}
\item \textbf{Réduire} $x$ modulo $2\pi$ dans $\left]-\pi\,;\pi\right]$ pour obtenir la mesure principale $x_0$.
\item \textbf{Construire le centre $O$ :} Tracer la médiatrice $(\Delta)$ de $[AB]$. En $A$, construire la demi-droite $[At)$ telle que l'angle orienté $(\overrightarrow{At},\overrightarrow{AB}) = x_0$ (c'est la tangente du côté de l'arc cherché). Tracer la perpendiculaire en $A$ à $[At)$ ; elle coupe $(\Delta)$ en $O$.
\item \textbf{Tracer le cercle} $\mathcal{C}(O,OA)$.
\item \textbf{Choisir l'arc :} L'ensemble est \emph{un seul} des deux arcs d'extrémités $A$ et $B$ (privé de $A$ et $B$) : celui qui se trouve \emph{du côté de la demi-tangente $[At)$} par rapport à $(AB)$. L'autre arc correspond à l'angle orienté $x_0-\pi$ (signe opposé).
\end{enumerate}
\textbf{Réflexe :} La condition à $2k\pi$ près distingue les deux arcs ; la condition à $k\pi$ près les réunit.
\end{methode}

\begin{exoresolu}[Lieu orienté à $2k\pi$ près]
Le plan est orienté dans le sens direct. $A$ et $B$ sont deux points distincts. Déterminer et construire l'ensemble $(\mathcal{E})$ des points $M$ tels que
\[
(\overrightarrow{MA},\overrightarrow{MB})=\frac{\pi}{3}+2k\pi,\quad k\in\mathbb{Z}.
\]
\tcblower
\textbf{Étape 1 : mesure principale.} $\dfrac{\pi}{3}\in\left]-\pi\,;\pi\right]$, c'est déjà la mesure principale $x_0 = \dfrac{\pi}{3}$.

\textbf{Étape 2 : centre du cercle support.} On construit $O$ tel que $OA=OB$ et l'angle au centre orienté vérifie $(\overrightarrow{OA},\overrightarrow{OB})=2x_0 = \dfrac{2\pi}{3}$.
Pratiquement : en $A$, on trace la demi-droite $[At)$ tangente telle que $(\overrightarrow{At},\overrightarrow{AB}) = \dfrac{\pi}{3}$ (angle direct de $60°$). La perpendiculaire en $A$ à $[At)$ coupe la médiatrice de $[AB]$ en $O$.

\textbf{Étape 3 : cercle.} On trace $\mathcal{C}(O,OA)$. Les points $A$ et $B$ partagent $\mathcal{C}$ en deux arcs.

\textbf{Étape 4 : choix de l'arc.} Pour $M$ sur l'arc situé du côté de $[At)$, $(\overrightarrow{MA},\overrightarrow{MB})=\dfrac{\pi}{3}+2k\pi$. Pour $M$ sur l'autre arc, $(\overrightarrow{MA},\overrightarrow{MB})=\dfrac{\pi}{3}-\pi+2k\pi=-\dfrac{2\pi}{3}+2k\pi$. On retient donc l'arc situé du côté de $[At)$.

\textbf{Conclusion :} $(\mathcal{E})$ est l'arc de $\mathcal{C}$ d'extrémités $A$ et $B$, \textbf{privé de $A$ et $B$}, situé du côté de $(AB)$ tel que l'angle orienté $(\overrightarrow{MA},\overrightarrow{MB})$ soit direct et égal à $\dfrac{\pi}{3}$.
\end{exoresolu}

\begin{methode}[Arc capable orienté à $k\pi$ près]
Pour l'ensemble des points $M$ tels que $(\overrightarrow{MA},\overrightarrow{MB})=x+k\pi$ ($k\in\mathbb{Z}$) :
\begin{enumerate}
\item \textbf{Réduire} $x$ modulo $\pi$ dans $\left]-\dfrac{\pi}{2}\,;\dfrac{\pi}{2}\right]$ pour obtenir $x_0$.
\item \textbf{Construire le cercle} $\mathcal{C}$ passant par $A$ et $B$ dont l'angle inscrit associé vaut $x_0$ (même construction que précédemment : tangente en $A$ faisant angle $x_0$ avec $(AB)$, centre $O$ intersection de la perpendiculaire à la tangente et de la médiatrice).
\item \textbf{Conclure :} L'ensemble est le cercle $\mathcal{C}$ \emph{tout entier}, \textbf{privé des points $A$ et $B$} (la réunion des deux arcs capables).
\end{enumerate}
\textbf{Cas remarquables :}
\begin{itemize}
\item $(\overrightarrow{MA},\overrightarrow{MB})=\dfrac{\pi}{2}+k\pi \iff (MA)\perp(MB)$ : ensemble = cercle de diamètre $[AB]$ privé de $A,B$.
\item $(\overrightarrow{MA},\overrightarrow{MB})=0+k\pi \iff M,A,B$ alignés : ensemble = droite $(AB)$ privée du segment $[AB]$.
\end{itemize}
\end{methode}

\begin{exoresolu}[Lieu orienté à $k\pi$ près]
Le plan est orienté. $A$ et $B$ sont deux points distincts avec $AB=8$ cm. Déterminer l'ensemble $(\mathcal{F})$ des points $M$ tels que
\[
(\overrightarrow{MA},\overrightarrow{MB})=\frac{\pi}{6}+k\pi,\quad k\in\mathbb{Z},
\]
et calculer le rayon du cercle support.
\tcblower
\textbf{Étape 1 : réduction modulo $\pi$.} $\dfrac{\pi}{6}\in\left]-\dfrac{\pi}{2}\,;\dfrac{\pi}{2}\right]$ : c'est la valeur réduite $x_0 = \dfrac{\pi}{6}$.

\textbf{Étape 2 : cercle support.} On construit $O$ sur la médiatrice de $[AB]$ tel que l'angle au centre vérifie $\operatorname{mes}(\widehat{AOB})=2\times\dfrac{\pi}{6}=\dfrac{\pi}{3}$. Le triangle $AOB$ est isocèle en $O$ avec un angle au sommet de $\dfrac{\pi}{3}$ : il est donc \textbf{équilatéral}, d'où $R=OA=AB=8$ cm.

Vérification par la formule : dans le triangle $AOB$ isocèle, $AB=2R\sin\dfrac{\widehat{AOB}}{2}$, donc
\[
R=\frac{AB}{2\sin\frac{\pi}{6}}=\frac{8}{2\times\frac12}=8\ \text{cm}.
\]

\textbf{Étape 3 : description de l'ensemble.} La condition étant donnée à $k\pi$ près, les deux arcs conviennent : si $M$ est sur un arc, $(\overrightarrow{MA},\overrightarrow{MB})=\dfrac{\pi}{6}+2k\pi$ ; sur l'autre, $(\overrightarrow{MA},\overrightarrow{MB})=\dfrac{\pi}{6}-\pi+2k\pi=-\dfrac{5\pi}{6}+2k\pi \equiv \dfrac{\pi}{6} \pmod{\pi}$. Dans les deux cas la condition est satisfaite.

\textbf{Conclusion :} $(\mathcal{F})$ est le cercle de centre $O$, de rayon $8$ cm, passant par $A$ et $B$, \textbf{privé des points $A$ et $B$}.
\end{exoresolu}

\courssub{Applications : lieux géométriques et cocyclicité}

\begin{methode}[Angle inscrit, angle au centre, tangente — Résumé des formules]
Dans un cercle $\mathcal{C}$ de centre $O$, $A,B,M \in \mathcal{C}$ distincts, $[At)$ tangente en $A$ du côté de $M$ :
\begin{enumerate}
\item \textbf{Angle inscrit / Angle au centre :} $\operatorname{mes}(\widehat{AOB})=2\,\operatorname{mes}(\widehat{AMB})$ (même arc intercepté). En orienté : $(\overrightarrow{OA},\overrightarrow{OB})=2(\overrightarrow{MA},\overrightarrow{MB})+2k\pi$.
\item \textbf{Corde et demi-tangente :} $\operatorname{mes}(\widehat{tAB})=\operatorname{mes}(\widehat{AMB})$. En orienté : $(\overrightarrow{At},\overrightarrow{AB})=(\overrightarrow{MA},\overrightarrow{MB})+2k\pi$.
\item \textbf{Deux angles inscrits} même arc $\Rightarrow$ égaux ; arcs opposés $\Rightarrow$ supplémentaires ($+\pi$).
\end{enumerate}
\textbf{Réflexe :} Repérer d'abord \emph{quel arc} est intercepté par l'angle considéré, puis choisir la bonne relation.
\end{methode}

\begin{exoresolu}[Calculs d'angles dans un cercle]
Soit $\mathcal{C}$ un cercle de centre $O$, $A$, $B$, $M$ trois points de $\mathcal{C}$ tels que $\operatorname{mes}(\widehat{AOB})=110°$ (angle au centre \textbf{rentrant}, interceptant l'arc mineur $\overset{\frown}{AB}$). La tangente en $A$ à $\mathcal{C}$ est notée $(At)$, avec $t$ du côté de $M$ (arc majeur).
\begin{enumerate}
\item Calculer $\operatorname{mes}(\widehat{AMB})$ lorsque $M$ est sur l'arc majeur.
\item Calculer $\operatorname{mes}(\widehat{tAB})$.
\item Que vaut $\operatorname{mes}(\widehat{ANB})$ si $N$ est sur l'arc mineur ?
\end{enumerate}
\tcblower
\textbf{1.} L'angle inscrit $\widehat{AMB}$ et l'angle au centre $\widehat{AOB}$ (rentrant, $110°$) interceptent le même arc (l'arc mineur $\overset{\frown}{AB}$ ne contenant pas $M$) :
\[
\operatorname{mes}(\widehat{AMB})=\frac12\operatorname{mes}(\widehat{AOB})=\frac{110°}{2}=55°.
\]

\textbf{2.} L'angle formé par la corde $[AB]$ et la demi-tangente $[At)$ est égal à l'angle inscrit interceptant le même arc (arc mineur) :
\[
\operatorname{mes}(\widehat{tAB})=\operatorname{mes}(\widehat{AMB})=55°.
\]

\textbf{3.} $M$ et $N$ sont sur des arcs opposés : les angles inscrits $\widehat{AMB}$ et $\widehat{ANB}$ interceptent des arcs complémentaires, ils sont donc supplémentaires :
\[
\operatorname{mes}(\widehat{ANB})=180°-55°=125°.
\]
(Vérification : l'angle au centre \emph{saillant} interceptant l'arc majeur vaut $360°-110°=250°$, et $\dfrac{250°}{2}=125°$.)

\textbf{Conclusion :} $\operatorname{mes}(\widehat{AMB})=\operatorname{mes}(\widehat{tAB})=55°$ et $\operatorname{mes}(\widehat{ANB})=125°$.
\end{exoresolu}

\begin{savaistu}[Le théodolite et l'arpentage au Cameroun]
Les géomètres-topographes de la MINDCAF (Ministère des Domaines, du Cadastre et des Affaires Foncières) utilisent le \cle{théodolite} pour mesurer des angles horizontaux et verticaux. Le principe de l'arc capable est à la base de la \emph{résolution de triangle} par intersection avant : pour localiser un point $M$ invisible, on mesure les angles $\widehat{AMB}$, $\widehat{BMC}$ depuis des points connus $A,B,C$. Le point $M$ est à l'intersection des arcs capables correspondants. C'est ainsi que sont nés les premiers plans cadastraux de Yaoundé et Douala !
\end{savaistu}

\begin{methode}[Prouver que quatre points sont cocycliques]
Pour montrer que $A$, $B$, $C$, $D$ sont sur un même cercle, on peut utiliser l'une des voies suivantes :
\begin{enumerate}
\item \textbf{Angles inscrits égaux (géométrique) :} Montrer que $\operatorname{mes}(\widehat{ACB})=\operatorname{mes}(\widehat{ADB})$ et que $C$, $D$ sont du \emph{même côté} de la droite $(AB)$. Alors $C$ et $D$ appartiennent au même arc capable de $[AB]$.
\item \textbf{Angles orientés :} Montrer que $(\overrightarrow{CA},\overrightarrow{CB})=(\overrightarrow{DA},\overrightarrow{DB})+k\pi$. Cela équivaut à la cocyclicité de $A,B,C,D$ (ou alignement, à exclure).
\item \textbf{Angles droits :} Si $\widehat{ACB}=\widehat{ADB}=\dfrac{\pi}{2}$, alors $C$ et $D$ sont sur le cercle de diamètre $[AB]$.
\item \textbf{Équidistance :} Exhiber un point $O$ tel que $OA=OB=OC=OD$.
\end{enumerate}
\textbf{Rédaction :} Toujours préciser que les points sont distincts et non alignés (sauf cas alignés) avant de conclure à la cocyclicité.
\end{methode}

\begin{exoresolu}[Cocyclicité au marché de Mokolo]
Un terrain a la forme d'un quadrilatère $ABCD$. On sait que $\operatorname{mes}(\widehat{ACB})=40°$ et $\operatorname{mes}(\widehat{ADB})=40°$, les points $C$ et $D$ étant du même côté de la droite $(AB)$. Montrer que $A$, $B$, $C$, $D$ sont cocycliques, puis calculer $\operatorname{mes}(\widehat{BAC})$ sachant que $\operatorname{mes}(\widehat{BDC})=30°$.
\tcblower
\textbf{Cocyclicité.} Les points $C$ et $D$ sont du même côté de $(AB)$ et
\[
\operatorname{mes}(\widehat{ACB})=\operatorname{mes}(\widehat{ADB})=40°.
\]
Donc $C$ et $D$ appartiennent au même arc capable du segment $[AB]$ (l'ensemble des points $M$ tels que $\operatorname{mes}(\widehat{AMB})=40°$ d'un côté fixé de $(AB)$ est un arc de cercle). Par suite, $A$, $B$, $C$, $D$ sont cocycliques.

Autre rédaction, en angles orientés : $(\overrightarrow{CA},\overrightarrow{CB})=(\overrightarrow{DA},\overrightarrow{DB})+k\pi$ car les deux angles géométriques sont égaux et $C$, $D$ du même côté de $(AB)$ ; donc les quatre points sont cocycliques.

\textbf{Calcul de $\widehat{BAC}$.} Dans le cercle passant par $A,B,C,D$, les angles inscrits $\widehat{BAC}$ et $\widehat{BDC}$ interceptent le même arc $\overset{\frown}{BC}$ (celui ne contenant ni $A$ ni $D$, les points étant dans l'ordre sur le cercle). Deux angles inscrits interceptant le même arc sont égaux :
\[
\operatorname{mes}(\widehat{BAC})=\operatorname{mes}(\widehat{BDC})=30°.
\]

\textbf{Conclusion :} Les quatre sommets du terrain sont cocycliques et $\operatorname{mes}(\widehat{BAC})=30°$.
\end{exoresolu}

\begin{methode}[Problème de lieu avec angle constant]
Pour trouver le lieu des points $M$ vérifiant une condition d'angle :
\begin{enumerate}
\item \textbf{Identifier} les points fixes (souvent $A$ et $B$) et l'angle constant $x$.
\item \textbf{Traduire} la condition : $\operatorname{mes}(\widehat{AMB})=x$ (géométrique), ou $(\overrightarrow{MA},\overrightarrow{MB})=x+2k\pi$, ou $+k\pi$ (orienté).
\item \textbf{Reconnaître} un arc capable : un arc (condition géométrique ou à $2k\pi$ près) ou le cercle entier (à $k\pi$ près), toujours \textbf{privé de $A$ et $B$}.
\item \textbf{Construire} le lieu : médiatrice de $[AB]$, demi-tangente en $A$ faisant l'angle $x$ (ou $x_0$ réduit), centre $O$, cercle.
\item \textbf{Conclure} en précisant les points exclus et, si le problème impose une région du plan (demi-plan, intérieur d'une figure), l'intersection du lieu avec cette région.
\end{enumerate}
\end{methode}

\begin{exoresolu}[Lieu d'un point d'observation]
Deux poteaux de la ligne électrique ENEO sont plantés en $A$ et $B$, distants de $10$ m. Un technicien placé en $M$ doit voir le segment $[AB]$ sous un angle droit : $\operatorname{mes}(\widehat{AMB})=\dfrac{\pi}{2}$. Déterminer et construire le lieu $(\mathcal{L})$ des positions possibles de $M$.
\tcblower
\textbf{Étape 1 : points fixes et angle.} $A$ et $B$ fixes, $AB=10$ m, $x=\dfrac{\pi}{2}$.

\textbf{Étape 2 : traduction.} $\operatorname{mes}(\widehat{AMB})=\dfrac{\pi}{2}$ signifie que le triangle $AMB$ est rectangle en $M$.

\textbf{Étape 3 : reconnaissance.} D'après le théorème de l'angle inscrit (cas $x=\dfrac{\pi}{2}$), $M$ appartient au cercle $\mathcal{C}$ de diamètre $[AB]$ : l'angle inscrit qui intercepte un diamètre est droit. Réciproquement, tout point $M$ de ce cercle, distinct de $A$ et $B$, vérifie $\operatorname{mes}(\widehat{AMB})=\dfrac{\pi}{2}$.

\textbf{Étape 4 : construction.} Le centre de $\mathcal{C}$ est le milieu $I$ de $[AB]$ et son rayon est
\[
R=\frac{AB}{2}=\frac{10}{2}=5\ \text{m}.
\]
On trace le cercle de centre $I$ et de rayon $5$ m.

\textbf{Étape 5 : conclusion.} $(\mathcal{L})$ est le cercle de diamètre $[AB]$, \textbf{privé des points $A$ et $B$} (en $A$ ou $B$, l'angle $\widehat{AMB}$ n'est pas défini). Le technicien peut se placer en tout point de ce cercle, à $5$ m du milieu de $[AB]$.
\end{exoresolu}

\begin{methode}[Réduction modulo $2\pi$ ou $\pi$ avant de conclure]
Face à une condition comme $(\overrightarrow{MA},\overrightarrow{MB})=\alpha+2k\pi$ avec $\alpha$ quelconque :
\begin{enumerate}
\item Écrire $\alpha$ sous la forme $\alpha_0+2p\pi$ avec $\alpha_0\in\left]-\pi\,;\pi\right]$ (mesure principale) : ajouter ou retrancher des multiples de $2\pi$.
\item Pour une condition à $k\pi$ près, réduire dans $\left]-\dfrac{\pi}{2}\,;\dfrac{\pi}{2}\right]$.
\item Appliquer alors les résultats sur les arcs capables avec la valeur réduite $\alpha_0$.
\end{enumerate}
\textbf{Piège classique :} $\dfrac{7\pi}{3}$ et $\dfrac{\pi}{3}$ définissent le \emph{même} lieu à $2k\pi$ près ; ne pas construire un angle de $\dfrac{7\pi}{3}$ ! De même, ne jamais mélanger degrés et radians dans un même calcul ($\sin 45 \neq \frac{\sqrt2}{2}$, c'est $\sin 45°$ ou $\sin\frac{\pi}{4}$).
\end{methode}

\begin{exoresolu}[Réduction puis lieu]
Le plan est orienté. $A$ et $B$ sont deux points distincts. Déterminer l'ensemble $(\mathcal{G})$ des points $M$ tels que
\[
(\overrightarrow{MA},\overrightarrow{MB})=\frac{9\pi}{4}+2k\pi,\quad k\in\mathbb{Z}.
\]
\tcblower
\textbf{Étape 1 : réduction modulo $2\pi$.} On cherche $p\in\mathbb{Z}$ tel que $\dfrac{9\pi}{4}-2p\pi\in\left]-\pi\,;\pi\right]$ :
\[
\frac{9\pi}{4}-2\pi=\frac{9\pi}{4}-\frac{8\pi}{4}=\frac{\pi}{4}\in\left]-\pi\,;\pi\right].
\]
La condition équivaut donc à $(\overrightarrow{MA},\overrightarrow{MB})=\dfrac{\pi}{4}+2k\pi$.

\textbf{Étape 2 : construction.} On construit la demi-droite $[At)$ telle que $(\overrightarrow{At},\overrightarrow{AB})=\dfrac{\pi}{4}+2k\pi$ (angle direct de $45°$), puis la perpendiculaire en $A$ à $[At)$ qui coupe la médiatrice de $[AB]$ en $O$. On trace le cercle $\mathcal{C}(O,OA)$.

\textbf{Étape 3 : choix de l'arc.} La condition étant à $2k\pi$ près, un seul arc convient : l'arc d'extrémités $A$ et $B$ du côté de $[At)$, pour lequel l'angle orienté $(\overrightarrow{MA},\overrightarrow{MB})$ est direct et vaut $\dfrac{\pi}{4}$.

\textbf{Conclusion :} $(\mathcal{G})$ est l'arc capable orienté associé à l'angle $\dfrac{\pi}{4}$, situé du côté de $[At)$ par rapport à $(AB)$, \textbf{privé de $A$ et $B$}. C'est le même lieu que pour la condition $(\overrightarrow{MA},\overrightarrow{MB})=\dfrac{\pi}{4}+2k\pi$.
\end{exoresolu}

\begin{attention}[Les 5 erreurs qui coûtent des points au Probatoire]
\begin{enumerate}
\item \textbf{Oublier d'exclure $A$ et $B$.} En $M=A$ ou $M=B$, les vecteurs $\overrightarrow{MA}$ ou $\overrightarrow{MB}$ sont nuls et l'angle n'est pas défini. Toute conclusion de lieu doit se terminer par : « ... privé des points $A$ et $B$ ». C'est la première chose que le correcteur vérifie.
\item \textbf{Confondre les conditions à $2k\pi$ près et à $k\pi$ près.} À $2k\pi$ près : \emph{un seul} arc du cercle. À $k\pi$ près : le \emph{cercle entier} (les deux arcs). Écrire « le cercle » pour une condition à $2k\pi$ près est une erreur de cours éliminatoire.
\item \textbf{Se tromper de côté pour la demi-tangente.} L'arc capable est du \emph{même côté} que la demi-tangente $[At)$ par rapport à $(AB)$. Vérifier systématiquement en testant un point de l'arc construit.
\item \textbf{Oublier le facteur 2 entre angle inscrit et angle au centre.} L'angle au centre vaut le \emph{double} de l'angle inscrit : $\operatorname{mes}(\widehat{AOB})=2\operatorname{mes}(\widehat{AMB})$, et non l'inverse. Dans les constructions, c'est $2x$ qu'il faut placer en $O$ (ou $x$ à la tangente).
\item \textbf{Ne pas réduire l'angle avant de construire.} Pour $(\overrightarrow{MA},\overrightarrow{MB})=\dfrac{9\pi}{4}+2k\pi$, on réduit d'abord en $\dfrac{\pi}{4}$. Construire un angle non réduit, ou mélanger degrés et radians dans un même calcul, fausse tout l'exercice.
\end{enumerate}
\end{attention}

\begin{exercice}[Entraînement 1 : Construction et calcul]
Soit $[AB]$ un segment de longueur $10$ cm. On note $(\Gamma)$ l'ensemble des points $M$ tels que $\operatorname{mes}(\widehat{AMB})=60°$.
\begin{enumerate}
\item Construire $(\Gamma)$ à la règle et au compas (laisser les traces).
\item Calculer le rayon du cercle support de $(\Gamma)$.
\item $M$ est un point de $(\Gamma)$ tel que $MA=6$ cm. Calculer $MB$.
\end{enumerate}
\end{exercice}

\begin{corrige}[Exercice 1]
\begin{enumerate}
\item Construction standard (Méthode 1) : médiatrice, tangente $60°$ en $A$, perpendiculaire, centre $O$, cercle, arc du côté de la tangente privé de $A,B$.
\item Angle au centre $= 120°$. Triangle $AOB$ isocèle en $O$, angle sommet $120°$. Hauteur depuis $O$ : $R\cos 60° = \frac{AB}{2} = 5 \Rightarrow R = 10$ cm. (Ou $AB = 2R\sin 60° \Rightarrow 10 = 2R\frac{\sqrt{3}}{2} \Rightarrow R = \frac{10}{\sqrt{3}}$). \textbf{Attention :} $\widehat{AOB}=120°$, la moitié est $60°$. $AB = 2R\sin 60° \Rightarrow R = \frac{10}{\sqrt{3}} \approx 5,77$ cm.
\item Dans le triangle $AMB$, on connaît $AB=10$, $MA=6$, $\widehat{AMB}=60°$. Loi des cosinus : $AB^2 = MA^2 + MB^2 - 2\cdot MA\cdot MB\cos 60°$. $100 = 36 + MB^2 - 6\cdot MB$. $MB^2 - 6MB - 64 = 0$. $\Delta = 36 + 256 = 292$. $MB = \frac{6+\sqrt{292}}{2} = 3+\sqrt{73}$ cm (longueur positive).
\end{enumerate}
\end{corrige}

\begin{exercice}[Entraînement 2 : Lieu orienté]
Le plan est orienté directement. $A(0;0)$, $B(4;0)$. Déterminer l'ensemble des points $M(x;y)$ tels que :
\[
(\overrightarrow{MA},\overrightarrow{MB}) = -\frac{\pi}{3} + 2k\pi.
\]
Donner une équation cartésienne du cercle support et préciser l'arc (inéquation).
\end{exercice}

\begin{corrige}[Exercice 2]
Mesure principale $x_0 = -\frac{\pi}{3}$ (négatif $\Rightarrow$ angle indirect).
Construction : Tangente en $A$ telle que $(\overrightarrow{At},\overrightarrow{AB}) = -\frac{\pi}{3}$. Perpendiculaire en $A$ coupe médiatrice $x=2$ en $O$.
Coordonnées : $A(0,0)$, $B(4,0)$. Tangente en $A$ : pente $\tan(-\pi/3)=-\sqrt{3}$. Équation $y=-\sqrt{3}x$. Perpendiculaire en $A$ : pente $\frac{\sqrt{3}}{3}$, équation $y=\frac{\sqrt{3}}{3}x$. Intersection avec $x=2$ : $O(2; \frac{2\sqrt{3}}{3})$.
Rayon $R = OA = \sqrt{4 + \frac{4}{3}} = \sqrt{\frac{16}{3}} = \frac{4}{\sqrt{3}}$.
Cercle : $(x-2)^2 + (y-\frac{2\sqrt{3}}{3})^2 = \frac{16}{3}$.
L'arc est celui situé \emph{en dessous} de l'axe $(AB)$ (côté de la tangente $y=-\sqrt{3}x$ pour $x>0$), c'est-à-dire $y < 0$ (sauf $A,B$).
Lieu : arc du cercle $(x-2)^2 + (y-\frac{2\sqrt{3}}{3})^2 = \frac{16}{3}$ avec $y < 0$, privé de $A(0,0)$ et $B(4,0)$.
\end{corrige}

\newpage

\coursec{Exercices}

\courssub{Je vérifie mes connaissances}

\begin{exercice}[QCM : mesures principales]
Pour chaque question, une seule réponse est exacte. Laquelle ?
\begin{enumerate}
\item La mesure principale de l'angle $\dfrac{17\pi}{3}$ est :
\begin{itemize}
\item a) $\dfrac{5\pi}{3}$ \qquad b) $-\dfrac{\pi}{3}$ \qquad c) $\dfrac{\pi}{3}$ \qquad d) $\dfrac{2\pi}{3}$
\end{itemize}
\item La mesure principale de l'angle $-\dfrac{11\pi}{4}$ est :
\begin{itemize}
\item a) $\dfrac{5\pi}{4}$ \qquad b) $-\dfrac{3\pi}{4}$ \qquad c) $\dfrac{3\pi}{4}$ \qquad d) $-\dfrac{\pi}{4}$
\end{itemize}
\item La mesure principale de $\dfrac{25\pi}{6}$ est :
\begin{itemize}
\item a) $\dfrac{\pi}{6}$ \qquad b) $-\dfrac{\pi}{6}$ \qquad c) $\dfrac{7\pi}{6}$ \qquad d) $\dfrac{11\pi}{6}$
\end{itemize}
\item Placer ensuite sur un cercle trigonométrique les points images des trois mesures principales obtenues.
\end{enumerate}
\end{exercice}

\begin{exercice}[Vrai ou faux, en justifiant]
$A$, $B$, $M$, $N$ sont quatre points d'un cercle $(\mathcal{C})$ de centre $O$, tous distincts de $A$ et $B$. Dire si chacune des affirmations suivantes est vraie ou fausse, en justifiant.
\begin{enumerate}
\item Si $\mathrm{mes}\,\widehat{AMB} = 40^{\circ}$, alors $\mathrm{mes}\,\widehat{AOB} = 80^{\circ}$ (angle au centre interceptant le même arc).
\item Si $M$ et $N$ sont sur le même arc d'extrémités $A$ et $B$, alors $\mathrm{mes}\,\widehat{AMB} = \mathrm{mes}\,\widehat{ANB}$.
\item Si $M$ et $N$ sont sur des arcs différents d'extrémités $A$ et $B$, alors $\mathrm{mes}\,\widehat{AMB} + \mathrm{mes}\,\widehat{ANB} = 180^{\circ}$.
\item L'ensemble des points $M$ tels que $\mathrm{mes}\,\widehat{AMB} = 90^{\circ}$ est le cercle de diamètre $[AB]$ privé des points $A$ et $B$.
\end{enumerate}
\end{exercice}

\begin{exercice}[Angle inscrit et demi-tangente]
Soit $(\mathcal{C})$ un cercle de centre $O$, $[AB]$ une corde et $(AT)$ la demi-tangente en $A$ telle que $\mathrm{mes}\,\widehat{TAB} = 55^{\circ}$.
\begin{enumerate}
\item Calculer la mesure de l'angle au centre $\widehat{AOB}$ interceptant l'arc $\wideparen{AB}$ ne contenant pas le secteur adjacent à $(AT)$.
\item En déduire la mesure de tout angle inscrit $\widehat{AMB}$ où $M$ appartient à l'arc capable associé à $(AT)$.
\item Que vaut $\mathrm{mes}\,\widehat{ANB}$ si $N$ appartient à l'autre arc d'extrémités $A$ et $B$ ?
\end{enumerate}
\end{exercice}

\begin{exercice}[Lecture d'angles orientés]
Dans le plan orienté, $A$ et $B$ sont deux points distincts. Pour chaque condition, dire (sans construction) si l'ensemble cherché est un arc de cercle, un cercle privé de deux points, une réunion d'arcs, ou une droite privée de points :
\begin{enumerate}
\item $\mathrm{mes}(\overrightarrow{MA}, \overrightarrow{MB}) = 0 \; [\pi]$ ;
\item $\mathrm{mes}(\overrightarrow{MA}, \overrightarrow{MB}) = \dfrac{\pi}{2} \; [\pi]$ ;
\item $\mathrm{mes}(\overrightarrow{MA}, \overrightarrow{MB}) = \dfrac{\pi}{3} \; [2\pi]$ ;
\item $\mathrm{mes}(\overrightarrow{MA}, \overrightarrow{MB}) = \dfrac{\pi}{3} \; [\pi]$.
\end{enumerate}
\end{exercice}

\courssub{Je m'entraîne}

\begin{exercice}[Construction d'un arc capable]
On donne un segment $[AB]$ tel que $AB = \SI{8}{cm}$.
\begin{enumerate}
\item Calculer le rayon $R$ du cercle dont un arc est l'ensemble des points $M$ tels que $\mathrm{mes}\,\widehat{AMB} = 75^{\circ}$. On donnera la valeur exacte en fonction de $\sin 75^{\circ}$, puis une valeur approchée à $\SI{0,1}{cm}$ près.
\item Construire à la règle et au compas cet arc capable : on tracera la demi-droite $[At)$ telle que $\mathrm{mes}\,\widehat{tAB} = 75^{\circ}$, puis la perpendiculaire en $A$ à $[At)$ et la médiatrice de $[AB]$.
\item Vérifier sur la figure qu'un point $M$ de l'arc construit satisfait bien $\mathrm{mes}\,\widehat{AMB} = 75^{\circ}$.
\end{enumerate}
\end{exercice}

\begin{exercice}[Arcs capables orientés]
Dans le plan orienté, on donne deux points $A$ et $B$ distincts.
\begin{enumerate}
\item Déterminer et représenter en rouge l'ensemble $(\mathcal{E}_1)$ des points $M$ tels que $\mathrm{mes}(\overrightarrow{MA}, \overrightarrow{MB}) = \dfrac{\pi}{4} \; [2\pi]$.
\item Déterminer et représenter en bleu l'ensemble $(\mathcal{E}_2)$ des points $M$ tels que $\mathrm{mes}(\overrightarrow{MA}, \overrightarrow{MB}) = \dfrac{\pi}{4} \; [\pi]$.
\item Comparer $(\mathcal{E}_1)$ et $(\mathcal{E}_2)$ : lequel est inclus dans l'autre ? Décrire $(\mathcal{E}_2) \setminus (\mathcal{E}_1)$ par une condition d'angle orienté.
\end{enumerate}
\end{exercice}

\begin{exercice}[Angle droit et rapport de distances]
Dans le plan orienté, $A$ et $B$ sont deux points tels que $AB = \SI{6}{cm}$.
\begin{enumerate}
\item Déterminer l'ensemble $(\mathcal{C})$ des points $M$ tels que $\mathrm{mes}(\overrightarrow{MA}, \overrightarrow{MB}) = \dfrac{\pi}{2} \; [\pi]$. Le construire.
\item On cherche de plus les points $M$ vérifiant $MA = 2\,MB$. Montrer que ces points appartiennent à un cercle $(\Gamma)$ dont on précisera le diamètre (on pourra utiliser les points $I$ et $J$ de la droite $(AB)$ définis par $\overrightarrow{IA} = -2\overrightarrow{IB}$ et $\overrightarrow{JA} = 2\overrightarrow{JB}$).
\item Construire $(\Gamma)$, puis déterminer graphiquement les points $M$ solutions des deux conditions simultanées. Combien y en a-t-il ?
\end{enumerate}
\end{exercice}

\begin{exercice}[Cocyclicité au marché]
Quatre villages $A$, $B$, $C$, $D$ de la région de l'Ouest sont modélisés par quatre points du plan. Un topographe mesure : $\mathrm{mes}\,\widehat{ACB} = 42^{\circ}$ et $\mathrm{mes}\,\widehat{ADB} = 42^{\circ}$, les points $C$ et $D$ étant situés du même côté de la route $(AB)$.
\begin{enumerate}
\item Justifier que les quatre villages sont situés sur un même cercle.
\item Le chef de village $E$ affirme que depuis chez lui on voit le segment $[AB]$ sous un angle de $138^{\circ}$, en étant de l'autre côté de la route $(AB)$. Le point $E$ appartient-il au même cercle ? Justifier.
\item Construire le cercle passant par $A$, $B$ et $C$, puis placer un point $D'$ possible.
\end{enumerate}
\end{exercice}

\courssub{Je me prépare au Probatoire}

\begin{exercice}[Orthocentre et cercle circonscrit]
Soit $ABC$ un triangle acutangle, $H$ son orthocentre et $(\mathcal{C})$ son cercle circonscrit. On note $H_1$ le symétrique de $H$ par rapport à la droite $(BC)$, et $A'$ le pied de la hauteur issue de $A$ (donc $A' \in (BC)$ et $(AA') \perp (BC)$).

\textbf{Partie A : questions préliminaires}
\begin{enumerate}
\item Faire une figure soignée. Montrer que $\mathrm{mes}(\overrightarrow{HB}, \overrightarrow{HC}) = \mathrm{mes}(\overrightarrow{AB}, \overrightarrow{AC}) + \pi \; [2\pi]$. \emph{Indication : utiliser $(HB) \perp (AC)$ et $(HC) \perp (AB)$.}
\item En déduire que $\mathrm{mes}(\overrightarrow{HB}, \overrightarrow{HC}) = \pi - \mathrm{mes}\,\widehat{BAC}$ en mesures géométriques.
\end{enumerate}

\textbf{Partie B : symétrique de l'orthocentre}
\begin{enumerate}
\item Justifier que $\mathrm{mes}(\overrightarrow{H_1B}, \overrightarrow{H_1C}) = -\mathrm{mes}(\overrightarrow{HB}, \overrightarrow{HC}) \; [2\pi]$.
\item En déduire que $\mathrm{mes}(\overrightarrow{H_1B}, \overrightarrow{H_1C}) = \mathrm{mes}(\overrightarrow{AB}, \overrightarrow{AC}) \; [\pi]$.
\item Conclure que $H_1$ appartient au cercle $(\mathcal{C})$.
\item On note $H_2$ et $H_3$ les symétriques de $H$ par rapport à $(CA)$ et $(AB)$. Que peut-on dire de $H_2$ et $H_3$ ? Justifier brièvement.
\end{enumerate}

\textbf{Partie C : application numérique}

On suppose $\mathrm{mes}\,\widehat{BAC} = 60^{\circ}$ et $BC = \SI{9}{cm}$. Calculer le rayon du cercle $(\mathcal{C})$.
\end{exercice}

\begin{exercice}[Deux cercles sécants et angle constant]
Deux cercles $(\mathcal{C}_1)$ et $(\mathcal{C}_2)$ se coupent en deux points $A$ et $B$. Une droite variable $(\Delta)$ passant par $A$ recoupe $(\mathcal{C}_1)$ en $M$ et $(\mathcal{C}_2)$ en $N$ (avec $M \neq A$ et $N \neq A$). Le plan est orienté.

\textbf{Partie A : conjecture}
\begin{enumerate}
\item Faire deux figures correspondant à deux positions différentes de $(\Delta)$. Mesurer à chaque fois l'angle $\widehat{MBN}$. Quelle conjecture peut-on formuler ?
\end{enumerate}

\textbf{Partie B : démonstration}
\begin{enumerate}
\item Justifier que $\mathrm{mes}(\overrightarrow{BM}, \overrightarrow{BA}) = \mathrm{mes}(\overrightarrow{MA}, \overrightarrow{MB})$... plus précisément, montrer à l'aide des angles inscrits dans $(\mathcal{C}_1)$ que $\mathrm{mes}(\overrightarrow{BM}, \overrightarrow{BA})$ est constant modulo $\pi$ lorsque $(\Delta)$ varie.
\item Montrer de même que $\mathrm{mes}(\overrightarrow{BA}, \overrightarrow{BN})$ est constant modulo $\pi$.
\item En déduire, par la relation de Chasles, que $\mathrm{mes}(\overrightarrow{BM}, \overrightarrow{BN})$ est constant modulo $\pi$. Exprimer cette constante en fonction des angles que forment les tangentes en $A$ aux deux cercles.
\end{enumerate}

\textbf{Partie C : lieu géométrique}

On suppose désormais que $N$ décrit $(\mathcal{C}_2)$ privé de $A$ et $B$, et que $M$ est défini comme ci-dessus.
\begin{enumerate}
\item Montrer que $M$ décrit $(\mathcal{C}_1)$ privé de $A$ et $B$.
\item On fixe $B$ et l'on fait varier $(\Delta)$ : quel est l'ensemble des milieux $I$ de $[MN]$ lorsque les deux cercles sont concentriques de centre $O$ ? (On admettra que $I$ décrit un cercle ; préciser lequel.)
\end{enumerate}
\end{exercice}

\begin{exercice}[Navigation entre deux phares]
Deux phares $F_1$ et $F_2$ sont situés sur la côte atlantique, à proximité de Kribi. Sur une carte à l'échelle $1/100\,000$, la distance $F_1F_2$ est représentée par $\SI{12}{cm}$. Un banc de rochers dangereux se trouve au voisinage du segment $[F_1F_2]$ : pour naviguer en sécurité, le capitaine doit maintenir son bateau $B$ de sorte que les deux phares soient vus sous un angle $\widehat{F_1BF_2}$ \emph{supérieur ou égal à} $30^{\circ}$, du côté mer de la côte.

\textbf{Partie A : étude théorique}
\begin{enumerate}
\item Quelle est la distance réelle entre les deux phares, en kilomètres ?
\item Calculer le rayon $R$ de l'arc capable correspondant à l'angle $30^{\circ}$ construit sur $[F_1F_2]$, d'abord sur la carte (en cm), puis en réalité (en km). On rappelle que $\sin 30^{\circ} = \dfrac{1}{2}$.
\item Décrire précisément la zone de navigation sûre : s'agit-il de l'intérieur ou de l'extérieur du segment de cercle ? Justifier en comparant les angles inscrits pour un point intérieur, sur, ou extérieur à l'arc capable.
\end{enumerate}

\textbf{Partie B : construction sur la carte}
\begin{enumerate}
\item Construire à la règle et au compas, sur une figure où $F_1F_2 = \SI{12}{cm}$, l'arc capable d'angle $30^{\circ}$ du côté mer.
\item Hachurer la zone de navigation sûre.
\item Le bateau se trouve en un point $B$ tel que $BF_1 = \SI{10}{cm}$ (sur la carte) et $\mathrm{mes}\,\widehat{F_1BF_2} = 30^{\circ}$. Calculer la distance $BF_2$ sur la carte, puis la distance réelle du bateau au phare $F_2$. \emph{Indication : utiliser la loi des sinus dans le triangle $F_1BF_2$ après avoir déterminé la nature ou les angles du triangle.}
\end{enumerate}

\textbf{Partie C : prolongement}

Le capitaine reçoit l'ordre de se rapprocher de la côte tout en gardant un angle de visée \emph{exactement égal à} $45^{\circ}$. Décrire et construire la trajectoire que doit suivre le bateau.
\end{exercice}

\newpage

\coursec{Corrigés des exercices}

\coursec{Angles orientés et arcs capables}

\courssub{Rappels sur les angles orientés de vecteurs}

\begin{definition}[Angle orienté de deux vecteurs]
Soient $\vec{u}$ et $\vec{v}$ deux vecteurs non nuls du plan. L'\cle{angle orienté} $(\vec{u}, \vec{v})$ est la mesure de la rotation qui transforme $\vec{u}$ en $\vec{v}$ dans le sens direct (trigonométrique). On note $\mathrm{mes}(\vec{u}, \vec{v})$ une de ses mesures. Toute mesure est définie à $2\pi$ près : si $\alpha$ est une mesure, alors $\alpha + 2k\pi\ (k \in \mathbb{Z})$ en sont d'autres.
\end{definition}

\begin{definition}[Mesure principale]
La \cle{mesure principale} d'un angle orienté est l'unique mesure appartenant à l'intervalle $]-\pi\,;\,\pi]$. Elle se note $\mathrm{mes}_p(\vec{u}, \vec{v})$.
\end{definition}

\begin{propriete}[Propriétés fondamentales]
Pour tous vecteurs non nuls $\vec{u}, \vec{v}, \vec{w}$ :
\begin{enumerate}
\item \textbf{Relation de Chasles :} $\mathrm{mes}(\vec{u}, \vec{w}) = \mathrm{mes}(\vec{u}, \vec{v}) + \mathrm{mes}(\vec{v}, \vec{w})\ [2\pi]$.
\item \textbf{Antisymétrie :} $\mathrm{mes}(\vec{v}, \vec{u}) = -\mathrm{mes}(\vec{u}, \vec{v})\ [2\pi]$.
\item \textbf{Colinéarité :} $\mathrm{mes}(\vec{u}, \vec{v}) = 0\ [\pi] \iff \vec{u} \parallel \vec{v}$.
\item \textbf{Orthogonalité :} $\mathrm{mes}(\vec{u}, \vec{v}) = \dfrac{\pi}{2}\ [\pi] \iff \vec{u} \perp \vec{v}$.
\end{enumerate}
\end{propriete}

\begin{attention}[Pièges classiques]
\begin{itemize}
\item Ne jamais écrire $\mathrm{mes}(\vec{u}, \vec{v}) = \alpha$ sans préciser le modulo ($[2\pi]$ ou $[\pi]$). Une égalité sans modulo est fausse en général.
\item La mesure principale est dans $]-\pi\,;\,\pi]$ : $\pi$ est admis, $-\pi$ ne l'est pas (on écrit $\pi$).
\item Modulo $\pi$, on perd l'orientation : $\alpha\ [\pi]$ regroupe $\alpha\ [2\pi]$ et $-\alpha\ [2\pi]$.
\end{itemize}
\end{attention}

\begin{methode}[Ramener un angle à sa mesure principale]
Pour un angle $\alpha$ donné (en radians) :
\begin{enumerate}
\item Diviser $\alpha$ par $2\pi$ : $\alpha = 2k\pi + r$ avec $r \in [0, 2\pi[$ (division euclidienne sur les réels).
\item Si $r \le \pi$, la mesure principale est $r$.
\item Si $r > \pi$, la mesure principale est $r - 2\pi$ (qui appartient à $]-\pi, 0[$).
\end{enumerate}
\end{methode}

\begin{exemplebox}[Calculs de mesures principales]
\begin{align*}
\frac{17\pi}{3} &= 6\pi - \frac{\pi}{3} = 3\times 2\pi - \frac{\pi}{3} &\Rightarrow\ \mathrm{mes}_p &= -\frac{\pi}{3} \\
-\frac{11\pi}{4} &= -3\pi + \frac{\pi}{4} = -2\pi - \pi + \frac{\pi}{4} = -2\pi - \frac{3\pi}{4} &\Rightarrow\ \mathrm{mes}_p &= -\frac{3\pi}{4} \\
\frac{25\pi}{6} &= 4\pi + \frac{\pi}{6} = 2\times 2\pi + \frac{\pi}{6} &\Rightarrow\ \mathrm{mes}_p &= \frac{\pi}{6}
\end{align*}
\end{exemplebox}

\courssub{Angle inscrit, angle au centre et demi-tangente}

\begin{definition}[Angles au centre et inscrit]
Soit un cercle de centre $O$. Pour trois points $A, B, M$ distincts sur ce cercle :
\begin{itemize}
\item L'\cle{angle au centre} $\widehat{AOB}$ intercepte l'arc $\wideparen{AB}$ ne contenant pas $O$ (ou l'arc opposé à $M$).
\item L'\cle{angle inscrit} $\widehat{AMB}$ intercepte l'arc $\wideparen{AB}$ ne contenant pas $M$.
\end{itemize}
\end{definition}

\begin{propriete}[Théorème de l'angle au centre]
Dans un même cercle, la mesure d'un angle au centre est le double de la mesure de tout angle inscrit interceptant le même arc :
\[ \mathrm{mes}\,\widehat{AOB} = 2 \times \mathrm{mes}\,\widehat{AMB}. \]
\end{propriete}

\begin{propriete}[Conséquences immédiates]
\begin{enumerate}
\item Deux angles inscrits interceptant le même arc ont même mesure.
\item Deux angles inscrits interceptant des arcs supplémentaires (extrémités communes, arcs opposés) sont supplémentaires (somme $= \pi$ ou $180^\circ$).
\item \textbf{Angle inscrit dans un demi-cercle :} Si $[AB]$ est un diamètre, alors pour tout $M \neq A,B$ sur le cercle, $\mathrm{mes}\,\widehat{AMB} = \dfrac{\pi}{2}\ (90^\circ)$. Réciproquement, si $\mathrm{mes}\,\widehat{AMB} = \dfrac{\pi}{2}$, alors $M$ appartient au cercle de diamètre $[AB]$.
\end{enumerate}
\end{propriete}

\begin{definition}[Angle formé par une corde et une demi-tangente]
Soit un cercle de centre $O$, un point $A$ sur ce cercle et la tangente $(AT)$ en $A$. Pour un point $B \neq A$ sur le cercle, l'angle $\widehat{TAB}$ (ou $\widehat{BAT}$) est l'angle formé par la demi-tangente $(AT)$ et la corde $[AB]$.
\end{definition}

\begin{propriete}[Théorème de l'angle entre tangente et corde]
La mesure de l'angle formé par une corde et la tangente en l'une de ses extrémités est égale à la mesure de l'angle inscrit interceptant l'arc opposé à cet angle (arc situé de l'autre côté de la corde par rapport à la tangente).
\[ \mathrm{mes}\,\widehat{TAB} = \mathrm{mes}\,\widehat{AMB} \quad \text{où } M \text{ est sur l'arc opposé à } (AT). \]
\end{propriete}

\begin{experience}[Découverte : angle tangente/corde]
\begin{enumerate}
\item Sur un cercle de centre $O$, tracer un diamètre $[AB]$. La tangente en $A$ est perpendiculaire à $(AB)$.
\item Choisir $M$ sur le cercle (arc supérieur). Mesurer $\widehat{AMB}$ : c'est un angle droit ($90^\circ$).
\item Mesurer $\widehat{TAB}$ entre la tangente $(AT)$ (perpendiculaire à $AB$) et la corde $[AB]$ : c'est aussi $90^\circ$.
\item Déplacer $M$ sur l'arc supérieur : $\widehat{AMB}$ reste $90^\circ$. L'égalité $\widehat{TAB} = \widehat{AMB}$ se vérifie pour tout $M$ sur l'arc opposé à la tangente.
\end{enumerate}
\end{experience}

\courssub{Arc capable : ensemble des points M tels que mes(AMB) = x}

\begin{definition}[Arc capable (angle géométrique)]
Soient $A, B$ deux points distincts et $x \in [0, \pi]$. L'\cle{arc capable} d'angle $x$ et d'extrémités $A, B$ est l'ensemble des points $M$ du plan (différents de $A$ et $B$) tels que l'angle géométrique $\widehat{AMB}$ a pour mesure $x$.
\begin{itemize}
\item Si $x = 0$ : c'est la droite $(AB)$ privée de $A$ et $B$ (demi-droites $[AB)$ et $(BA)$).
\item Si $x = \pi$ : c'est le segment $[AB]$ privé de $A$ et $B$.
\item Si $x \in ]0, \pi[$ : ce sont \textbf{deux arcs de cercle} symétriques par rapport à la droite $(AB)$, d'extrémités $A$ et $B$, privés de $A$ et $B$.
\end{itemize}
\end{definition}

\begin{propriete}[{Construction de l'arc capable (angle $x \in ]0, \pi[$)}]
Pour construire l'arc capable d'extrémités $A, B$ et d'angle $x$ du côté choisi de $(AB)$ :
\begin{enumerate}
\item Tracer le segment $[AB]$.
\item En $A$, construire la demi-droite $[At)$ telle que $\mathrm{mes}\,\widehat{tAB} = x$ (du côté choisi).
\item Tracer la perpendiculaire à $[At)$ en $A$ (c'est le rayon $OA$).
\item Tracer la médiatrice de $[AB]$.
\item L'intersection $O$ de ces deux droites est le centre du cercle. Le rayon est $R = OA = OB$.
\item L'arc capable est l'arc de cercle d'extrémités $A, B$ situé du même côté que $[At)$ par rapport à $(AB)$, privé de $A$ et $B$.
\end{enumerate}
\end{propriete}

\begin{aretenir}[Rayon de l'arc capable]
Dans le triangle isocèle $OAB$ ($OA=OB=R$), l'angle au centre $\widehat{AOB}$ intercepte l'arc $AB$ opposé à $M$. Comme $\mathrm{mes}\,\widehat{AMB} = x$, on a $\mathrm{mes}\,\widehat{AOB} = 2x$ (si $x \le \pi/2$) ou $2(\pi-x)$ (si $x > \pi/2$). En notant $I$ le milieu de $[AB]$, le triangle $OIA$ est rectangle en $I$ et :
\[ \sin x = \frac{AB/2}{R} \quad \Longrightarrow \quad R = \frac{AB}{2\sin x}. \]
(Valable pour tout $x \in ]0, \pi[$ car $\sin x = \sin(\pi-x)$).
\end{aretenir}

\begin{exemplebox}[Construction : $AB = 8$ cm, $x = 75^\circ$]
$R = \dfrac{8}{2\sin 75^\circ} = \dfrac{4}{\sin 75^\circ}$. $\sin 75^\circ = \dfrac{\sqrt{6}+\sqrt{2}}{4} \approx 0,966 \Rightarrow R \approx 4,14$ cm.
\begin{popfigure}
\begin{tikzpicture}[scale=0.8, >=Stealth]
\coordinate (A) at (-4,0);
\coordinate (B) at (4,0);
\coordinate (O) at (0, {4/tan(75)}); % R*cos(75) = 4/tan(75) approx 1.07
\coordinate (I) at (0,0);
\draw[popblue, thick] (A) -- (B) node[midway, below=2pt] {$AB=8$};
\draw[poporange, dashed] (A) -- ++(75:5) node[pos=0.7, above right] {$[At)$};
\draw[popgreen, thick] (A) -- (O) node[midway, left] {$R$};
\draw[popgreen, thick] (O) -- (B);
\draw[popmuted, thin] (O) circle ({4/sin(75)});
\draw[popblue, very thick] (A) arc (180+75: -75: {4/sin(75)});
\draw[poporange, thick] (A) -- ++(75:1.5);
\draw[popgreen, thick] (A) -- ++(-15:1.5); % perpendicular to tangent
\fill (A) circle (2pt) node[below left] {$A$};
\fill (B) circle (2pt) node[below right] {$B$};
\fill (O) circle (2pt) node[above] {$O$};
\fill (I) circle (2pt) node[below] {$I$};
\node at (0, -1.5) {Arc capable $75^\circ$};
\end{tikzpicture}
\legende{Construction de l'arc capable d'angle $75^\circ$ sur $[AB]=8$ cm.}
\end{popfigure}
\end{exemplebox}

\courssub{Arcs capables orientés : mes(MA, MB) = x + 2k\pi et x + k\pi}

\begin{definition}[Angle orienté de deux droites / vecteurs position]
Pour trois points $A, B, M$ distincts, on note $\mathrm{mes}(\overrightarrow{MA}, \overrightarrow{MB})$ la mesure de l'angle orienté des vecteurs $\overrightarrow{MA}$ et $\overrightarrow{MB}$. C'est l'angle orienté des droites $(MA)$ et $(MB)$.
\end{definition}

\begin{propriete}[{Lieu : $\mathrm{mes}(\overrightarrow{MA}, \overrightarrow{MB}) = \alpha\ [2\pi]$}]
Soient $A, B$ distincts et $\alpha \in ]-\pi, \pi]$, $\alpha \neq 0$.
L'ensemble $\mathcal{E}_{2\pi} = \{ M \in \mathcal{P} \setminus \{A,B\} \mid \mathrm{mes}(\overrightarrow{MA}, \overrightarrow{MB}) = \alpha\ [2\pi] \}$ est \textbf{un unique arc de cercle} d'extrémités $A, B$ (privé de $A, B$), situé du côté de $(AB)$ où la rotation de $\overrightarrow{MA}$ vers $\overrightarrow{MB}$ s'effectue dans le sens direct (trigonométrique).
\begin{itemize}
\item Si $\alpha > 0$ : l'arc est du côté où l'on tourne dans le sens direct.
\item Si $\alpha < 0$ : l'arc est du côté opposé (sens indirect), équivalent à $|\alpha|$ du côté direct.
\end{itemize}
\end{propriete}

\begin{propriete}[{Lieu : $\mathrm{mes}(\overrightarrow{MA}, \overrightarrow{MB}) = \alpha\ [\pi]$}]
L'ensemble $\mathcal{E}_{\pi} = \{ M \in \mathcal{P} \setminus \{A,B\} \mid \mathrm{mes}(\overrightarrow{MA}, \overrightarrow{MB}) = \alpha\ [\pi] \}$ est la \textbf{réunion de deux arcs de cercle} symétriques par rapport à $(AB)$, privés de $A$ et $B$.
C'est exactement la réunion des lieux pour $\alpha\ [2\pi]$ et $-\alpha\ [2\pi]$.
Cas particuliers :
\begin{itemize}
\item $\alpha = 0\ [\pi]$ : la droite $(AB)$ privée de $A, B$.
\item $\alpha = \pi/2\ [\pi]$ : le cercle de diamètre $[AB]$ privé de $A, B$ (deux arcs réunis).
\end{itemize}
\end{propriete}

\begin{aretenir}[Tableau récapitulatif des lieux]
\begin{center}
\begin{tabularx}{\linewidth}{|l|X|X|}\hline
\textbf{Condition} & \textbf{Nature du lieu} & \textbf{Construction} \\ \hline
$\mathrm{mes}(\overrightarrow{MA}, \overrightarrow{MB}) = 0\ [\pi]$ & Droite $(AB) \setminus \{A,B\}$ & Tracer $(AB)$ \\ \hline
$\mathrm{mes}(\overrightarrow{MA}, \overrightarrow{MB}) = \pi/2\ [\pi]$ & Cercle diamètre $[AB] \setminus \{A,B\}$ & Cercle centre milieu $[AB]$, rayon $AB/2$ \\ \hline
$\mathrm{mes}(\overrightarrow{MA}, \overrightarrow{MB}) = \alpha\ [2\pi]\ (\alpha \neq 0)$ & \textbf{Un seul arc} de cercle (orienté) & Tangente en $A$ faisant $\alpha$ avec $(AB)$, $\perp$ en $A$, médiatrice $[AB]$ \\ \hline
$\mathrm{mes}(\overrightarrow{MA}, \overrightarrow{MB}) = \alpha\ [\pi]\ (\alpha \neq 0, \pi/2)$ & \textbf{Deux arcs} symétriques (arc capable géométrique) & Même construction des deux côtés de $(AB)$ \\ \hline
\end{tabularx}
\end{center}
\end{aretenir}

\begin{attention}[Points exclus]
Dans \textbf{tous} les cas, les points $A$ et $B$ sont \textbf{exclus} du lieu car les vecteurs $\overrightarrow{MA}$ ou $\overrightarrow{MB}$ seraient nuls et l'angle non défini. On note souvent $\mathcal{E}^* = \mathcal{E} \setminus \{A,B\}$.
\end{attention}

\courssub{Applications : lieux géométriques et cocyclicité}

\begin{methode}[Montrer que 4 points sont cocycliques (angles orientés)]
Pour montrer que $A, B, C, D$ sont cocycliques (dans cet ordre ou non), il suffit de prouver l'une des égalités modulo $\pi$ :
\[ \mathrm{mes}(\overrightarrow{CA}, \overrightarrow{CB}) = \mathrm{mes}(\overrightarrow{DA}, \overrightarrow{DB})\ [\pi] \]
ou $\mathrm{mes}(\overrightarrow{AC}, \overrightarrow{BC}) = \mathrm{mes}(\overrightarrow{AD}, \overrightarrow{BD})\ [\pi]$, etc.
C'est la réciproque du théorème de l'angle inscrit orienté.
\end{methode}

\begin{methode}[Résoudre un problème de lieu géométrique]
\begin{enumerate}
\item Traduire la condition géométrique en égalité d'angles orientés (modulo $\pi$ ou $2\pi$).
\item Identifier la nature du lieu (droite, cercle, arc capable orienté, réunion de deux arcs).
\item Construire le lieu (centre, rayon, arc concerné).
\item Vérifier les points exclus ($A, B$ et éventuellement intersections avec droites interdites).
\end{enumerate}
\end{methode}

\begin{savaistu}[Le cercle d'Apollonius]
L'ensemble des points $M$ tels que $\dfrac{MA}{MB} = k\ (k>0, k\neq 1)$ est un cercle de diamètre $[IJ]$ où $I, J$ sont les pieds des bissectrices intérieure et extérieure de l'angle $\widehat{AMB}$ (ou points de division harmonique de $[AB]$). Ce cercle est orthogonal au cercle de diamètre $[AB]$. C'est un outil puissant pour coupler condition d'angle et condition de distance (voir Exercice 7).
\end{savaistu}

\courssub{Je vérifie mes connaissances}

\begin{corrige}[Exercice 1 — QCM : mesures principales]
La mesure principale d'un angle est l'unique mesure dans l'intervalle $]-\pi\,;\,\pi]$. On ajoute ou on retranche des multiples de $2\pi$.
\begin{enumerate}
\item $\dfrac{17\pi}{3} = \dfrac{18\pi - \pi}{3} = 6\pi - \dfrac{\pi}{3}$. Comme $6\pi = 3 \times 2\pi$, la mesure principale est $-\dfrac{\pi}{3}$. \textbf{Réponse b)}.
\item $-\dfrac{11\pi}{4} = -\dfrac{11\pi}{4} + 2\pi = -\dfrac{11\pi}{4} + \dfrac{8\pi}{4} = -\dfrac{3\pi}{4}$. Comme $-\dfrac{3\pi}{4} \in \left]-\pi\,;\,\pi\right]$, c'est la mesure principale. \textbf{Réponse b)}.
\item $\dfrac{25\pi}{6} = \dfrac{24\pi + \pi}{6} = 4\pi + \dfrac{\pi}{6}$. Comme $4\pi = 2 \times 2\pi$, la mesure principale est $\dfrac{\pi}{6}$. \textbf{Réponse a)}.
\item Sur le cercle trigonométrique :
\begin{itemize}
\item $-\dfrac{\pi}{3}$ : point image dans le quadrant IV, symétrique du point de $\dfrac{\pi}{3}$ par rapport à l'axe des abscisses ;
\item $-\dfrac{3\pi}{4}$ : point image dans le quadrant III, sur la bissectrice des axes négatifs ;
\item $\dfrac{\pi}{6}$ : point image dans le quadrant I, premier point de la graduation régulière en douzièmes de tour.
\end{itemize}
\end{enumerate}
\begin{attention}
La mesure principale appartient à $]-\pi\,;\,\pi]$ : $\pi$ est admis, mais pas $-\pi$. Une réponse comme $\dfrac{5\pi}{3}$ n'est jamais une mesure principale car $\dfrac{5\pi}{3} > \pi$.
\end{attention}
\end{corrige}

\begin{corrige}[Exercice 2 — Vrai ou faux, en justifiant]
\begin{enumerate}
\item \textbf{Vrai.} C'est le théorème de l'angle au centre : l'angle au centre interceptant le même arc qu'un angle inscrit mesure le double de celui-ci. Donc $\mathrm{mes}\,\widehat{AOB} = 2 \times 40^{\circ} = 80^{\circ}$.
\item \textbf{Vrai.} Deux angles inscrits interceptant le même arc ont la même mesure : si $M$ et $N$ sont sur le même arc d'extrémités $A$ et $B$, alors $\widehat{AMB}$ et $\widehat{ANB}$ interceptent l'arc opposé, d'où $\mathrm{mes}\,\widehat{AMB} = \mathrm{mes}\,\widehat{ANB}$.
\item \textbf{Vrai.} Le quadrilatère $AMBN$ est alors inscriptible avec $M$ et $N$ de part et d'autre de $(AB)$ : les angles $\widehat{AMB}$ et $\widehat{ANB}$ interceptent deux arcs dont la réunion est le cercle entier, donc la somme des angles au centre vaut $360^{\circ}$ et celle des angles inscrits vaut $180^{\circ}$ : $\mathrm{mes}\,\widehat{AMB} + \mathrm{mes}\,\widehat{ANB} = 180^{\circ}$.
\item \textbf{Vrai.} D'après le théorème de l'angle inscrit dans un demi-cercle (et sa réciproque) : $\mathrm{mes}\,\widehat{AMB} = 90^{\circ}$ si et seulement si $M$ appartient au cercle de diamètre $[AB]$, avec $M \neq A$ et $M \neq B$ (sinon l'angle n'est pas défini).
\end{enumerate}
\end{corrige}

\begin{corrige}[Exercice 3 — Angle inscrit et demi-tangente]
\begin{enumerate}
\item L'angle formé par la demi-tangente $(AT)$ et la corde $[AB]$ a pour mesure la moitié de l'angle au centre interceptant l'arc $\wideparen{AB}$ situé du côté opposé à $(AT)$ :
\[ \mathrm{mes}\,\widehat{TAB} = \frac{1}{2}\,\mathrm{mes}\,\widehat{AOB} \quad\Longrightarrow\quad \mathrm{mes}\,\widehat{AOB} = 2 \times 55^{\circ} = 110^{\circ}. \]
\item Tout angle inscrit $\widehat{AMB}$ où $M$ est sur l'arc capable associé à $(AT)$ (l'arc situé du même côté que $(AT)$ par rapport à $(AB)$) intercepte l'arc $\wideparen{AB}$ de mesure $110^{\circ}$, donc :
\[ \mathrm{mes}\,\widehat{AMB} = \frac{110^{\circ}}{2} = 55^{\circ}. \]
On retrouve bien l'égalité $\mathrm{mes}\,\widehat{AMB} = \mathrm{mes}\,\widehat{TAB}$, propriété fondamentale de la construction de l'arc capable.
\item Si $N$ appartient à l'autre arc, les angles $\widehat{AMB}$ et $\widehat{ANB}$ sont supplémentaires (question 3 de l'exercice 2) :
\[ \mathrm{mes}\,\widehat{ANB} = 180^{\circ} - 55^{\circ} = 125^{\circ}. \]
\end{enumerate}
\end{corrige}

\begin{corrige}[Exercice 4 — Lecture d'angles orientés]
\begin{enumerate}
\item $\mathrm{mes}(\overrightarrow{MA}, \overrightarrow{MB}) = 0 \; [\pi]$ signifie que les vecteurs $\overrightarrow{MA}$ et $\overrightarrow{MB}$ sont colinéaires, c'est-à-dire $M \in (AB)$. L'ensemble est \textbf{la droite $(AB)$ privée des points $A$ et $B$}.
\item $\mathrm{mes}(\overrightarrow{MA}, \overrightarrow{MB}) = \dfrac{\pi}{2} \; [\pi]$ signifie $(MA) \perp (MB)$ : l'ensemble est \textbf{le cercle de diamètre $[AB]$ privé de $A$ et $B$} (les deux arcs capables réunis).
\item $\mathrm{mes}(\overrightarrow{MA}, \overrightarrow{MB}) = \dfrac{\pi}{3} \; [2\pi]$ : l'ensemble est \textbf{un seul arc de cercle} d'extrémités $A$ et $B$ (privé de $A$ et $B$), celui situé du côté où le passage de $\overrightarrow{MA}$ à $\overrightarrow{MB}$ se fait dans le sens direct.
\item $\mathrm{mes}(\overrightarrow{MA}, \overrightarrow{MB}) = \dfrac{\pi}{3} \; [\pi]$ regroupe les conditions $\dfrac{\pi}{3}\;[2\pi]$ et $-\dfrac{\pi}{3}\;[2\pi]$ : l'ensemble est \textbf{la réunion de deux arcs de cercle} symétriques par rapport à $(AB)$, privés de $A$ et $B$.
\end{enumerate}
\end{corrige}

\courssub{Je m'entraîne}

\begin{corrige}[Exercice 5 — Construction d'un arc capable]
\begin{enumerate}
\item Le triangle $OAB$ est isocèle en $O$ avec $\mathrm{mes}\,\widehat{AOB} = 2 \times 75^{\circ} = 150^{\circ}$. En notant $I$ le milieu de $[AB]$, le triangle $OIA$ est rectangle en $I$ et $\mathrm{mes}\,\widehat{AOI} = 75^{\circ}$, donc :
\[ \sin 75^{\circ} = \frac{AI}{OA} = \frac{AB/2}{R} \quad\Longrightarrow\quad R = \frac{4}{\sin 75^{\circ}}. \]
Avec $\sin 75^{\circ} = \dfrac{\sqrt{6}+\sqrt{2}}{4} \approx 0{,}966$ : $R \approx \dfrac{4}{0{,}966} \approx \SI{4,14}{cm}$.
\item Construction :
\begin{enumerate}
\item tracer le segment $[AB]$ de longueur $\SI{8}{cm}$ ;
\item tracer la demi-droite $[At)$ telle que $\mathrm{mes}\,\widehat{tAB} = 75^{\circ}$ (au rapporteur, ou par construction du triangle équilatéral et bissections) ;
\item élever en $A$ la perpendiculaire à $[At)$ ;
\item tracer la médiatrice de $[AB]$ ;
\item l'intersection $O$ de ces deux droites est le centre du cercle ; tracer le cercle de centre $O$ passant par $A$ (et $B$) ;
\item l'arc capable est l'arc d'extrémités $A$ et $B$ situé du même côté que $[At)$, privé de $A$ et $B$.
\end{enumerate}
\item Vérification : on choisit un point $M$ quelconque de l'arc construit et on mesure $\widehat{AMB}$ au rapporteur : on trouve $75^{\circ}$ (aux erreurs de tracé près). C'est cohérent avec le théorème : l'angle entre la tangente et la corde vaut l'angle inscrit, et la perpendiculaire en $A$ à $[At)$ passe par le centre car la tangente est perpendiculaire au rayon.
\end{enumerate}
\end{corrige}

\begin{corrige}[Exercice 6 — Arcs capables orientés]
\begin{enumerate}
\item $(\mathcal{E}_1)$ : $\mathrm{mes}(\overrightarrow{MA}, \overrightarrow{MB}) = \dfrac{\pi}{4} \; [2\pi]$. C'est l'arc capable orienté : un seul arc de cercle d'extrémités $A$ et $B$ (privé de $A$ et $B$), situé du côté de $(AB)$ où la rotation de $\overrightarrow{MA}$ vers $\overrightarrow{MB}$ s'effectue dans le sens direct (sens trigonométrique). On le construit en traçant $[At)$ avec $\mathrm{mes}(\overrightarrow{AB}, \overrightarrow{At}) = \dfrac{\pi}{4}$, la perpendiculaire en $A$ à $[At)$ et la médiatrice de $[AB]$. \textbf{Tracé en rouge.}
\item $(\mathcal{E}_2)$ : condition modulo $\pi$, c'est-à-dire $\dfrac{\pi}{4}\;[2\pi]$ ou $-\dfrac{\pi}{4}\;[2\pi]$. C'est la réunion de deux arcs de cercle symétriques par rapport à la droite $(AB)$, privés de $A$ et $B$. \textbf{Tracé en bleu.}
\item L'arc rouge $(\mathcal{E}_1)$ est l'un des deux arcs bleus : $\boxed{(\mathcal{E}_1) \subset (\mathcal{E}_2)}$. Le complémentaire $(\mathcal{E}_2) \setminus (\mathcal{E}_1)$ est l'arc symétrique de $(\mathcal{E}_1)$ par rapport à $(AB)$, caractérisé par :
\[ \mathrm{mes}(\overrightarrow{MA}, \overrightarrow{MB}) = -\frac{\pi}{4} \; [2\pi]. \]
\end{enumerate}
\end{corrige}

\begin{corrige}[Exercice 7 — Angle droit et rapport de distances]
\begin{enumerate}
\item $\mathrm{mes}(\overrightarrow{MA}, \overrightarrow{MB}) = \dfrac{\pi}{2} \; [\pi]$ équivaut à $(MA) \perp (MB)$. D'après le théorème de l'angle droit inscrit, $(\mathcal{C})$ est le cercle de diamètre $[AB]$, privé de $A$ et $B$ : centre le milieu de $[AB]$, rayon $\dfrac{AB}{2} = \SI{3}{cm}$.
\item Cherchons les points $M$ tels que $MA = 2\,MB$, c'est-à-dire $\dfrac{MA}{MB} = 2$. Les points $I$ et $J$ de $(AB)$ définis par $\overrightarrow{IA} = -2\overrightarrow{IB}$ (barycentre, $I$ entre $A$ et $B$) et $\overrightarrow{JA} = 2\overrightarrow{JB}$ ($J$ à l'extérieur, côté de $B$) vérifient $\dfrac{IA}{IB} = \dfrac{JA}{JB} = 2$. Ce sont les pieds des bissectrices intérieure et extérieure issues de $M$ dans le triangle $MAB$ : pour tout point $M$ tel que $MA = 2\,MB$, $(MI)$ et $(MJ)$ sont les bissectrices, donc $(MI) \perp (MJ)$, c'est-à-dire $\mathrm{mes}(\overrightarrow{MI}, \overrightarrow{MJ}) = \dfrac{\pi}{2}\;[\pi]$. Par le théorème de l'angle droit, $M$ appartient au cercle $(\Gamma)$ de diamètre $[IJ]$. Réciproquement, tout point de ce cercle (privé de $(AB)$) vérifie $MA = 2\,MB$ (cercle d'Apollonius).

Positionnons $I$ et $J$ sur une graduation où $A$ a pour abscisse $0$ et $B$ l'abscisse $6$ : $I$ vérifie $IA = 2\,IB$ avec $IA + IB = 6$, d'où $IB = 2$, $IA = 4$ ; $J$ vérifie $JA = 2\,JB$ avec $JA - JB = 6$, d'où $JB = 6$, $JA = 12$. Donc $IJ = 8$ et $(\Gamma)$ a pour rayon $4$.
\item Construction : placer $I$ à $\SI{4}{cm}$ de $A$ sur $[AB]$, $J$ à $\SI{6}{cm}$ de $B$ à l'extérieur ; tracer le cercle de diamètre $[IJ]$. Les intersections de $(\mathcal{C})$ et $(\Gamma)$ donnent les points solutions. Le centre de $(\mathcal{C})$ est le milieu de $[AB]$ (abscisse $3$), le centre de $(\Gamma)$ est le milieu de $[IJ]$ (abscisse $8$) : les centres sont distants de $\mathbf{5}$ cm. Les rayons sont $3$ et $4$. Comme $|4-3| = 1 < 5 < 7 = 4+3$, les deux cercles se coupent en \textbf{deux points} symétriques par rapport à $(AB)$. Il y a donc \textbf{deux points $M$ solutions}.
\end{enumerate}
\end{corrige}

\begin{corrige}[Exercice 8 — Cocyclicité au marché]
\begin{enumerate}
\item Les points $C$ et $D$ sont du même côté de $(AB)$ et vérifient $\mathrm{mes}\,\widehat{ACB} = \mathrm{mes}\,\widehat{ADB} = 42^{\circ}$. D'après la réciproque du théorème des angles inscrits (propriété de l'arc capable), $C$ et $D$ appartiennent au même arc capable d'extrémités $A$ et $B$, donc au même cercle passant par $A$ et $B$. Les quatre points $A$, $B$, $C$, $D$ sont \textbf{cocycliques}.
\item $E$ est de l'autre côté de $(AB)$ et $\mathrm{mes}\,\widehat{AEB} = 138^{\circ}$. Or $42^{\circ} + 138^{\circ} = 180^{\circ}$ : les angles sont supplémentaires. Pour un quadrilatère inscriptible, deux points situés de part et d'autre de $(AB)$ voient $[AB]$ sous des angles supplémentaires. Donc \textbf{oui}, $E$ appartient au même cercle (il est sur l'arc capable complémentaire).
\item Construction : le centre $O$ est l'intersection des médiatrices de $[AB]$ et $[AC]$ (ou de $[BC]$). Tracer le cercle de centre $O$ passant par $A$, $B$, $C$. Un point $D'$ possible est n'importe quel point de l'arc contenant $C$ (distinct de $A$, $B$, $C$) : on vérifiera $\mathrm{mes}\,\widehat{AD'B} = 42^{\circ}$.
\end{enumerate}
\end{corrige}

\courssub{Je me prépare au Probatoire}

\begin{corrige}[Exercice 9 — Orthocentre et cercle circonscrit]
\textbf{Partie A}
\begin{enumerate}
\item Figure : triangle $ABC$ acutangle, hauteurs $(AA')$, $(BB')$, $(CC')$ concourantes en $H$ intérieur au triangle.

$(HB)$ est portée par la hauteur issue de $B$, donc $(HB) \perp (AC)$ ; de même $(HC) \perp (AB)$. Pour deux couples de droites perpendiculaires, les angles orientés sont égaux modulo $\pi$ :
\[ \mathrm{mes}(\overrightarrow{HB}, \overrightarrow{HC}) = \mathrm{mes}(\overrightarrow{AC}, \overrightarrow{AB}) \; [\pi] = -\mathrm{mes}(\overrightarrow{AB}, \overrightarrow{AC}) \; [\pi]. \]
Comme le triangle est acutangle, $H$ est intérieur et $B$, $C$ sont vus depuis $H$ dans la configuration opposée à celle de $A$ : en mesures modulo $2\pi$, on obtient
\[ \mathrm{mes}(\overrightarrow{HB}, \overrightarrow{HC}) = \mathrm{mes}(\overrightarrow{AB}, \overrightarrow{AC}) + \pi \; [2\pi]. \]
\emph{Justification du $+\pi$ :} en passant aux angles géométriques, $\mathrm{mes}\,\widehat{BHC} = 180^{\circ} - \mathrm{mes}\,\widehat{BAC}$ (voir question 2), ce qui impose le décalage de $\pi$.
\item En mesures géométriques, $\mathrm{mes}\,\widehat{BHC} = \left|\mathrm{mes}(\overrightarrow{HB}, \overrightarrow{HC})\right|$ ramenée dans $[0\,;\,\pi]$. Comme $\mathrm{mes}\,\widehat{BAC} \in \left]0\,;\,\dfrac{\pi}{2}\right[$ (triangle acutangle), $\mathrm{mes}(\overrightarrow{AB}, \overrightarrow{AC}) + \pi$ a pour valeur absolue ramenée dans $[0\,;\,\pi]$ : $\pi - \mathrm{mes}\,\widehat{BAC}$. Donc :
\[ \mathrm{mes}\,\widehat{BHC} = \pi - \mathrm{mes}\,\widehat{BAC} = 180^{\circ} - \mathrm{mes}\,\widehat{BAC}. \]
\end{enumerate}

\textbf{Partie B}
\begin{enumerate}
\item La symétrie orthogonale d'axe $(BC)$ fixe $B$ et $C$ et envoie $H$ sur $H_1$. Une symétrie axiale transforme un angle orienté en son opposé, donc :
\[ \mathrm{mes}(\overrightarrow{H_1B}, \overrightarrow{H_1C}) = -\mathrm{mes}(\overrightarrow{HB}, \overrightarrow{HC}) \; [2\pi]. \]
\item D'après A.1, on a $\mathrm{mes}(\overrightarrow{HB}, \overrightarrow{HC}) = -\mathrm{mes}(\overrightarrow{AB}, \overrightarrow{AC})\ [\pi]$. Par conséquent :
\[ \mathrm{mes}(\overrightarrow{H_1B}, \overrightarrow{H_1C}) = -(-\mathrm{mes}(\overrightarrow{AB}, \overrightarrow{AC})) \; [\pi] = \mathrm{mes}(\overrightarrow{AB}, \overrightarrow{AC}) \; [\pi]. \]
\item Cette égalité modulo $\pi$ exprime exactement la cocyclicité des quatre points $A$, $B$, $C$, $H_1$ (réciproque du théorème de l'angle inscrit orienté). Comme $A$, $B$, $C$ définissent $(\mathcal{C})$, on conclut : $\boxed{H_1 \in (\mathcal{C})}$.
\item Le même raisonnement appliqué aux symétries par rapport à $(CA)$ et $(AB)$ montre que $H_2$ et $H_3$ appartiennent aussi à $(\mathcal{C})$ : \textbf{les trois symétriques de l'orthocentre par rapport aux côtés du triangle sont sur le cercle circonscrit}.
\end{enumerate}

\textbf{Partie C}

Par le théorème de l'angle au centre, $\mathrm{mes}\,\widehat{BOC} = 2 \times 60^{\circ} = 120^{\circ}$. Dans le triangle isocèle $OBC$, en notant $I$ le milieu de $[BC]$ : $\sin 60^{\circ} = \dfrac{BC/2}{R}$, donc
\[ R = \frac{BC}{2\sin 60^{\circ}} = \frac{9}{2 \times \frac{\sqrt{3}}{2}} = \frac{9}{\sqrt{3}} = 3\sqrt{3} \approx \SI{5,2}{cm}. \]

\textbf{Barème indicatif (10 points) :} figure 1 pt ; A.1 : 2 pts ; A.2 : 1 pt ; B.1 : 1 pt ; B.2 : 1,5 pt ; B.3 : 1 pt ; B.4 : 0,5 pt ; C : 2 pts.

\textbf{Points de vigilance :} citer explicitement l'hypothèse « triangle acutangle » (elle place $H$ à l'intérieur et fixe le signe) ; en B.3, énoncer la réciproque du théorème de l'angle inscrit ; en C, ne pas oublier que l'angle au centre est le double de l'angle inscrit interceptant le même arc.
\end{corrige}

\begin{corrige}[Exercice 10 — Deux cercles sécants et angle constant]
\textbf{Partie A}
\begin{enumerate}
\item Sur les deux figures, la mesure de $\widehat{MBN}$ est la même (aux erreurs de mesure près). \textbf{Conjecture :} l'angle $\widehat{MBN}$ est constant lorsque la droite $(\Delta)$ varie.
\end{enumerate}

\textbf{Partie B}
\begin{enumerate}
\item Dans $(\mathcal{C}_1)$, les points $A$, $B$, $M$ sont sur le cercle. Par le théorème de l'angle inscrit orienté, pour tout point $M$ de $(\mathcal{C}_1)$ distinct de $A$ et $B$, l'angle orienté $\mathrm{mes}(\overrightarrow{BM}, \overrightarrow{BA})$ est constant modulo $\pi$ (il intercepte l'arc $AB$ fixe). Or $M \in (\Delta)$ et $A \in (\Delta)$, donc la direction de $(AM)$ est celle de $(\Delta)$. En utilisant la tangente $(AT_1)$ en $A$ à $(\mathcal{C}_1)$, le théorème de l'angle entre tangente et corde donne :
\[ \mathrm{mes}(\overrightarrow{BM}, \overrightarrow{BA}) = \mathrm{mes}(\overrightarrow{AT_1}, \overrightarrow{AB}) \; [\pi], \]
quantité qui ne dépend que de $(\mathcal{C}_1)$ et de la corde $[AB]$, donc \textbf{constante modulo $\pi$} quand $(\Delta)$ varie.
\item De même dans $(\mathcal{C}_2)$, avec $(AT_2)$ tangente en $A$ à $(\mathcal{C}_2)$ :
\[ \mathrm{mes}(\overrightarrow{BA}, \overrightarrow{BN}) = \mathrm{mes}(\overrightarrow{BA}, \overrightarrow{AT_2}) \; [\pi], \quad \text{constante modulo } \pi. \]
\item Par la relation de Chasles :
\[ \mathrm{mes}(\overrightarrow{BM}, \overrightarrow{BN}) = \mathrm{mes}(\overrightarrow{BM}, \overrightarrow{BA}) + \mathrm{mes}(\overrightarrow{BA}, \overrightarrow{BN}) \; [\pi]. \]
Somme de deux constantes modulo $\pi$, c'est une constante modulo $\pi$. En utilisant les expressions avec les tangentes :
\[ \mathrm{mes}(\overrightarrow{BM}, \overrightarrow{BN}) = \mathrm{mes}(\overrightarrow{AT_1}, \overrightarrow{AB}) + \mathrm{mes}(\overrightarrow{AB}, \overrightarrow{AT_2}) = \mathrm{mes}(\overrightarrow{AT_1}, \overrightarrow{AT_2}) \; [\pi]. \]
La constante est donc \textbf{l'angle des deux tangentes en $A$} aux deux cercles.
\end{enumerate}

\textbf{Partie C}
\begin{enumerate}
\item Quand $N$ décrit $(\mathcal{C}_2) \setminus \{A, B\}$, la droite $(AN) = (\Delta)$ prend toutes les directions de sécantes par $A$ (sauf $(AB)$ et la tangente commune éventuelle), donc $M$, second point d'intersection de $(\Delta)$ avec $(\mathcal{C}_1)$, décrit $(\mathcal{C}_1)$ privé de $A$ et $B$.
\item Si les deux cercles sont concentriques de centre $O$ (rayons $r_1$ et $r_2$), $M$ et $N$ sont sur la même droite passant par $A$, avec $OM = r_1$ et $ON = r_2$. Le milieu $I$ de $[MN]$ : la projection orthogonale de $O$ sur $(\Delta)$ est le milieu commun des cordes, et l'on montre que $OI$ est constant ou que $I$ se déduit par homothétie. On admet que $I$ décrit \textbf{le cercle de centre $O$ et de rayon $\dfrac{r_1 + r_2}{2}$} (cercle médian, privé des points correspondant aux positions exclues).
\end{enumerate}

\textbf{Barème indicatif (10 points) :} figures et conjecture 1,5 pt ; B.1 : 2 pts ; B.2 : 1 pt ; B.3 : 2 pts ; C.1 : 1,5 pt ; C.2 : 2 pts.

\textbf{Points de vigilance :} travailler modulo $\pi$ (les points $M$, $N$ peuvent changer d'arc) ; citer le théorème « angle inscrit = angle entre corde et tangente » ; dans Chasles, vérifier l'ordre des vecteurs.
\end{corrige}

\begin{corrige}[Exercice 11 — Navigation entre deux phares]
\textbf{Partie A}
\begin{enumerate}
\item Échelle $1/100\,000$ : $\SI{1}{cm}$ sur la carte représente $\SI{100\,000}{cm} = \SI{1}{km}$. Donc $F_1F_2 = 12 \times 1 = \SI{12}{km}$.
\item Sur la carte : $R = \dfrac{F_1F_2}{2\sin 30^{\circ}} = \dfrac{12}{2 \times \frac{1}{2}} = \SI{12}{cm}$. En réalité : $R = \SI{12}{km}$ (l'échelle conserve les angles et multiplie les longueurs).
\item Pour un point $B$ situé \textbf{sur} l'arc capable, $\mathrm{mes}\,\widehat{F_1BF_2} = 30^{\circ}$. Pour un point \textbf{intérieur} au segment de cercle (entre la corde $[F_1F_2]$ et l'arc), l'angle est \textbf{plus grand} que $30^{\circ}$ (le point est plus proche du segment, on le voit sous un angle plus ouvert) ; pour un point \textbf{extérieur}, l'angle est plus petit que $30^{\circ}$. La condition de sécurité « angle $\geq 30^{\circ}$ » correspond donc au \textbf{segment de cercle (intérieur de l'arc capable, côté mer), arc inclus}, privé éventuellement du voisinage immédiat de $[F_1F_2]$.
\end{enumerate}

\textbf{Partie B}
\begin{enumerate}
\item Construction : tracer $[F_1F_2]$ de $\SI{12}{cm}$ ; tracer $[F_1t)$ avec $\mathrm{mes}\,\widehat{tF_1F_2} = 30^{\circ}$ du côté mer ; élever la perpendiculaire en $F_1$ à $[F_1t)$ ; tracer la médiatrice de $[F_1F_2]$ ; leur intersection $O$ est le centre ($OF_1 = \SI{12}{cm}$, on vérifie que le triangle $OF_1F_2$ est équilatéral, ce qui est cohérent : angle au centre $60^{\circ}$) ; tracer l'arc d'extrémités $F_1$, $F_2$ du côté mer.
\item Hachurer la région comprise entre la corde $[F_1F_2]$ et l'arc (côté mer).
\item Dans le triangle $F_1BF_2$ : $\mathrm{mes}\,\widehat{F_1BF_2} = 30^{\circ}$, $F_1F_2 = 12$, $BF_1 = 10$. Par la loi des sinus :
\[ \frac{F_1F_2}{\sin \widehat{B}} = \frac{BF_1}{\sin \widehat{F_2}} \quad\Longrightarrow\quad \sin \widehat{F_2} = \frac{10 \times \sin 30^{\circ}}{12} = \frac{10 \times 0{,}5}{12} = \frac{5}{12} \approx 0{,}4167. \]
D'où $\widehat{F_2} \approx 24{,}6^{\circ}$, puis $\widehat{F_1} \approx 180^{\circ} - 30^{\circ} - 24{,}6^{\circ} = 125{,}4^{\circ}$. À nouveau la loi des sinus :
\[ BF_2 = \frac{F_1F_2 \times \sin \widehat{F_1}}{\sin \widehat{B}} = \frac{12 \times \sin 125{,}4^{\circ}}{0{,}5} \approx \frac{12 \times 0{,}816}{0{,}5} \approx \SI{19,6}{cm}. \]
Distance réelle : $19{,}6 \times 1 \approx \SI{19,6}{km}$. (Il existe une seconde configuration avec $\widehat{F_2} \approx 155{,}4^{\circ}$, impossible car la somme dépasserait $180^{\circ}$ avec l'angle de $30^{\circ}$... en fait $155{,}4 + 30 = 185{,}4 > 180$, donc bien une seule solution.)
\end{enumerate}

\textbf{Partie C}

La trajectoire est l'\textbf{arc capable d'angle $45^{\circ}$} construit sur $[F_1F_2]$, du côté mer : ensemble des points $B$ tels que $\mathrm{mes}\,\widehat{F_1BF_2} = 45^{\circ}$. Construction : rayon $R' = \dfrac{12}{2\sin 45^{\circ}} = \dfrac{12}{\sqrt{2}} = 6\sqrt{2} \approx \SI{8,5}{cm}$ sur la carte ; tracer $[F_1u)$ avec $\mathrm{mes}\,\widehat{uF_1F_2} = 45^{\circ}$, la perpendiculaire en $F_1$ à $[F_1u)$, la médiatrice de $[F_1F_2]$ ; le centre $O'$ est leur intersection ; l'arc d'extrémités $F_1$, $F_2$ situé du côté mer (privé de $F_1$ et $F_2$) est la trajectoire cherchée. Cet arc est plus proche de la côte que l'arc à $30^{\circ}$, ce qui correspond bien à l'ordre de se rapprocher.

\textbf{Barème indicatif (10 points) :} A.1 : 1 pt ; A.2 : 1 pt ; A.3 : 1,5 pt ; B.1 : 1,5 pt ; B.2 : 0,5 pt ; B.3 : 2,5 pts ; C : 2 pts.

\textbf{Points de vigilance :} conversion d'échelle ($\SI{1}{cm} \leftrightarrow \SI{1}{km}$, pas $\SI{1000}{km}$ !) ; dans la loi des sinus, vérifier la cohérence de la somme des angles pour écarter le cas ambigu ; pour le lieu, préciser « arc privé de $F_1$ et $F_2$ » car l'angle n'y est pas défini.
\end{corrige}

\newpage

\coursec{Fiche bilan}

\begin{popfigure}
\begin{center}
\begin{tikzpicture}[
  mindmap,
  grow cyclic,
  every node/.style={concept, text=white, font=\small},
  root concept/.append style={concept color=popdark, font=\bfseries},
  level 1/.append style={level distance=3.4cm, sibling angle=60},
  level 1 concept/.append style={font=\footnotesize},
  scale=0.92, transform shape
]
\node[root concept]{Angles orientés et arcs capables}
  child[concept color=popblue]{ node[align=center] {Angles orientés de vecteurs\\ mesure principale\\ relation de Chasles} }
  child[concept color=poporange]{ node[align=center] {Angle inscrit\\ $=$ moitié de\\ l'angle au centre} }
  child[concept color=popgreen]{ node[align=center] {Angle corde--tangente\\ $=$ angle inscrit\\ interceptant la corde} }
  child[concept color=poppurple]{ node[align=center] {Arc capable\\ $\operatorname{mes}(\widehat{AMB})=x$\\ arc de cercle} }
  child[concept color=poppink]{ node[align=center] {Arcs orientés\\ $(\overrightarrow{MA},\overrightarrow{MB})=x+2k\pi$\\ ou $x+k\pi$} }
  child[concept color=popgold]{ node[align=center] {Applications\\ lieux géométriques\\ cocyclicité} };
\end{tikzpicture}
\end{center}
\legende{Carte mentale de la leçon : les six idées forces autour des arcs capables.}
\end{popfigure}

\begin{bilan}
\textbf{Définitions clés}
\begin{itemize}
  \item \cle{Mesure principale} d'un angle orienté : unique mesure appartenant à $\left]-\pi\,;\,\pi\right]$.
  \item \cle{Angle inscrit} : angle $\widehat{AMB}$ dont le sommet $M$ est sur un cercle et dont les côtés coupent le cercle en $A$ et $B$.
  \item \cle{Arc capable} : ensemble des points $M$ d'un demi-plan de frontière $(AB)$ tels que $\operatorname{mes}(\widehat{AMB})=x$ : c'est un arc de cercle passant par $A$ et $B$ (extrémités exclues).
\end{itemize}

\textbf{Théorèmes et formules à connaître par c\oe ur}
\begin{center}
\begin{tabularx}{\linewidth}{|l|X|}
\hline
\rowcolor{popblueL} \textbf{Propriété} & \textbf{Énoncé} \\
\hline
Angle inscrit / angle au centre & $\operatorname{mes}(\widehat{AMB})=\dfrac{1}{2}\operatorname{mes}(\widehat{AOB})$ pour $A,B,M$ sur le cercle de centre $O$. \\
\hline
Corde et demi-tangente & L'angle entre une corde $[AB]$ et la demi-tangente en $A$ égale tout angle inscrit interceptant $[AB]$ de l'autre côté. \\
\hline
Arc capable $\operatorname{mes}(\widehat{AMB})=x$ & Arc de cercle d'extrémités $A$ et $B$, symétrique possible par rapport à $(AB)$ ; $x=\dfrac{\pi}{2}$ donne le cercle de diamètre $[AB]$. \\
\hline
$(\overrightarrow{MA},\overrightarrow{MB})=x+2k\pi$ & Un \emph{seul} arc orienté de $A$ vers $B$ (les points $A$ et $B$ exclus). \\
\hline
$(\overrightarrow{MA},\overrightarrow{MB})=x+k\pi$ & Le cercle entier passant par $A$ et $B$, privé de $A$ et $B$ (réunion des deux arcs). \\
\hline
Cocyclicité & $A,B,M,N$ cocycliques $\iff (\overrightarrow{MA},\overrightarrow{MB})=(\overrightarrow{NA},\overrightarrow{NB})+k\pi$. \\
\hline
\end{tabularx}
\end{center}

\textbf{Méthodes express}
\begin{itemize}
  \item \cle{Construire un arc capable} : tracer la demi-droite $[At)$ telle que $\operatorname{mes}(\widehat{BAt})=x$ ; élever la perpendiculaire en $A$ à $[At)$ ; elle coupe la médiatrice de $[AB]$ au centre $O$ ; tracer l'arc de cercle de centre $O$ situé du côté opposé à $[At)$.
  \item \cle{Prouver une cocyclicité} : montrer l'égalité de deux angles inscrits (ou orientés modulo $\pi$) interceptant la même corde.
  \item \cle{Trouver un lieu} : repérer un segment fixe $[AB]$ et un angle constant sous lequel il est vu, puis conclure « arc capable » ou « cercle de diamètre $[AB]$ » si l'angle est droit.
\end{itemize}
\end{bilan}

\begin{aretenir}[Checklist avant le Probatoire]
\begin{itemize}
  \item[$\square$] Je sais déterminer la mesure principale d'un angle orienté (dans $\left]-\pi\,;\,\pi\right]$).
  \item[$\square$] Je sais utiliser la relation de Chasles sur les angles orientés de vecteurs.
  \item[$\square$] Je connais le théorème de l'angle inscrit : $\operatorname{mes}(\widehat{AMB})=\dfrac{1}{2}\operatorname{mes}(\widehat{AOB})$.
  \item[$\square$] Je sais que l'angle entre une corde et la demi-tangente égale l'angle inscrit associé.
  \item[$\square$] Je sais construire l'arc capable $\operatorname{mes}(\widehat{AMB})=x$ avec la tangente et la médiatrice de $[AB]$.
  \item[$\square$] Je sais que $\operatorname{mes}(\widehat{AMB})=\dfrac{\pi}{2}$ caractérise le cercle de diamètre $[AB]$.
  \item[$\square$] Je distingue bien $x+2k\pi$ (un seul arc orienté) et $x+k\pi$ (le cercle privé de $A$ et $B$).
  \item[$\square$] Je n'oublie jamais d'\emph{exclure} les points $A$ et $B$ des lieux géométriques.
  \item[$\square$] Je sais prouver que quatre points sont cocycliques par égalité d'angles modulo $\pi$.
  \item[$\square$] Je pense à l'arc symétrique par rapport à $(AB)$ quand l'énoncé ne précise pas de demi-plan.
  \item[$\square$] Je rédige proprement : « l'ensemble cherché est l'arc de cercle \dots, privé de $A$ et $B$ ».
\end{itemize}
\end{aretenir}

\findecours

\end{document}
